% !TEX root = ../../tree_notes.tex
% =====================================================================
%  Chapter 8 (map v2) -- Balanced Search Trees       Status: WRITTEN
%  (folder keeps its ch06_balanced_trees name; printed as Chapter 8,
%   label c06.)
%
%  Files, in reading order:
%    00_opener.tex              glance + badges, reading paths, ladder
%                              staircase, the two running examples (the chain
%                              C and the zig-zag Z), roadmap, recognition guide,
%                              notation
%    01_why_rotations.tex       8.1 Why Balance? The Rotation       sec:c06-rotations
%    02_avl.tex                 8.2 AVL Trees                        sec:c06-avl
%    03_treaps.tex              8.3 Treaps, Split/Merge, Implicit    sec:c06-treaps
%    04_redblack.tex            8.4 Red-Black Trees (conceptual)     sec:c06-redblack
%    05_splay.tex               8.5 Splay Trees (intro, amortized)   sec:c06-splay
%    99_chapter_end.tex         the 15-part revision contract + Summit (LC 493)
%
%  This chapter OWNS "make the height small BY CONSTRUCTION." Chapter 7 ended
%  with the rotting descent -- every BST routine is O(h), and h can reach n.
%  Here comes the machinery that forces h = O(log n) forever: the rotation (the
%  one local, inorder-preserving move), AVL trees (rebalance on every edit, the
%  provable Fibonacci bound), treaps (random priorities instead of case work,
%  plus the split/merge kit and implicit keys), red-black trees (why std::set is
%  O(log n) -- the guarantee Chapter 7 deferred), and a first look at splay
%  trees (amortised self-adjusting; the full block lives in Chapter 25).
%
%  Anti-repetition -- deliberately NOT re-taught here:
%    * the BST invariant, search / insert / delete, inorder = sorted
%        -> Chapter 7 (\chref{c05}); we build ON the descent, never re-derive it
%    * std::set / std::map API, lower_bound/upper_bound, Range Module, calendars
%        -> 7.4 (\secref{c05-contests}); here we explain WHY they cost O(log n)
%           (red-black) and when you must roll your own, without re-listing the API
%    * the three DFS orders / Morris
%        -> Chapter 3 (\chref{c03}); the rotation only needs inorder, already proved
%    * Fenwick trees and segment trees (LC 307, CSES 1190 also live there)
%        -> Chapters 18--19 (\chref{c14}, \chref{c13}); we present those dynamic
%           range tasks through the implicit-treap lens and hand the engines forward
%    * lazy propagation as a general theory
%        -> Chapter 20 (\chref{c15}); the treap's lazy reverse is developed lightly
%    * HEAPS and priority_queue (the median-from-stream engines)
%        -> Chapter 10 (\chref{c07}); referenced, not re-taught
%    * CARTESIAN TREES and the monotonic stack (LC 84, subarray minimums,
%      CF 817D Showmen, next-smaller links)
%        -> Chapter 9 (\chref{c26}). THIS OVERRIDES THE OLD ch06.tex STUB, which
%           still listed Cartesian trees as its topic 4; per topics_by_chapter.txt
%           the Cartesian/monotonic-stack material moved OUT of this folder into
%           the new Chapter 9. It is intentionally ABSENT here.
%    * full splay / link-cut / Euler-tour dynamic trees
%        -> Chapter 25 (\chref{c23}); 8.5 is the intro that \chref{c23} already
%           forward-references ("splay rotations, from Chapter 8")
%    * AVL/red-black/treap amortised-and-deamortised proofs, join-based purely
%      functional sets
%        -> beyond this book's promise; noted, not developed
% =====================================================================
\bookchapter{Balanced Search Trees}{c06}
\chapterinfo{LC Medium--Hard, CF 1700--2100, GFG Hard}{Written}

% !TEX root = ../../tree_notes.tex
% =====================================================================
%  Chapter 8 opener -- "Chapter at a Glance", reading paths + ladder
%  staircase, the two running examples (the degenerate chain C and its
%  balanced twin; the zig-zag Z that forces a double rotation), the
%  "who balances, how" comparison, roadmap, recognition guide, notation.
%  NO numbered boxes here.
% =====================================================================

\begin{chaptertheorybox}{Chapter at a Glance}
Chapter~7 (\chref{c05}) spent one invariant -- $\text{left}<\text{node}<\text{right}$ -- to turn a
binary tree into a search index, and closed on a debt: every routine costs $\bigO{h}$, and $h$ is
\emph{not yours to control}. Feed the keys in sorted order and the BST rots into a linked list, so
``search a million times'' silently becomes ``scan a million times.'' This chapter pays that debt by
\emph{construction}: it keeps the height logarithmic, forever, no matter what order the keys arrive
in. There are exactly four ways to do it, and this chapter teaches all four dialects of one sentence
-- \emph{restore the shape without ever breaking the inorder line}:
\begin{itemize}[nosep]
  \item \textbf{The rotation is the only legal move.} A single local pointer surgery re-weights a
        subtree while leaving the sorted inorder sequence \emph{identical}. Every balanced tree, from
        AVL to red-black to splay, is a policy for \emph{when} to rotate and \emph{how much} to care.
        \hfill \lv{2}--\lv{3}
  \item \textbf{Three philosophies of staying short.} \emph{Deterministic rebalancing} (AVL) enforces
        a strict height rule and rotates on every edit; \emph{randomisation} (treap) trades case work
        for a coin flip and wins $\bigO{\log n}$ in expectation; \emph{colour invariants} (red-black)
        relax the rule just enough that $\le 2$ recolours/rotisons suffice -- which is \emph{why} your
        \code{std::set} is fast. \hfill \lv{3}--\lv{5}
  \item \textbf{A height you control unlocks augmentation.} Once $h=\bigO{\log n}$ by construction, a
        size-augmented or implicit-keyed tree answers rank, select, and range questions on a
        \emph{dynamic sequence} -- the promised, guaranteed version of the fragile trick the Chapter~7
        Summit (\probref{c05-summit}) could only limp through. \hfill \lv{4}--\lv{5}
\end{itemize}

\textbf{You will learn to}
\begin{itemize}[nosep]
  \item perform and reason about the four rotation shapes (LL, RR, LR, RL) as the unique
        inorder-preserving fixes for a height violation; \hfill\lv{2}--\lv{3}
  \item insert and delete in an AVL tree, tracking balance factors up the access path and stopping at
        the first restored ancestor; \hfill\lv{3}--\lv{4}
  \item build a treap by assigning random priorities, and wield \code{split}/\code{merge} as the
        universal toolkit for insert, erase, rank, and range operations; \hfill\lv{4}
  \item explain, from the red-black invariants, why \code{std::set}/\code{std::map} guarantee
        $\bigO{\log n}$ -- the promise Chapter~7 borrowed and now repays; \hfill\lv{3}--\lv{4}
  \item recognise an \emph{implicit-key} treap as a dynamic array (insert/delete by index, reverse a
        range) and a splay tree as an amortised, locality-friendly container. \hfill\lv{4}--\lv{5}
\end{itemize}

\textbf{Prerequisites.} Chapter~7 (\chref{c05}) in full -- the BST descent, insertion, three-case
deletion, and the $\bigO{h}$-is-the-whole-cost warning are assumed, never re-derived. The inorder
walk is a black box from Chapter~3 (\chref{c03}); here we only \emph{lean on} the fact that a
rotation leaves it unchanged.
\end{chaptertheorybox}

% ---------------- Station 1b: reading paths + ladder position ----------------
\begin{chaptertheorybox}{How to Read This Chapter (two paths up the ladder)}
Level badges $\lv1$--$\lv6$: $\lv2$ is ``do the rotation by hand,'' $\lv3$ is a solid interview
medium, $\lv4$ is implementing a whole balanced structure cold, and $\lv5$ is an LC Hard where the
naive BST's height is the enemy the structure exists to slay.

\textbf{Beginner path ($\lv2$--$\lv3$).} Read \secref{c06-rotations} end to end -- it fixes \emph{why}
height is cost and gives you the four rotations as the one legal repair. Then \secref{c06-avl} through
insertion and the four rebalance cases on the running examples $C$ and $Z$. Take every
$\lv2$--$\lv{3}$ solved problem. You may \emph{skip on a first pass}: treap \code{split}/\code{merge}
and implicit keys, the red-black invariant proof sketch, splay amortisation, and the Summit. By the
end you can keep a search tree of a million keys at depth about $20$ no matter how they arrive.

\textbf{Full path.} Everything, in order -- it ends at the scaled Summit (LC Hard 493 \emph{Reverse
Pairs}: count $i<j$ with $\mathit{nums}[i]>2\cdot\mathit{nums}[j]$, where the guaranteed-$\bigO{\log
n}$ rank query of a balanced, augmented tree is the honest upgrade over Chapter~7's fragile descent,
and merge/Fenwick are the array-indexed rivals). The reduction gallery shows that ``AVL,'' ``treap,''
``red-black'' and ``splay'' are the same rotation-policy question with different answers, and that an
implicit treap is just a balanced BST whose keys are \emph{positions}. The forward rope is Chapter~9
(\chref{c26}, Cartesian trees and the monotonic stack) for the array$\to$tree bridge this chapter
deliberately leaves out, and Chapter~25 (\chref{c23}) for the full splay/link-cut machinery.
\end{chaptertheorybox}

\begin{figure}[H]
\centering
\begin{tikzpicture}[diagram, x=9mm, y=7mm]
  \foreach \st/\lab in {1/lv1,2/lv2,3/lv3,4/lv4,5/lv5,6/lv6} {
    \fill[\lab!22, draw=\lab, line width=0.8pt] (\st-1,0) rectangle (\st,\st*0.52);
    \node[font=\scriptsize\bfseries, text=\lab] at (\st-0.5,\st*0.52-0.42) {L\st};}
  \node[annot, anchor=south west, font=\scriptsize]
    at (0.05,0.62) {8.1 rotations};
  \node[annot, anchor=south west, font=\scriptsize]
    at (1.05,1.14) {8.1/8.2 the four rebalance cases};
  \node[annot, anchor=south west, font=\scriptsize]
    at (2.05,1.66) {8.2 AVL insert/delete cold};
  \node[annot, anchor=south west, font=\scriptsize]
    at (3.05,2.18) {8.3 treaps; split/merge; implicit keys};
  \node[annot, anchor=south west, font=\scriptsize, text=lv5]
    at (4.05,2.70) {8.4 red-black proof; 8.5 splay amortisation};
  \node[annot, anchor=south west, font=\scriptsize, text=lv5]
    at (5.05,3.22) {scaled summit (493)};
\end{tikzpicture}
\caption*{Chapter~8 on the six-step ladder: doing a rotation by hand ($\lv2$) climbs to a full AVL
insert/delete ($\lv3$); treaps and their split/merge engine and the red-black guarantee sit at
$\lv3$--$\lv4$; implicit-sequence treaps and splay amortisation reach $\lv4$; the Summit (493) makes
the controlled height pay for itself against an LC Hard, at $\lv5$. Beginner path $=$ steps 1--3.}
\end{figure}

% ---------------- The running examples ----------------
\subsection*{The Running Examples: A Chain That Must Be Saved, and a Case That Needs Two Turns}

This chapter has \emph{two} running examples, used over and over. The first is the villain: insert
the keys $1,2,\dots,7$ \emph{in increasing order} into an ordinary BST of Chapter~7 and watch it
collapse into a right-going chain of height $6$ -- search degrades to a scan, exactly the cliffhanger
that chapter ended on. The second is the interesting one: a violation that no \emph{single} rotation
can fix, which forces the LR double rotation. Every topic revisits them -- the rotation ($8.1$), AVL
rebalancing ($8.2$), treap priorities and split/merge ($8.3$), red-black recolouring ($8.4$), and the
splay access ($8.5$).

\begin{figure}[H]
\centering
\begin{tikzpicture}[diagram, scale=0.82, transform shape]
  % left: the chain C (ascending insert), height 6
  \node[tleaf]  (c1) at (0,0)     {$1$};
  \node[tleaf]  (c2) at (1.15,-0.9) {$2$};
  \node[tleaf]  (c3) at (2.3,-1.8)  {$3$};
  \node[tleaf]  (c4) at (3.45,-2.7) {$4$};
  \node[tleaf]  (c5) at (4.6,-3.6)  {$5$};
  \node[tleaf]  (c6) at (5.75,-4.5) {$6$};
  \node[tnode]  (c7) at (6.9,-5.4)  {$7$};
  \foreach \a/\b in {c1/c2,c2/c3,c3/c4,c4/c5,c5/c6,c6/c7}
    \draw[tedgehl] (\a) -- (\b);
  \node[annot, align=left, font=\scriptsize] at (1.4,0.9)
    {\textbf{Chain $C$}: insert $1,2,\dots,7$ in order\\ height $=6=n-1$\\ every descent is now
     $\bigO{n}$ -- the rot};
  % right: the balanced twin
  \node[troot]  (r4) at (12,0)   {$4$};
  \node[tnode]  (r2) at (10.4,-1.3) {$2$};
  \node[tnode]  (r6) at (13.6,-1.3) {$6$};
  \node[tleaf]  (r1) at (9.6,-2.6)  {$1$};
  \node[tleaf]  (r3) at (11.2,-2.6) {$3$};
  \node[tleaf]  (r5) at (12.8,-2.6) {$5$};
  \node[tleaf]  (r7) at (14.4,-2.6) {$7$};
  \foreach \a/\b in {r4/r2,r4/r6,r2/r1,r2/r3,r6/r5,r6/r7}
    \draw[tedgehl] (\a) -- (\b);
  \node[annot, align=left, font=\scriptsize] at (12,1.1)
    {\textbf{Balanced twin}: same seven keys\\ height $=2$\\ a rotation policy \emph{keeps} it here};
  \draw[flow, double] (7.6,-2.6) -- (9.2,-2.6) node[midway, above, font=\scriptsize] {balance};
\end{tikzpicture}
\caption{The villain, $C$: the \emph{same} set of keys stored as a degenerate chain (left, height $6$)
and as a height-balanced tree (right, height $2$). Nothing changes about the \emph{contents} or the
inorder line -- both read $1,2,3,4,5,6,7$ -- only the \emph{shape}. Balance is the art of moving from
left to right and \emph{staying} there, on every future edit, using the one legal move: the rotation.}\label{fig:c06-chain}
\end{figure}

\begin{figure}[H]
\centering
\begin{tikzpicture}[diagram, scale=0.86, transform shape]
  % Z before: 30 -> (10 -> (right 20 -> right 25)), 40
  \node[troot] (z30) at (0,0)     {$30$};
  \node[tnode] (z10) at (-1.6,-1.3) {$10$};
  \node[tnode] (z40) at (1.6,-1.3)  {$40$};
  \node[tnode] (z20) at (-0.5,-2.6) {$20$};
  \node[tleaf] (z25) at (0.4,-3.9)  {$25$};
  \draw (z30) -- (z10); \draw (z30) -- (z40);
  \draw (z10) -- (z20); \draw (z20) -- (z25);
  \node[annot, align=left, font=\scriptsize] at (-4.6,-2.6)
    {\textbf{$Z$, before}: insert $30,10,40,$\\ $20,25$. Node $30$ has\\ balance factor $+2$, its\\
     heavy child $10$ leans \emph{right}\\ $\Rightarrow$ the LR (zig-zag) case};
  % arrow
  \draw[flow, double] (3.0,-2.0) -- (4.6,-2.0) node[midway, above, font=\scriptsize, align=center] {rotate-left $10$,\\ then right $30$};
  % Z after: 20 -> 10,30 ; 30 -> 25,40
  \node[troot] (y20) at (7.2,-0.4) {$20$};
  \node[tleaf] (y10) at (5.8,-1.7) {$10$};
  \node[tnode] (y30) at (8.6,-1.7) {$30$};
  \node[tleaf] (y25) at (7.8,-3.0) {$25$};
  \node[tleaf] (y40) at (9.4,-3.0) {$40$};
  \draw (y20) -- (y10); \draw (y20) -- (y30);
  \draw (y30) -- (y25); \draw (y30) -- (y40);
  \node[annot, align=left, font=\scriptsize] at (7.6,0.7)
    {\textbf{$Z$, after}: height $2$, balanced\\ \emph{inorder unchanged}: $10,20,25,30,40$};
\end{tikzpicture}
\caption{The interesting case, $Z$: a Left-Right (zig-zag) violation. When the heavy child leans the
\emph{opposite} way, one rotation only converts the problem into a single-side case -- you need
\emph{two}: a rotation on the child first, then on the anchor. Both figures keep the identical sorted
inorder; this is the crux of \thmref{c06-rotation}.}\label{fig:c06-zigzag}
\end{figure}

\medskip
\noindent Keep the four \emph{strategies} below in mind as variations on one theme; each is proved or
used, never merely asserted, in the section named after it:
\begin{center}
\renewcommand{\arraystretch}{1.25}
\setlength{\tabcolsep}{5pt}
\begin{tabular}{@{}p{0.18\linewidth} p{0.44\linewidth} p{0.30\linewidth}@{}}
\toprule
\textbf{Scheme} & \textbf{How it keeps $h=\bigO{\log n}$} & \textbf{Where it lives} \\
\midrule
AVL          & strict rule: $\abs{\mathit{BF}}\le 1$ per node; rotate on every edit that breaks it & \secref{c06-avl} \\
Treap        & random priority per node; BST on keys \emph{and} heap on priorities is short \emph{in expectation} & \secref{c06-treaps} \\
Red-black    & colour invariants; longest path $\le 2\times$ shortest $\Rightarrow h\le 2\log_2(n{+}1)$ & \secref{c06-redblack} \\
Splay        & no stored rule: every access rotates to the root; amortised $\bigO{\log n}$ & \secref{c06-splay} \\
\bottomrule
\end{tabular}
\end{center}

% ---------------- Roadmap diagram ----------------
\begin{figure}[H]
\centering
\begin{tikzpicture}[diagram, node distance=3.5mm,
    top/.style={boxnode, fill=primaryblue!10, draw=primaryblue, text width=19mm,
                minimum height=10mm, font=\small},
    res/.style={annot, text width=19mm, anchor=north, font=\scriptsize}]
  \node[top] (t1) {\textbf{8.1} Why \& Rotations};
  \node[top, right=of t1] (t2) {\textbf{8.2} AVL};
  \node[top, right=of t2] (t3) {\textbf{8.3} Treaps};
  \node[top, right=of t3] (t4) {\textbf{8.4} Red-Black};
  \node[top, right=of t4] (t5) {\textbf{8.5} Splay};
  \foreach \a/\b in {t1/t2, t2/t3, t3/t4, t4/t5} \draw[flow] (\a) -- (\b);
  \node[res] at ([yshift=-2mm]t1.south) {the one legal move; four cases};
  \node[res] at ([yshift=-2mm]t2.south) {balance factors; $\le2$ turns fix};
  \node[res] at ([yshift=-2mm]t3.south) {random priorities; split/merge};
  \node[res] at ([yshift=-2mm]t4.south) {the std::set guarantee};
  \node[res] at ([yshift=-2mm]t5.south) {self-adjust; amortised};
\end{tikzpicture}
\caption*{The five topics are one question at increasing abstraction. 8.1 gives the \emph{atomic
move}; 8.2 uses it \emph{greedily and provably} (AVL); 8.2 hands 8.3's treap the same job via
\emph{chance} instead of case work, and the split/merge engine; 8.4 explains the \emph{relaxed} rule
the standard library actually ships; 8.5 drops the stored invariant entirely for an \emph{amortised}
self-adjusting tree that anticipates the dynamic-trees chapter.}
\end{figure}

% ---------------- Topic table ----------------
\begin{center}
\renewcommand{\arraystretch}{1.25}
\setlength{\tabcolsep}{4pt}
\begin{tabular}{@{}p{0.235\linewidth} p{0.365\linewidth} c p{0.23\linewidth}@{}}
\toprule
\textbf{Topic} & \textbf{Key results} & \textbf{Levels} & \textbf{Solved problem} \\
\midrule
\secref{c06-rotations}\ \ Why Balance? The Rotation
  & op cost $=h$; the rotation preserves inorder; LL/RR/LR/RL cases & $\lv2$--\lv{3}
  & \probref{c06-treecons} \\
\addlinespace
\secref{c06-avl}\ \ AVL Trees
  & balance factor; Fibonacci height bound; insert/delete with $\le2$ turns & $\lv3$--\lv{4}
  & \probref{c06-avl} \\
\addlinespace
\secref{c06-treaps}\ \ Treaps, Split/Merge, Implicit
  & BST+heap on random priorities; $\bigO{\log n}$ expected; split/merge kit & $\lv4$--\lv{5}
  & \probref{c06-rpmq} \\
\addlinespace
\secref{c06-redblack}\ \ Red-Black Trees (conceptual)
  & five colour invariants; $h\le 2\log_2(n{+}1)$ -- the std::set promise & $\lv3$--\lv{4}
  & (references \chref{c05}) \\
\addlinespace
\secref{c06-splay}\ \ Splay Trees (intro)
  & move-to-root; zig/zig-zig/zig-zag; amortised $\bigO{\log n}$ & $\lv4$--\lv{5}
  & (intro; full $\to$ \chref{c23}) \\
\bottomrule
\end{tabular}
\end{center}

% ---------------- Problem recognition guide ----------------
\subsection*{Problem Recognition Guide --- When You See \emph{this}, Think \emph{that}}

Balancing problems announce themselves when a \emph{plain BST or a hand-rolled ordered structure
must survive adversarial order}, or when you need a \emph{dynamic sequence} with range operations.
Match the phrasing; every row points at the section that builds the tool.

\begin{center}
\renewcommand{\arraystretch}{1.3}
\setlength{\tabcolsep}{4pt}
\begin{tabular}{@{}p{0.44\linewidth} p{0.34\linewidth} l@{}}
\toprule
\textbf{When the problem says\ldots} & \textbf{\textcolor{deepblue}{\textbf{Think\ldots}}} & \textbf{Section} \\
\midrule
``insert these keys; the tree must stay efficient'', keys arrive \emph{sorted}
  & plain BST $=$ a chain -- balance by construction & \secref{c06-rotations} \\
\addlinespace
``implement a self-balancing BST / AVL'', ``rebalance after insert''
  & track balance factors up the path; the four cases & \secref{c06-avl} \\
\addlinespace
``cut/paste a sequence, reverse a range, $k$th by position''
  & implicit-key treap; \code{split}/\code{merge} & \secref{c06-treaps} \\
\addlinespace
``rank / count-smaller / order statistic, online, worst-case fast''
  & size-augmented \emph{balanced} tree (or Fenwick) & \secref{c06-treaps}, Summit \\
\addlinespace
``why is \code{std::set} $\bigO{\log n}$?'' / a guarantee you must prove
  & red-black colour invariants & \secref{c06-redblack} \\
\addlinespace
``frequently accessed items should be cheap / no stored balance''
  & splay: move-to-root, amortised & \secref{c06-splay} \\
\bottomrule
\end{tabular}
\end{center}

\begin{center}
\renewcommand{\arraystretch}{1.25}
\setlength{\tabcolsep}{4pt}
\begin{tabular}{@{}p{0.27\linewidth} p{0.45\linewidth} p{0.20\linewidth}@{}}
\toprule
\textbf{Constraint you see} & \textbf{What it forces} & \textbf{Chapter-8 tool} \\
\midrule
keys arrive in (near-)sorted order & never trust a plain BST's height & \secref{c06-avl} \\
an ordered set/map needed under edits & reach for \code{set}/\code{map}, know it is red-black & \secref{c06-redblack} \\
a \emph{sequence} with range reverse/sum & implicit treap, keys are positions & \secref{c06-treaps} \\
rank/count queries, worst-case $\bigO{\log n}$ & augment with subtree size on a balanced tree & Summit (\probref{c06-summit}) \\
dense integer keys, array-index ranks & Fenwick / segment beat a pointer tree & \chref{c14}, \chref{c13} \\
\bottomrule
\end{tabular}
\end{center}

\begin{notebox}[Not this chapter?]
If you just need the \emph{ordered descent itself} -- search, insert, delete, kth-smallest, the
\code{std::set} API, interval management -- that is Chapter~7 (\chref{c05}); this chapter assumes all
of it and only asks \emph{how the height stays small}. If the tree must support \emph{link/cut of
whole subtrees} and path queries, the full splay and link-cut machinery is Chapter~25
(\chref{c23}) -- \secref{c06-splay} is only its prerequisite intro. If you want the
\emph{array$\to$tree bridge} -- Cartesian trees, the monotonic stack, largest rectangle, subarray
minimums, next-smaller links -- that is Chapter~9 (\chref{c26}), which this chapter deliberately
leaves out. And if keys are dense integers and you only need ranks by index, the Fenwick tree
(\chref{c14}) and segment tree (\chref{c13}) are faster engines for the very same sums. This chapter
owns \emph{keeping the ordered tree short by construction, and what that guaranteed height buys you}.
\end{notebox}

% ---------------- Notation ----------------
\subsection*{Notation Used in This Chapter}

\begin{center}
\renewcommand{\arraystretch}{1.2}
\setlength{\tabcolsep}{4pt}
\begin{tabular}{@{}p{0.20\linewidth} p{0.27\linewidth}@{\quad}p{0.17\linewidth} p{0.27\linewidth}@{}}
\toprule
\textbf{Symbol} & \textbf{Meaning} & \textbf{Symbol} & \textbf{Meaning} \\
\midrule
$h$                 & tree height; every descent costs $\bigO{h}$, balance forces $h=\bigO{\log n}$ & $\bal(v)$ & balance factor of $v$, $\height(L)-\height(R)$ \\
$C$                 & the degenerate chain (\figref{c06-chain})   & $Z$ & the zig-zag instance (\figref{c06-zigzag}) \\
$\mathit{p}(v)$     & random priority of $v$ in a treap (heap order) & $\mathit{size}(v)$ & node count of $v$'s subtree (augmentation) \\
$\mathrm{rot}_{L},\ \mathrm{rot}_{R}$ & left / right rotation about a node & $\mathrm{size}$-BF & implicit key $=$ left-subtree size \\
\code{set}, \code{map} & the \texttt{std} containers -- red-black under the hood (\secref{c06-redblack}) & $\lv{3}$ & difficulty badge, one of $\lv{1}$--$\lv{6}$ \\
\bottomrule
\end{tabular}
\end{center}

% !TEX root = ../../tree_notes.tex
% =====================================================================
%  8.1  Why Balance? The Rotation                sec:c06-rotations
%  Exports: def:c06-heightcost  thm:c06-rotation  def:c06-rotcases
%           fig:c06-rot  lst:c06-rotcode  ex:c06-rot-run
%           prob:c06-treecons  lst:c06-treecons-soln   (practice: LC 1382)
% =====================================================================
\section{Why Balance? The Rotation}\label{sec:c06-rotations}

\begin{previouslybox}
\begin{enumerate}[nosep, leftmargin=1.6em]
  \item What is the BST invariant, and why does one comparison \emph{discard a whole subtree} during
        a search? (\chref{c05}, \defref{c05-prop}, \propref{c05-search})
  \item Every Chapter-7 routine -- search, insert, the three-case delete -- costs how much, written
        in which single parameter? (\bigO{h}; \secref{c05-insert-delete})
  \item Insert $1,2,\dots,n$ into a plain BST. What shape do you get, and what does $h$ become?
        (a chain, $h=n-1$; \secref{c05-contests}' warning, \probref{c05-summit}' cliffhanger)
  \item Which traversal of a BST lists the keys in sorted order, and why will we lean on it being
        \emph{invariant under reshaping}? (\chref{c03}, \thmref{c05-inorder})
\end{enumerate}
\end{previouslybox}

\begin{topicintro}
Chapter~7 proved a beautiful thing and then quietly made you pay for it: a BST answers
membership, insertion, deletion and every ordered query in $\bigO{h}$, \emph{where $h$ is the height}
-- and $h$ is a property of the \emph{input order}, not of you. Insert sorted keys and the tree is a
linked list with spare pointers; the Summit \probref{c05-summit} went quadratic for exactly this
reason. This section installs the fix. There is only \emph{one} legal local operation on a binary
search tree -- the \term{rotation} -- and its magic is that it changes the shape (and hence the
height) while leaving the sorted inorder line \emph{byte for byte identical}. Once you accept the
rotation as the atomic move, every balanced tree in the rest of the chapter is a \emph{policy}: a
rule for when to rotate and how hard to keep things short. We finish on Codeforces 675D, which hands
you a degenerate insertion order and asks for each node's parent -- the cleanest possible way to feel
both the $\bigO{n^2}$ trap and the $\bigO{n\log n}$ escape.
\textbf{Prerequisites.} All of Chapter~7 (\chref{c05}); the descent, the invariant, and the
$\bigO{h}$-cost warning are assumed. Inorder from Chapter~3 (\chref{c03}) is used as a black box,
never re-derived.

\textbf{After this section you can}
\begin{itemize}[nosep]
  \item state precisely why every BST operation is $\bigO{h}$ and why $h=\bigO{\log n}$ is the goal; \hfill\lv{2}
  \item perform a left and a right rotation and prove the inorder sequence is unchanged; \hfill\lv{2}--\lv{3}
  \item classify a height violation into LL / RR / LR / RL and pick the one- or two-rotation fix; \hfill\lv{3}
  \item solve BST-construction-under-degenerate-order problems with a set, escaping $\bigO{n^2}$. \hfill\lv{4}
\end{itemize}
\end{topicintro}

\begin{recognitionbox}[Pattern card -- the shape is the problem]
\reckey{a search tree's \emph{efficiency}, not its answer, depends on insertion order; or you are
asked to \emph{rebuild/reshape} a BST while preserving its sorted contents.}
\recuse{
  \item ``insert these keys, tree must stay fast'' / keys arrive sorted -- balance by construction
        (the whole chapter; start \secref{c06-avl});
  \item ``give each node's parent in a BST built by this order'' -- CF 675D, the set-neighbour trick
        (\probref{c06-treecons});
  \item ``turn a skewed BST into a balanced one'' -- rebuild from inorder (Practice, LC 1382);
  \item any rotation-tracing or ``apply one left rotation at $x$'' warm-up -- \defref{c06-rotcases}.
}
\recnotuse{
  \item the tree is \emph{already} balanced and you just query it -- that is plain Chapter~7 descent
        (\chref{c05});
  \item keys are dense integers and you only need prefix ranks -- Fenwick \chref{c14} skips trees
        entirely;
  \item you need whole-subtree link/cut, not single-node rebalancing -- splay/link-cut, \chref{c23}.
}
One-line test: \emph{is the \textbf{cost} of a correct answer the thing at stake -- a height you must
force down without disturbing the sorted order? Then you need rotations, and this chapter.}
\end{recognitionbox}

\subsection{Intuition}

Picture the running chain $C$ of \figref{c06-chain}: keys $1,2,\dots,7$ inserted in increasing order.
Every new key is larger than everything present, so it always goes right -- a straight diagonal of
seven nodes, height $6$. A search for $7$ now touches all seven nodes: the tree has \emph{forgotten}
how to binary-search and degenerated into a linked list. Yet the \emph{contents} are perfectly fine --
read the chain's inorder (left, node, right) and you still get $1,2,3,4,5,6,7$, sorted. The problem is
purely geometric: same data, bad shape.

So we want a move that is to a BST what \emph{reshuffling a bookshelf} is to a library: change which
shelf is the central support without re-ordering a single book. That move is the \term{rotation}. It
re-parents \emph{three} pointers around one node, so its cost is $\bigO{1}$, and -- crucially --
\thmref{c06-rotation} guarantees it cannot disturb the inorder line. Do enough of them and the chain
$C$ becomes the balanced twin on the right of \figref{c06-chain}: height drops from $6$ to $2$, every
descent is $\bigO{\log n}$ again, and \emph{not one key moved in the sorted order}.

\begin{figure}[H]
\centering
\begin{tikzpicture}[diagram, scale=0.95, transform shape]
  % left rotation about x:  before
  \node[tnode] (xa) at (0,0) {$x$};
  \node[tleaf] (aA) at (-1.2,-1.3) {$A$};
  \node[tnode] (ay) at (1.2,-1.3) {$y$};
  \node[tleaf] (aB) at (0.3,-2.6) {$B$};
  \node[tleaf] (aC) at (2.1,-2.6) {$C$};
  \draw (xa) -- (aA); \draw (xa) -- (ay); \draw (ay) -- (aB); \draw (ay) -- (aC);
  \node[annot, font=\scriptsize] at (0.5,1.0) {left rotate about $x$};
  \node[annot, font=\scriptsize\bfseries] at (0,-4.0) {before};
  % arrow
  \draw[flow, double] (3.4,-1.4) -- (5.0,-1.4);
  % after
  \node[tnode] (yb) at (6.6,0) {$y$};
  \node[tnode] (bx) at (5.0,-1.3) {$x$};
  \node[tleaf] (cC) at (8.0,-1.3) {$C$};
  \node[tleaf] (bA) at (4.0,-2.6) {$A$};
  \node[tleaf] (bB) at (6.0,-2.6) {$B$};
  \draw (yb) -- (bx); \draw (yb) -- (cC); \draw (bx) -- (bA); \draw (bx) -- (bB);
  \node[annot, font=\scriptsize\bfseries] at (6.5,-4.0) {after};
  \node[annot, align=left, font=\scriptsize] at (8.6,-2.4)
    {inorder both:\\ $A,x,B,y,C$\\ \emph{identical} --\\ only $x,y,B$ re-linked};
\end{tikzpicture}
\caption{A left rotation about $x$ (right child $y$): $y$ is promoted to $x$'s former spot, $x$
becomes $y$'s left child, and $y$'s old left subtree $B$ is handed down to become $x$'s right child.
Exactly three subtree pointers move; the inorder sequence $A,x,B,y,C$ is preserved verbatim, which is
\thmref{c06-rotation}. A right rotation is the exact mirror image.}\label{fig:c06-rot}
\end{figure}

\subsection{Definitions and Notation}

\begin{definition}{Height is the whole cost}{c06-heightcost}
For a BST of $n$ distinct keys and height $h$, every one-descent operation of Chapter~7 --
membership, the failed-search insert, the successor-splice delete, a kth-smallest walk -- costs
$\bigO{h}$. A tree is \term{height-balanced} if there is a constant $c$ with $h \le c\,\lg(n+1)$
regardless of insertion order. Because $2 \le h+1 \le n$ always, ``balanced'' means the descent is
guaranteed logarithmic, and \emph{un}balanced means it can be linear: the entire difference between
the Summit's $\bigO{n^2}$ failure and a $\bigO{n\log n}$ success.
\end{definition}

\noindent\textbf{In plain words.} Stop thinking ``$\bigO{n}$'' for a search. Think ``$\bigO{h}$'',
and make keeping $h$ tiny the \emph{data structure's job}, so the algorithm never has to.

\begin{theorem}{A rotation preserves the inorder sequence}{c06-rotation}
Let $x$ be a node with right child $y$, left subtree $A$, and $y$'s left subtree $B$, right subtree
$C$ (so the subtree's inorder is $A,x,B,y,C$). The left rotation about $x$ producing $y\to(x\to(A,B),C)$
yields the \emph{same} inorder sequence $A,x,B,y,C$, and its inverse (right rotation about $y$) restores
the original. Symmetrically for a right rotation. Hence rotations preserve both the \emph{sorted
order} and the BST property of the affected subtree.
\end{theorem}
\begin{proofbox}
Read the inorder of the \emph{after} picture. The new root is $y$; its left subtree is rooted at $x$,
whose own children are $A$ (left) and $B$ (right), and whose right child under $y$ is $C$. Visit $y$'s
left subtree first: within $x$, visit $A$, then $x$, then $B$ -- giving $A,x,B$. Then visit $y$. Then
visit $y$'s right subtree $C$. Total: $A,x,B,y,C$ -- term for term the \emph{before} inorder. No key
crossed any other, so the BST order (every subtree a contiguous sorted run) is intact on both sides.
Undoing is the mirror right rotation about $y$, which we just showed preserves inorder too, so the two
operations are exact inverses. \end{proofbox}
\begin{insightbox}
\thmref{c06-rotation} is the \emph{licence to reshape}. It says geometry and semantics are decoupled
in a BST: the sorted line is fixed, and the \emph{tree that realises it} is a choice. Height is a
property of the realisation, not the data -- so we are free to re-choose it (rotate) whenever we
like, and every balancing scheme in Sections 8.2--8.5 is exactly a rule for when to spend a rotation.
\end{insightbox}

\begin{definition}{The four violation cases: LL, RR, LR, RL}{c06-rotcases}
A node $v$ is \term{unbalanced} when its two subtrees differ in height too much (the precise threshold
defines AVL vs red-black vs nothing; here ``one side is a chain and the other short'' suffices). Look
at the heavy child and which side \emph{it} leans:
\begin{itemize}[nosep]
  \item \textbf{LL} -- heavy child leans Left: one \emph{right} rotation about $v$ fixes it.
  \item \textbf{RR} -- heavy child leans Right: one \emph{left} rotation about $v$.
  \item \textbf{LR} -- heavy child leans Right (zig-zag): first a left rotation about the child turns
        it into LL, then a right rotation about $v$. (This is $Z$, \figref{c06-zigzag}.)
  \item \textbf{RL} -- heavy child leans Left (zig-zag): first right about the child, then left about $v$.
\end{itemize}
The single-rotation cases are \emph{zig-zig} (both bends the same way); the double are \emph{zig-zag}
(bends differ), needing exactly two rotations and three relinkings in total.
\end{definition}

\noindent\textbf{In plain words.} A same-direction lean is one rotation; an opposite-direction lean
(the child sticks out the far side) is \emph{two}, because the first rotation converts the zig-zag
into a zig-zig, then the second finishes. Node $Z$ needs two; you cannot fix it with one, which is
the whole lesson of \figref{c06-zigzag}.

\subsection{The algorithm: one rotation, written once}

Every balancing scheme reuses the same two four-line routines. Keep pointers to the three moving
subtrees ($A,B,C$), relink twice, and recompute the two affected nodes' heights -- done. That
$\bigO{1}$ cost is why a single rotation is ``free,'' and why the schemes differ only in
\emph{decision}, never in \emph{mechanics}.

\begin{lstlisting}[caption={Left and right rotations about a node -- the atomic move reused by AVL, red-black and splay}, label={lst:c06-rotcode}]
// Left rotation about x (x must have a right child y). Returns the new subtree root y.
Node* rotateLeft(Node* x) {
    Node* y = x->right;          // y is promoted
    x->right = y->left;          // B moves under x as its right child
    y->left  = x;                // x becomes y's left child
    updHeight(x); updHeight(y);  // x changed first, then its new parent y
    return y;                    // y is the root of the reshaped subtree
}
// Right rotation about y (mirror image; y must have a left child x). Returns new root x.
Node* rotateRight(Node* y) {
    Node* x = y->left;
    y->left  = x->right;         // B moves under y as its left child
    x->right = y;
    updHeight(y); updHeight(x);
    return x;
}
\end{lstlisting}
\begin{complexitybox}
Time $\bigO{1}$ -- a fixed number of pointer assignments and two height recomputes; memory $\bigO{1}$.
This is the point: rotations are constant-time, so a balancing policy that performs a \emph{constant
number} of them per edit (AVL: $\le 2$; red-black: $\le 2$ rotations plus recolouring; splay: $O(1)$
amortised) keeps every Chapter-7 operation at its $\bigO{\log n}$ ideal. Verdict: the atomic
$\bigO{1}$ move \emph{toward} balance.
\end{complexitybox}

\subsection{Worked Example}

\begin{example}{Rotating the chain $C$ and resolving the zig-zag $Z$}{c06-rot-run}
\emph{Chain $C$, top triple.} Focus on the head $1\to 2\to 3$ (each going right) of $1,2,\dots,7$.
Apply \texttt{rotateLeft(1)}: $2$ is promoted, $1$ becomes $2$'s left child, and $2$'s (empty) left
subtree $B$ becomes $1$'s right. Local shape $1\to 2\to(\text{nothing},3)$ becomes $2\to(1,3)$ -- a
height-$1$ balanced triple instead of a height-$2$ chain. Repeat on the next triple and, after
$\bigO{1}$ rotations per level, the diagonal flattens toward the balanced twin of
\figref{c06-chain}. Inorder is untouched: at every step the run reads $1,2,3,\dots$.\\
\emph{Zig-zag $Z$.} At node $30$ the left subtree ($10\to 20\to 25$, all going right) has height $2$
while the right ($40$) has height $0$: an LL-vs-RL mismatch -- specifically the heavy left child $10$
leans \emph{right}, the \textbf{LR} case. One rotation cannot fix it. First \texttt{rotateLeft(10)}:
$20$ is promoted, giving $30\to(20\to(10,25),40)$ -- now the heavy child $20$ leans \emph{left} at
$30$'s subtree, a clean LL. Then \texttt{rotateRight(30)}: $20$ reaches the top, final shape
$20\to(10, 30\to(25,40))$, height $2$, balanced. Inorder $10,20,25,30,40$ before and after --
\thmref{c06-rotation} honoured twice.
\end{example}
\begin{tracebox}[The two rotations on $Z$, pointer by pointer]
\textbf{Start:} $30$'s left $=10$, $10$'s right $=20$, $20$'s right $=25$, $30$'s right $=40$.
\textbf{rotateLeft(10):} $B := 20$'s left (null). Set $10\to$right $=B=$ null; $20\to$left $=10$.
Now $30\to$left $=20$, $20\to$left $=10$, $20\to$right $=25$. Heavy child of $30$ is now a
\emph{left-leaning} node ($20$ with left child $10$) -- the LL case.
\textbf{rotateRight(30):} $B := 20$'s right ($=25$). Set $30\to$left $=B=25$; $20\to$right $=30$.
Now root $=20$, $20\to$left $=10$, $20\to$right $=30$, $30\to$left $=25$, $30\to$right $=40$. Done:
height $2$, both subtrees of $20$ height $1$. \checkmark Two rotations, inorder $10,20,25,30,40$
preserved -- see \figref{c06-zigzag}.
\end{tracebox}

\subsection{Variations and Twists}

\begin{itemize}[nosep]
  \item \textbf{Double rotation as two singles}: LR $=$ left-then-right and RL $=$ right-then-left
        factor into two single rotations about adjacent nodes -- implement \code{rotateLeft} /
        \code{rotateRight} once and compose; never hand-code a ``double.'' \hfill\lv{3}
  \item \textbf{Rebuild instead of rotate}: dump the BST's inorder into an array and build a balanced
        tree by recursive halving -- the LC 1382 approach; $\bigO{n}$ with no rotations at all.
        \hfill\lv{3}
  \item \textbf{Rotation distance}: any two BSTs on the same keys are connected by $\bigO{n}$
        rotations (Catalan-flavoured; a research-grade fact, cited not used here). \hfill\lv{5}
  \item \textbf{Rotations preserve subtree \emph{sets}}: the three subtrees $A,B,C$ keep their exact
        key contents; only the parentage changes -- this is why an augmented field like
        $\mathit{size}()$ needs \emph{recomputing} on the two moved nodes, not a rebuild. \hfill\lv{4}
\end{itemize}

\subsection{Common Pitfalls}\label{sec:c06-rotations-pitfalls}

\begin{warningbox}[Trap 1 -- rotating the wrong way relative to the lean]
A heavy-left subtree (LL) needs a \emph{right} rotation; a common slip is to left-rotate and deepen
the chain further. Decide the direction from the side that is \emph{too tall}: pull the tall side up.
If after a rotation the height went the wrong way, you rotated about the wrong node or side.
\end{warningbox}
\begin{lstlisting}[caption={\textcolor{warningred}{WRONG}: left-rotate an LL (heavy-left) node -- deepens the left chain}, language=C++, basicstyle=\ttfamily\small]
// v's LEFT subtree is too tall (LL). rotateLeft(v) requires a right child and
// makes the LEFT side even heavier -- the opposite of the fix you need.
Node* badFix(Node* v) { return rotateLeft(v); }   // WRONG direction for LL
\end{lstlisting}
\begin{lstlisting}[caption={\textcolor{successgreen}{RIGHT}: LL is repaired by one right rotation about the anchor}, language=C++, basicstyle=\ttfamily\small]
Node* fixLL(Node* v) { return rotateRight(v); }     // pull the heavy LEFT child up
// LR (zig-zag): first straighten the child, then rotate the anchor.
Node* fixLR(Node* v) { v->left = rotateLeft(v->left); return rotateRight(v); }
\end{lstlisting}

\begin{warningbox}[Trap 2 -- updating heights in the wrong order]
After \texttt{rotateLeft(x)} the node $x$ ends up \emph{below} $y$. If you recompute $y$'s height
before $x$'s, $y$ reads a stale child height. Update the node that moved \emph{down} first, then the
one promoted: exactly the order in \code{rotateLeft} (\coderef{c06-rotcode}). A stale balance
factor silently turns AVL into the same old chain.
\end{warningbox}
\begin{lstlisting}[caption={\textcolor{warningred}{WRONG}: promoted node recomputed before its new child}, language=C++, basicstyle=\ttfamily\small]
x->right = y->left; y->left = x;
updHeight(y);            // y still sees x's OLD height here -- WRONG
updHeight(x);
\end{lstlisting}
\begin{lstlisting}[caption={\textcolor{successgreen}{RIGHT}: fix the lower node $x$ first, then its parent $y$}, language=C++, basicstyle=\ttfamily\small]
x->right = y->left; y->left = x;
updHeight(x);            // x's children are final
updHeight(y);            // now reads x's corrected height -- RIGHT
\end{lstlisting}

\begin{warningbox}[Trap 3 -- simulating degenerate insertion with an $\bigO{h}$ descent]
``Insert $n$ keys into a BST and report each parent'' done by walking down the tree is
$\bigO{h}$ per key, i.e.\ $\bigO{n^2}$ on the chain $C$ -- and CF 675D's $n\le 10^{5}$ will
\tle. The escape (\probref{c06-treecons}) needs \emph{no tree at all}: a \code{std::set} finds each
key's neighbours in $\bigO{\log n}$ and the parent falls out of their insertion times.
\end{warningbox}

\subsection{Problem Lab}

The section's anchor is Codeforces 675D -- build a BST by a given insertion order and report each
node's parent. Simulated honestly it is the degenerate-height trap in the raw ($\bigO{n^2}$); solved
cleverly it is a two-line set query. It is chosen precisely because it \emph{is} the reason this
chapter exists.

% ---------------------------------------------------------------------------
\solvedproblem{c06-treecons}{Tree Construction}{Codeforces 675D}
  {https://codeforces.com/problemset/problem/675/D}{CF 1700 · L4}
  {BST insertion order, set-based neighbour lookup, degenerate-height escape, later-inserted parent}

\begin{problemstatement}[Tree Construction]
You are given a sequence $a_1,\dots,a_n$ of $n$ \textbf{distinct} integers, the order in which they
are inserted into an initially empty binary search tree (each inserted by the usual
$\bigO{h}$ descent of Chapter~7). For every element \emph{except the first}, output the value of its
\textbf{parent} node in the finished BST. $3\le n\le 10^{5}$ and $1\le a_i\le 10^{9}$. Output the
parents one per line, in the order the corresponding children were inserted.
\end{problemstatement}

\begin{attemptgate}[Attempt it yourself first -- \lv{4}, give yourself 25 minutes]
The honest simulation (actually build the tree) is correct but $\bigO{n^2}$ on sorted input -- the
chain $C$ again. The insight is to \emph{not} build the tree.
\peek{When you insert a fresh key $x$ into a BST, its parent is an \emph{already-present} key. Which
two candidates could it possibly be, in terms of the sorted order of the keys already inserted?}
\nudge{$x$'s parent must be its predecessor or its successor \emph{among the keys already present}
(the \code{std::set} can find both in $\bigO{\log n}$). The one inserted \emph{later} is deeper and
becomes the parent -- the earlier one already had the later on the path.}
\fullhint{Keep a \code{set} of present keys and a \code{map} key$\to$insertion-index \code{when}. For
$x$: \code{auto r = s.lower\_bound(x)}; the successor is \code{*r} (when \code{r} is not \code{s.end()}), the predecessor is
\code{*prev(r)} (when \code{r} is not \code{s.begin()}). Parent $=$ whichever of them has the \emph{larger} \code{when}.
Print it for $i\ge 2$, then \code{s.insert(x)} and set \code{when[x]=i}. Every step is
$\bigO{\log n}$, total $\bigO{n\log n}$ -- immune to the degenerate order that sinks the simulation.}
\end{attemptgate}

\solstep{Understanding the Problem}
We must report parent pointers of a BST defined purely by insertion order. The naive route walks the
real tree down from the root -- but on sorted data that tree is the chain $C$, so each walk is
$\bigO{n}$ and the whole is $\bigO{n^2}$, too slow for $n=10^{5}$. The reframing: the finished shape
is a \emph{function of the sorted contents and the arrival times}, not of any pointer we must build.
A new key $x$ hangs as a leaf; its parent is the last-inserted of the two keys that $x$ sits
\emph{between} in sorted order. So we only need, for each $x$, its predecessor and successor among the
already-inserted keys and their insertion order -- a job for \code{std::set} in $\bigO{\log n}$
(\secref{c05-contests}).

\solstep{Key Observations}
\begin{observation}
$x$'s parent is one of exactly two already-present keys: the largest key $<x$ (predecessor) or the
smallest key $>x$ (successor). Any other key is separated from $x$ by one of these two and could not
be the leaf's direct parent.
\end{observation}
\begin{observation}
Among the predecessor $p$ and successor $s$, the parent is whichever was inserted \textbf{later}.
The earlier one had the later already on the descent path, so $x$ attaches below the later one.
Equivalently: at every node, the child on the side $x$ belongs to was decided last among the two
neighbours.
\end{observation}
\begin{observation}
Both neighbours and their insertion times are available from a sorted \code{set} plus a
\code{map} key$\to$time in $\bigO{\log n}$ -- \emph{no tree is built}, so the degenerate shape never
costs us. This is the escape from the $\bigO{n^2}$ simulation.
\end{observation}

\solstep{Think It Through With Me}
\begin{thinkbox}
\deadend{Actually construct the BST with \secref{c05-insert-delete} insertions and print parents.}
{Correct but each insertion is an $\bigO{h}$ descent, and $h$ on sorted input is $n$, giving
$\bigO{n^2}$ -- the chain $C$ made manifest. For $n=10^{5}$ that is $\sim 10^{10}$ operations, a
\tle. The tree's \emph{bad shape} is exactly what we are trying to avoid; here we can avoid the tree.}
\deadend{Sort the array and answer neighbours by position.}
{A plain sort gives global neighbours but throws away \emph{insertion time}, and the parent depends on
which neighbour arrived \emph{later}. Sort order alone picks the wrong parent about half the time --
you need a structure that grows online, one key at a time.}
\aha{Maintain the present keys in a \code{std::set} (sorted, $\bigO{\log n}$ neighbours) and remember
each key's insertion index. $x$'s parent $=$ the predecessor or successor whose index is \emph{larger}.
One \code{lower\_bound} plus a step-back per key, $\bigO{n\log n}$ total, no tree, no degeneration.}
\end{thinkbox}

\solstep{Solution}
Iterate $i=1..n$ in insertion order. Before handling $a_i$, the \code{set} $s$ holds exactly the keys
already inserted (smaller than $i$ in time) and \code{when} maps each to its index. Compute
\code{r = s.lower\_bound(x)}. Candidate parent is \code{*r} (successor) when \code{r} is not
\code{s.end()}, and \code{*prev(r)} (predecessor) when \code{r} is not \code{s.begin()}; pick the
candidate with the larger \code{when}, which
is the later-inserted and hence the true parent. For $i\ge 2$ print that parent. Then insert $x$ into
$s$ and record \code{when[x]=i}. The root ($i=1$) prints nothing.

\solstep{Proof of Correctness}
\begin{proofbox}[The later-inserted neighbour is the parent]
Consider the moment $x$ is inserted and let $p<x<q$ be its predecessor and successor among the
already-present keys (one may be absent at an end). In the final BST every key of the tree lies in
exactly one of the open intervals determined by the sorted keys, and $x$'s search path is governed by
these: descending, the path is steered by $p$ and $q$ until it reaches whichever of them was inserted
last. Suppose $q$ was inserted after $p$ (the other case is symmetric). At the time $q$ arrived, $p$
was already present with $p<q$, so $p$ is in $q$'s left subtree; because no key lies strictly between
$p$ and $q$ (they are adjacent in the current sorted set), nothing between them can separate $x$ from
$q$, and $x$ (with $p<x<q$) descends at $q$ to the left child slot that is still empty -- hence $x$
becomes $q$'s left child, so $q$ is the parent. Conversely the earlier-inserted $p$ already had $q$
on the descent path below it, so it cannot be $x$'s direct parent. Therefore the parent is precisely
the neighbour with the greater insertion index, which the algorithm reports. The root ($i=1$) has no
neighbours and correctly prints nothing. \end{proofbox}

\solstep{Implementation}
\begin{lstlisting}[caption={Tree Construction -- set-based parents without building the tree (Codeforces 675D)}, label={lst:c06-treecons-soln}]
#include <bits/stdc++.h>
using namespace std;
int main(){
    ios::sync_with_stdio(false); cin.tie(nullptr);
    int n; cin >> n;
    vector<int> a(n);
    for (int i = 0; i < n; ++i) cin >> a[i];

    set<int> s;             // keys already inserted, kept sorted
    map<int,int> when;      // key -> insertion index (its "time")
    for (int i = 0; i < n; ++i) {
        int x = a[i];
        int parent = -1;
        auto r = s.lower_bound(x);            // first present key > x (distinct => never == x)
        // successor candidate *r, predecessor candidate *prev(r): pick the LATER-inserted one
        if (r != s.end())           parent = *r;
        if (r != s.begin()) {
            int p = *prev(r);                 // greatest key < x
            if (parent == -1 || when[p] > when[parent]) parent = p;
        }
        if (i > 0) cout << parent << '\n';    // root (i==0) has no parent
        s.insert(x);
        when[x] = i;
    }
    return 0;
}
\end{lstlisting}

\solstep{Trace}
\begin{tracebox}[Official-style sample: insert order $2, 1, 3, 5, 4$]
$i=0$, $x=2$: empty set, root, print nothing. $s=\{2\}$, \code{when[2]=0}.\\
$i=1$, $x=1$: \code{lower\_bound(1)}=iter to $2$ (successor). No predecessor ($r$ is begin). Parent
$=2$. Print $2$. $s=\{1,2\}$, \code{when[1]=1}.\\
$i=2$, $x=3$: \code{lower\_bound(3)}=end. Predecessor $=2$ (time $0$). Parent $=2$. Print $2$.
$s=\{1,2,3\}$, \code{when[3]=2}.\\
$i=3$, $x=5$: lb=end, pred $=3$ (time $2$). Parent $=3$. Print $3$. \code{when[5]=3}.\\
$i=4$, $x=4$: successor $=5$ (time $3$), predecessor $=3$ (time $2$). Later is $5$. Parent $=5$.
Print $5$. Output $2,2,3,5$ -- children $1,3,5,4$'s parents, matching the built tree. \checkmark
\end{tracebox}

\solstep{Complexity}
\begin{complexitybox}
Time $\bigO{n\log n}$: one \code{lower\_bound}, one \code{prev}, one \code{insert} per key, each
$\bigO{\log n}$ on a \code{std::set} of $\le n$ nodes (red-black -- \secref{c06-redblack} proves the
guarantee). Memory $\bigO{n}$ for the set and the time map. Crucially this is
$\bigO{n\log n}$ \emph{for every input order}, because we never let the tree's height enter the
cost -- the exact contrast with the $\bigO{n^2}$ simulation the section opened on. Verdict: \ac.
\end{complexitybox}

\solstep{Pitfalls and Variants}
\begin{warningbox}
Three ways to fall: (i) building the tree and descending (Trap~3) -- $\bigO{n^2}$ on sorted input.
(ii) forgetting distinctness -- here keys are distinct, so \code{lower\_bound} returns the first key
$\ge x$, which for a fresh $x$ is really the successor; with duplicates you must fix an
$\le$/$>$ convention. (iii) picking the \emph{earlier}-inserted neighbour -- inverted; it is the
\emph{later} one that is deeper and becomes the parent. \emph{Variants:} GNU \code{pb\_ds} or a hand
treap would also work but are overkill -- \code{set} already \emph{is} the balanced BST
(\secref{c06-redblack}). The same ``parent is the later neighbour'' idea reconstructs the whole tree,
not just the parent list.
\end{warningbox}

\begin{insightbox}[Key takeaways]
\begin{enumerate}[nosep, leftmargin=1.6em]
  \item \textbf{A rotation is the only legal move, and it is $\bigO{1}$ and inorder-preserving}
        (\thmref{c06-rotation}) -- every balanced tree is a \emph{policy} over rotations, never a new
        mechanic.
  \item \textbf{The shape is the enemy, not the data.} Same keys, same sorted line; a chain and a
        balanced tree differ only geometrically, and geometry decides $\bigO{h}$ vs $\bigO{n}$
        (\defref{c06-heightcost}).
  \item \textbf{You can often dodge the tree entirely.} When only sorted \emph{contents} and
        \emph{arrival times} matter, \code{std::set} gives $\bigO{\log n}$ without ever paying the
        degenerate height (\probref{c06-treecons}).
\end{enumerate}
\end{insightbox}

\begin{hardeningbox}
\begin{itemize}[nosep, leftmargin=1.4em]
  \item \textbf{Report depth instead of parent} -- with the same neighbours, a node's depth is one
        more than its (later-inserted) parent's; precompute times to answer both in one pass.
  \item \textbf{Reconstruct the whole tree} -- connect each key to its chosen parent and you have an
        $\bigO{n\log n}$ build of $C$, no descent.
  \item \textbf{Amortise later} -- splay (\secref{c06-splay}) rebuilds a BST from an order with the
        rotations themselves, giving $\bigO{n\log n}$ amortised \emph{while} physically holding a tree.
\end{itemize}
\end{hardeningbox}

\subsection{Practice Problems}

\begin{stuckbox}
\begin{itemize}[nosep, leftmargin=1.4em]
  \item Did you pick the \emph{later}-inserted neighbour, not the earlier, as the parent?
        (\probref{c06-treecons}.)
  \item Did you update heights \emph{bottom-first} after a rotation, so balance factors stay honest?
        (Trap~2, \coderef{c06-rotcode}.)
  \item For the rebuild practice: did you read the inorder (a sorted array) and then \emph{halve} it,
        rather than fiddling with rotations one by one?
\end{itemize}
\end{stuckbox}

\begin{practiceset}
  \practice{Balance a Binary Search Tree}{LeetCode 1382}
    {https://leetcode.com/problems/balance-a-binary-search-tree/}{LC Medium L3}
    {Given a possibly skewed BST, return \emph{any} height-balanced BST with the same values.
    The rotation-by-rotation route is fiddly; the clean route uses \thmref{c05-inorder}: dump the
    inorder into a sorted array, then build a balanced tree by recursively taking the middle element
    as the root (the same halving that makes a balanced tree from a sorted list). $\bigO{n}$ time and
    space, no rotations, and the sorted order is preserved by construction -- the extreme case of
    ``the shape is a choice'' (\thmref{c06-rotation}). Contrast with \probref{c06-treecons}, where the
    order was \emph{forced} and only the query was ours to answer.}
\end{practiceset}

\begin{recallbox}[10-second refresher -- 8.1 Why Balance? The Rotation]
Every BST operation is $\bigO{h}$ (\defref{c06-heightcost}), so keeping $h=\bigO{\log n}$ is the
whole game. The \term{rotation} re-links three pointers in $\bigO{1}$ and provably leaves the sorted
inorder unchanged (\thmref{c06-rotation}, \coderef{c06-rotcode}), so the shape is free to repair while
the data stays put. A same-side lean (LL, RR) needs one rotation; a zig-zag (LR, RL) needs two
(\defref{c06-rotcases}). And often the tree is not needed at all: CF 675D finds each parent as the
\emph{later-inserted} set neighbour in $\bigO{n\log n}$, sidestepping the degenerate height
(\probref{c06-treecons}).
\selftest{
  \item Draw the left rotation about $x$ and write the inorder of both pictures; why are they equal?
  \item The chain $1,2,3,4$ (all right children): how many rotations to balance, and which side each?
  \item Why is $Z$ (\figref{c06-zigzag}) not fixable by a \emph{single} rotation?
  \item In CF 675D, why is the \emph{earlier}-inserted neighbour provably not the parent?
  \item State the balance-height goal in one line: $h \le$ what, and why does it matter to search?
}
\end{recallbox}

% !TEX root = ../../tree_notes.tex
% =====================================================================
%  8.2  AVL Trees                                sec:c06-avl
%  Exports: def:c06-bal  thm:c06-avlbound  lem:c06-onerot  thm:c06-fixheight
%           fig:c06-avlfig  lst:c06-avlrot  lst:c06-avlins  ex:c06-avl-run
%           prob:c06-avl  lst:c06-avl-soln   (practice: GFG AVL deletion)
% =====================================================================
\section{AVL Trees}\label{sec:c06-avl}

\begin{previouslybox}
\begin{enumerate}[nosep, leftmargin=1.6em]
  \item State the one legal local move on a BST and what it provably preserves. (The rotation; the
        sorted inorder line is unchanged -- \thmref{c06-rotation}, \chref{c05}.)
  \item A node is a \emph{zig-zig} (single rotation) or \emph{zig-zag} (double rotation) violation --
        how do you tell which, from the heavy child's lean? (\defref{c06-rotcases},
        \figref{c06-zigzag})
  \item Every Chapter-7 routine costs $\bigO{h}$; if we force $h=\bigO{\log n}$, what do all of them
        become? ($\bigO{\log n}$ -- \defref{c06-heightcost}.)
\end{enumerate}
\end{previouslybox}

\begin{topicintro}
Section 8.1 gave you the rotation and the licence to reshape for free (\thmref{c06-rotation}); what it
did \emph{not} give you is a rule for \emph{when} to rotate so that the tree can never rot into the
chain $C$. The oldest such rule is due to Adelson-Velsky and Landis -- the \term{AVL tree} -- and it is
the most honest of the four strategies in this chapter: it \emph{demands} that every node be within
one level of balance, and it \emph{pays} for that demand with a rotation on the spot where any edit
breaks the promise. The reward is a hard worst-case bound: an AVL tree of $n$ keys has height at most
$1.44\,\lg(n+2)$ -- provably, for every input order (\thmref{c06-avlbound}). This section derives that
bound from a Fibonacci-style counting argument, shows the exact bookkeeping (a height field and a
balance factor carried up the single access path), and proves the elegant stopping rule: after one
insertion \emph{at most two rotations} are ever needed, and you stop at the first ancestor whose height
did not change (\thmref{c06-fixheight}). We end by implementing insertion cold on GeeksforGeeks'
AVL-insertion judge, which asks for the preorder of the tree you build.
\textbf{Prerequisites.} Sections 8.1 (rotations, \secref{c06-rotations}) and all of Chapter~7
(\chref{c05}); preorder is used only as an output format, defined in \chref{c03}.

\textbf{After this section you can}
\begin{itemize}[nosep]
  \item carry a per-node height and balance factor and update them on the way back up an insertion; \hfill\lv{3}
  \item classify a violation at the lowest unbalanced ancestor and apply the right one/two rotation; \hfill\lv{3}--\lv{4}
  \item prove the $\le 1.44\lg(n+2)$ height bound from the minimum-node recurrence; \hfill\lv{4}--\lv{5}
  \item implement AVL insertion (and reason about deletion) so the tree never degenerates. \hfill\lv{4}
\end{itemize}
\end{topicintro}

\begin{recognitionbox}[Pattern card -- must stay short, worst case]
\reckey{a keyed set/map that must answer \emph{every} operation in guaranteed $\bigO{\log n}$, with no
randomness and no amortisation -- you own the balance and refuse to let a bad input order hurt you.}
\recuse{
  \item ``insert/delete these keys, keep it fast, print a traversal'' -- AVL insertion, the GFG judge
        (\probref{c06-avl});
  \item database / filesystem indexes where worst-case latency matters -- the strict AVL rule beats a
        lazy one;
  \item you need rank/select on a \emph{keyed} set and can augment each node -- add a size field to the
        exact same machinery (\secref{c06-treaps} reuses it).
}
\recnotuse{
  \item you are happy with the library -- \code{std::set}/\code{map} already balance (red-black,
        \secref{c06-redblack}); reach for AVL only when you must customise the node;
  \item the keys arrive adversarially but you just want \emph{expected} $\bigO{\log n}$ with less case
        work -- a treap (\secref{c06-treaps}) trades the four cases for a coin flip;
  \item the tree models a \emph{sequence}, not a set of keys -- implicit treap (\secref{c06-treaps}),
        not AVL.
}
One-line test: \emph{do you need to \textbf{own} the height, exactly, on every single edit? Then pay
AVL's strict tax and rotate on the spot.}
\end{recognitionbox}

\subsection{Intuition}

The chain $C$ of \figref{c06-chain} is what happens when nobody is watching the height. To keep a BST
short, pick a property that (i) is \emph{local} -- checkable at one node from its two children -- so you
never scan the tree, and (ii) is \emph{strong enough} to force global logarithmic height. AVL's property
is blunt and perfect: at every node, the heights of the left and right subtrees differ by at most one.
Call that difference the \term{balance factor}. Enforce $|\text{BF}|\le 1$ everywhere and the tree
cannot grow a long spine: every time a descent goes two levels in one direction, a rotation kicks in and
bends it back.

Because the check is local, the \emph{only} place a rotation can be needed is along the single
root-to-leaf path you just walked to insert (or delete) a key -- the one path whose subtree heights
changed. You climb back up that path, and at the first node whose balance factor is now $\pm 2$, you fix
it with the one- or two-rotation recipe of \defref{c06-rotcases} and \emph{stop}: \thmref{c06-fixheight}
shows that the repaired subtree is exactly as tall as it was before your edit, so every higher ancestor
is untouched. This is why AVL is cheap -- a single descent, a single climb, and never more than two
rotations -- despite its fearsome reputation.

\begin{figure}[H]
\centering
\begin{tikzpicture}[diagram, scale=0.9, transform shape]
  % before: 1 -> 2 -> 3 chain, BF(1)=+2 violation (RR from node 1's perspective using BF=H(L)-H(R))
  \node[troot]  (b1) at (0,0)     {$1$};
  \node[tnode]  (b2) at (1.5,-1.2) {$2$};
  \node[tleaf]  (b3) at (3.0,-2.4) {$3$};
  \draw (b1) -- (b2); \draw (b2) -- (b3);
  \node[annot, font=\scriptsize,text=warningred] at (0.9,0.55) {$\bal=-2$};
  \node[annot, font=\scriptsize] at (-1.1,-1.2) {$\bal=-1$? no:\\ node $1$ heavy RIGHT};
  \node[annot, align=left, font=\scriptsize] at (1.5,-3.7) {insert $1,2,3$: node $1$ now at
    $\bal=-2$ (the RR case)};
  % arrow
  \draw[flow, double] (4.3,-1.4) -- (5.7,-1.4) node[midway, above, font=\scriptsize] {rotateLeft(1)};
  % after: 2(1,3)
  \node[troot]  (a2) at (7.2,0)     {$2$};
  \node[tleaf]  (a1) at (6.0,-1.2) {$1$};
  \node[tleaf]  (a3) at (8.4,-1.2) {$3$};
  \draw (a2) -- (a1); \draw (a2) -- (a3);
  \node[annot, font=\scriptsize,text=successgreen] at (7.2,0.6) {all $\bal\in\{-1,0,1\}$};
  \node[annot, align=left, font=\scriptsize] at (7.2,-2.6) {height $1$, inorder $1,2,3$\\ \emph{identical}
    -- \thmref{c06-rotation}};
\end{tikzpicture}
\caption{AVL's whole job in one picture. Inserting $1,2,3$ (a mini chain $C$) pushes node $1$'s balance
factor to $-2$: its two subtrees now differ in height by two, violating $|\bal|\le 1$. That is the RR
case, cured by a single \texttt{rotateLeft(1)} that lifts $2$ to the root: $2\to(1,3)$, every balance
factor legal, height back to $1$, inorder untouched. The same repair, run on the full $1,\dots,7$, yields
the balanced twin of \figref{c06-chain}.}\label{fig:c06-avlfig}
\end{figure}

\subsection{Definitions and Notation}

\begin{definition}{The balance factor and the AVL invariant}{c06-bal}
Store at each node $v$ the height $\height(v)$ of its subtree (with $\height(\nil)=-1$) and the
\term{balance factor} $\bal(v)=\height(v.\mathit{left})-\height(v.\mathit{right})$. A tree is
\term{AVL} if for \emph{every} node, $\bal(v)\in\{-1,0,+1\}$ -- the \term{AVL invariant}. The set of
nodes with $\bal(v)=\pm 2$ after an edit are exactly the ones needing repair.
\end{definition}

\noindent\textbf{In plain words.} Every node keeps a tiny integer that says ``how lopsided am I?'' AVL
forbids any node from being more than one level lopsided, and each node \emph{reports} its own lopsidedness
in $\bigO{1}$, so checking the whole tree is only checking the nodes you actually touched.

\begin{theorem}{An AVL tree has height $\le 1.44\,\lg(n+2)$}{c06-avlbound}
If $T$ is AVL with $n$ nodes and height $h$, then $h \le 1.44\,\lg(n+2)$ (all logs base $2$). Equivalently
$n \ge N(h)$ where $N(h)$ is the fewest nodes an AVL tree of height $h$ can have.
\end{theorem}
\begin{proofbox}[$N(h)$ is a shifted Fibonacci number]
Count the \emph{smallest} AVL tree of height $h$. Its two children must be AVL; to minimise nodes one
child has height $h-1$ and, because $|\bal|\le 1$, the other has height at least $h-2$. Making both as
small as possible gives the recurrence $N(h)=N(h-1)+N(h-2)+1$ with $N(0)=1$, $N(1)=2$. Subtracting $1$
throughout, $M(h)=N(h)+1$ satisfies $M(h)=M(h-1)+M(h-2)$, the Fibonacci recurrence, so
$N(h)=F_{h+3}-1$, where $F$ is the Fibonacci sequence. Fibonacci grows like $\varphi^k$ with
$\varphi=\tfrac{1+\sqrt5}{2}\approx 1.618$, so any AVL tree with $n$ nodes obeys
$n \ge F_{h+3}-1 > \varphi^{\,h+1}-1$, i.e.\ $\varphi^{\,h+1} < n+2$. Taking $\lg$,
$(h+1)\lg\varphi < \lg(n+2)$, hence $h < \lg(n+2)/\lg\varphi - 1 \le 1.4404\,\lg(n+2)$. The constant
$\lg\varphi\approx 0.694$ is exactly why $1.44=1/0.694$ appears. \end{proofbox}
\begin{insightbox}
The worst AVL shape -- the one that hides the most height behind the fewest nodes -- is a
Fibonacci-spiral of minimal subtrees, not a chain. AVL is strong enough to make the chain impossible but
does \emph{not} insist on a perfect binary tree; the gap between $1.44\lg n$ and the ideal $\lg n$ is the
price of the local, cheap rule. Red-black (\secref{c06-redblack}) pays a slightly bigger constant
($2\lg(n+1)$) for fewer rotations; treaps (\secref{c06-treaps}) hit $\bigO{\log n}$ only \emph{in
expectation} but need no case analysis at all.
\end{insightbox}

\begin{lemma}{A rotation makes the restructured subtree legal in one/two turns}{c06-onerot}
At a lowest unbalanced node $w$ (so $|\bal(w)|=2$) whose children are themselves balanced, the repair of
\defref{c06-rotcases} -- a single rotation for a zig-zig (LL, RR) and two for a zig-zag (LR, RL) --
produces a subtree rooted at the middle node of $\{w,\text{child},\text{grandchild}\}$ in which all
three top balance factors lie in $\{-1,0,+1\}$, preserving the inorder of $w$'s subtree.
\end{lemma}
\begin{proofbox}
By \thmref{c06-rotation} the rotation preserves inorder, so the three keys keep their relative order; the
new root is the median. In a zig-zig the three nodes were collinear on one side, so after the single
rotation they form a height-$1$ balanced triple (the two off-path subtrees were, by the choice of $w$ as
\emph{lowest} violator, already within one level, and the promotion evens them out). In a zig-zag the
child sits \emph{between} $w$ and the grandchild; the first rotation straightens that child into a
zig-zig, and the second is the zig-zig case just handled. Either way the three involved nodes end legal
and every off-path subtree $A,B,C$ keeps its exact contents. \end{proofbox}

\begin{theorem}{After one insertion, $\le 2$ rotations suffice and you stop early}{c06-fixheight}
Insert a key into an AVL tree. Walking from the insertion point up to the root, let $w$ be the first node
with $\bal(w)=\pm 2$. Then: (i) every node strictly below $w$ on the path is already balanced; (ii) the
\defref{c06-rotcases} repair at $w$ makes $w$'s restructured subtree exactly as tall as it was
\emph{before} the insertion; and (iii) consequently all of $w$'s ancestors keep their old balance factors
-- a single repair suffices and the walk may stop. Deletion can require more than one repair (a
shortened subtree propagates upward), but still $\bigO{\log n}$ rotations total.
\end{theorem}
\begin{proofbox}
Only the path to the inserted leaf changed height, so the first $\pm2$ node encountered going up is the
\emph{lowest} violator -- its children are balanced, exactly \lemref{c06-onerot}'s hypothesis. That
lemma's repair turns the offending collinear/zig-zag triple into a balanced triple whose height equals
the pre-insertion height of $w$'s subtree: before the insertion, $w$'s heavy side had height one more
than its light side (differing by at most $1$); the inserted key lengthened the heavy side by one (making
$\bal=\pm2$); re-rooting the triple at the median redistributes so the subtree is exactly the old heavy
height again. Since the height $w$'s parent sees is unchanged, the parent's balance factor is unchanged,
and by induction none of the ancestors becomes unbalanced -- so we stop. For deletion the subtree can
\emph{shrink} by one at the first violator, and a shrink does propagate, which is why deletion may need
several repairs up the path -- but each is $\bigO{1}$ and the path has length $\bigO{h}=\bigO{\log n}$ by
\thmref{c06-avlbound}. \end{proofbox}

\subsection{The algorithm: rotate, then re-ascend only until heights stop changing}

The whole structure fits in three small helpers: a height accessor, an update that recomputes one node,
and a \texttt{balance} routine that applies \lemref{c06-onerot}. Insertion is Chapter~7's recursive
insert (\secref{c05-insert-delete}) with one extra line -- call \texttt{balance} on the way back up. That
line is the entire difference between a rotting BST and a guaranteed $\bigO{\log n}$ one.

\begin{lstlisting}[caption={Balance factors and the four-case repair -- the reuse of Section 8.1's rotations with height upkeep}, label={lst:c06-avlrot}]
int H(Node* n)        { return n ? n->ht : -1; }            // height, null = -1
void upd(Node* n)     { n->ht = 1 + max(H(n->left), H(n->right)); }
int  BF(Node* n)      { return H(n->left) - H(n->right); }  // balance factor

Node* rotateLeft(Node* x){ Node* y = x->right;              // x->right = B, y->left = x
    x->right = y->left; y->left = x; upd(x); upd(y); return y; }
Node* rotateRight(Node* y){ Node* x = y->left;              // mirror image
    y->left = x->right; x->right = y; upd(y); upd(x); return x; }

// Assumes n's children are already AVL; returns the (possibly rotated) subtree root.
Node* balance(Node* n){
    upd(n);
    if (BF(n) == 2) {                                   // heavy LEFT
        if (BF(n->left) < 0) n->left = rotateLeft(n->left);   // LR -> LL
        return rotateRight(n);                                // then fix LL
    }
    if (BF(n) == -2) {                                  // heavy RIGHT
        if (BF(n->right) > 0) n->right = rotateRight(n->right); // RL -> RR
        return rotateLeft(n);                                   // then fix RR
    }
    return n;                                           // already balanced
}
\end{lstlisting}

\begin{lstlisting}[caption={AVL insertion -- Chapter 7's insert plus one balance call on the return path}, label={lst:c06-avlins}]
Node* insert(Node* n, int key){
    if (n == nullptr) return new Node(key);            // plant the leaf (as in Ch 7)
    if (key < n->key)      n->left  = insert(n->left,  key);
    else if (key > n->key) n->right = insert(n->right, key);
    else                   return n;                   // distinct keys: no re-insert
    return balance(n);                                  // fix THIS node on the way back up
}
\end{lstlisting}

\begin{complexitybox}
Time $\bigO{\log n}$: the recursion is a single root-to-leaf descent of length $\bigO{h}$, and
\thmref{c06-avlbound} is what makes $h=\bigO{\log n}$ instead of $n$; the work per level is a constant
$\bigO{1}$ \texttt{balance} (at most two rotations by \thmref{c06-fixheight}). Memory $\bigO{\log n}$ for
the recursion stack. Verdict: every one-descent operation is \emph{guaranteed} $\bigO{\log n}$, worst
case -- the promise the chain $C$ kept breaking.
\end{complexitybox}

\subsection{Worked Example}

\begin{example}{AVL insertions on the running examples $C$ and $Z$}{c06-avl-run}
\emph{Chain $C$ ($1,2,\dots,7$), eagerly balanced.} Insert $1,2$ -- fine. Insert $3$: the path
$1\!\to\!2\!\to\!3$ leaves node $1$ at $\bal=-2$ (an RR), so \texttt{balance(1)} does one
\texttt{rotateLeft(1)} giving $2\to(1,3)$ (\figref{c06-avlfig}). Insert $4$: tree $2\to(1,3\to(\_,4))$,
legal. Insert $5$: node $3$ hits $\bal=-2$, \texttt{rotateLeft(3)} $\Rightarrow$ $2\to(1,4\to(3,5))$.
Insert $6$: node $2$ hits $\bal=-2$ (RR), \texttt{rotateLeft(2)} $\Rightarrow$
$4\to(2\to(1,3),\,5\to(\_,6))$. Insert $7$: node $5$ hits $\bal=-2$, \texttt{rotateLeft(5)}
$\Rightarrow$ the \emph{balanced twin} $4\to(2\to(1,3),\,6\to(5,7))$ of \figref{c06-chain}, height $2$,
preorder $4,2,1,3,6,5,7$. Four single rotations, each at the lowest violator; the tree never grew past
height $2$.\\
\emph{Zig-zag order $30,10,40,20,25$.} Here eager balancing reveals a subtlety: by the time $25$ is
added, the path is $30\!\to\!10\!\to\!20\!\to\!25$, and the \emph{lowest} node that goes illegal is $10$
($\bal(10)=-2$, a right-right triple $10\!\to\!20\!\to\!25$), \emph{not} $30$. \texttt{balance} does one
\texttt{rotateLeft(10)} $\Rightarrow$ $20\to(10,25)$, so the whole tree becomes
$30\to(20\to(10,25),\,40)$ -- height $2$, and node $30$'s balance factor never even reached $\pm2$.
Section 8.1 froze this sequence as the plain-BST shape violating at $30$ to \emph{illustrate} the LR case
(\figref{c06-zigzag}); \thmref{c06-fixheight} is exactly why a running AVL repairs the lower node first
and stops. A genuine LR during insertion needs the zig-zag to be the \emph{lowest} violation, e.g.\
$30,10,20$: node $30$ hits $\bal=+2$ with child $10$ leaning right, so we \texttt{rotateLeft(10)} then
\texttt{rotateRight(30)} $\Rightarrow$ $20\to(10,30)$. Inorder $10,20,25,30,40$ is preserved throughout
all these edits, honouring \thmref{c06-rotation}.
\end{example}

\begin{tracebox}[Eager $1,2,3$ and a real LR at $30,10,20$, pointer by pointer]
\textbf{$1,2,3$:} after inserting $3$, $\text{BF}(1)=\height(\_)-\height(2\!\to\!3)=-1-1=-2$ (RR).
\texttt{rotateLeft(1)}: $y:=2$; $1\to$right $:=2$'s left (null); $2\to$left $:=1$; \texttt{upd(1)},
\texttt{upd(2)}. Root $=2\to(1,3)$, all $\bal\in\{-1,0,1\}$. \checkmark\\
\textbf{$30,10,20$:} insert $30$, then $10$ (left), then $20$ ($<\!30$ go left, $>\!10$ go right).
Going up: $\text{BF}(10)=-1$ (ok), $\text{BF}(30)=\height(10)-\height(\_)=1-(-1)=+2$ (heavy LEFT, child
$10$ leans right: LR). \texttt{n->left=rotateLeft(10)}: $20$ up, $10$ left, giving $30\to(20\to(10,\_),\_)$;
then \texttt{rotateRight(30)}: $20$ up, $30$ right, $20$'s right (null) to $30$'s left $\Rightarrow$
$20\to(10,30)$. Two rotations, inorder $10,20,30$ intact. \checkmark
\end{tracebox}

\subsection{Variations and Twists}

\begin{itemize}[nosep, leftmargin=1.4em]
  \item \textbf{Deletion is the hard half}: after splicing a node (\thmref{c05-delete}), a subtree can
        \emph{shrink}, and \thmref{c06-fixheight} only guarantees early stopping for \emph{heights that
        do not change}. You must re-balance every ancestor on the path -- still $\bigO{\log n}$ total, but
        the case split is nastier (a null balance factor needs care). \hfill\lv{5}
  \item \textbf{Size augmentation} turns AVL into an order-statistic tree (rank, kth) -- the same
        \texttt{upd} just also sums children's sizes; the rotations are unchanged. \hfill\lv{4}
  \item \textbf{Relaxed AVL} -- AVL-like rules ($|\bal|\le 2$, or Fibonacci-weight) cut the rotation count
        at the price of a worse constant, the ancestor of red-black's philosophy. \hfill\lv{4}--\lv{5}
  \item \textbf{Iterative rebalance}: remember the path in a stack, then retrace it calling
        \texttt{balance} -- the same asymptotics without deep recursion when $n$ is huge. \hfill\lv{4}
\end{itemize}

\subsection{Common Pitfalls}\label{sec:c06-avl-pitfalls}

\begin{warningbox}[Trap 1 -- forgetting to recompute heights on the way down, not up]
A rotation changes \emph{two} nodes' heights. In \texttt{rotateLeft} you must \texttt{upd(x)} \emph{then}
\texttt{upd(y)}, because $x$ moved below $y$; recomputing $y$ first reads $x$'s stale height and silently
corrupts every balance factor above. The fix is exactly the two lines' order in \coderef{c06-avlrot}.
\end{warningbox}
\begin{lstlisting}[caption={\textcolor{warningred}{WRONG}: recompute the promoted node before the demoted one}, basicstyle=\ttfamily\small, language=C++]
x->right = y->left; y->left = x;
upd(y); upd(x);   // y read x's OLD height -- corrupts BF all the way up
\end{lstlisting}
\begin{lstlisting}[caption={\textcolor{successgreen}{RIGHT}: fix the lower node first, then its new parent}, basicstyle=\ttfamily\small, language=C++]
x->right = y->left; y->left = x;
upd(x); upd(y);   // x final, then y sees x's new height
\end{lstlisting}

\begin{warningbox}[Trap 2 -- rotating at the wrong ancestor, or not stopping]
\thmref{c06-fixheight}: fix the \emph{lowest} node with $|\bal|=2$ and (for insertion) then stop -- its
subtree returns to the old height, so ancestors are already fine. Rotating at a higher node first is
wrong: the lowest violator's children may not be balanced, breaking \lemref{c06-onerot}'s hypothesis.
Conversely for \emph{deletion} you must \emph{not} stop -- a shrinking subtree propagates (\secref{c06-avl-pitfalls}, Var.~1).
\end{warningbox}

\begin{warningbox}[Trap 3 -- the LR/RL child rotation forgets to reattach]
In \texttt{balance}, the first rotation of a zig-zag must be written back: \texttt{n->left =
rotateLeft(n->left)}. Drop the assignment and the promoted child is orphaned -- the tree loses a whole
subtree and the judge prints a short traversal (a \textbf{WA} on the GFG preorder check).
\end{warningbox}

\subsection{Problem Lab}

The anchor is GeeksforGeeks' AVL insertion judge: insert keys one by one into an AVL tree and print the
preorder of the result. It is the purest possible AVL drill -- no queries, just ``build it balanced and
show me the shape'' -- so a single misplaced rotation or stale height is instantly visible in the
traversal.

% ---------------------------------------------------------------------------
\solvedproblem{c06-avl}{AVL Tree Insertion}{GeeksforGeeks}
  {https://www.geeksforgeeks.org/problems/avl-tree-insertion/1}{GFG Hard · L4}
  {per-node height and balance factor, four-case rebalancing, recursion unwinds the single access path}

\begin{problemstatement}[AVL Tree Insertion]
You are given $t$ test cases. Each gives an integer $n$ and a sequence of $n$ \textbf{distinct} keys
inserted \emph{one by one, in order,} into an initially empty AVL tree. After each insertion the tree is
rebalanced so the AVL invariant $|\bal(v)|\le 1$ holds everywhere. For each test case print the
\textbf{preorder traversal} (root, left, right) of the resulting AVL tree, keys separated by spaces.
Typical limits: $1\le n\le 10^{3}$ per case, keys fit in a 32-bit integer.
\end{problemstatement}

\begin{attemptgate}[Attempt it yourself first -- \lv{4}, give yourself 30 minutes]
This is \coderef{c06-avlins} and \coderef{c06-avlrot} assembled into a program with I/O and a preorder.
Do not hand-code four separate cases -- write one \texttt{balance} and call it on the way up.
\peek{When you recursively insert into the left or right child, the child's subtree root may change
(rotation). Where must you re-balance, and how do you carry the returned height upward?}
\nudge{Write \texttt{insert} so it \emph{returns} the (new) subtree root, and \texttt{balance} the current
node right after the recursive call. \texttt{balance} recomputes the height, checks $|\bal|$, and applies
\defref{c06-rotcases}: heavy side and that child's lean decide single vs double.}
\fullhint{Use \texttt{H}, \texttt{upd}, \texttt{BF} and the two rotations exactly as in \coderef{c06-avlrot}.
\texttt{balance(n)}: \texttt{upd(n)}; if \texttt{BF(n)==2} then (if \texttt{BF(n->left)<0}
\texttt{n->left=rotateLeft(n->left)}) return \texttt{rotateRight(n)}; symmetric for $-2$; else return
\texttt{n}. Then \texttt{insert} is Ch.7's descent with \texttt{return balance(n)} on the way up, and the
driver reads keys, calls \texttt{root=insert(root,x)}, and prints \texttt{preorder(root)}.}
\end{attemptgate}

\solstep{Understanding the Problem}
We must reproduce the AVL insertion algorithm faithfully and then emit a specific traversal, so the
judge can compare the exact tree we built. Two things are load-bearing: (i) the tree must be balanced
after \emph{every} key, not just at the end -- though because AVL rotations are deterministic given the
arrival order, building balanced-on-each-insert and reading the final preorder is the same as inserting
and rebalancing throughout; and (ii) our rotations must update heights in the correct order so balance
factors above stay honest. There is no query to answer -- the \emph{shape is the answer} -- which makes
this an ideal correctness drill.

\solstep{Key Observations}
\begin{observation}
Only the recursion path from the root to the freshly-planted leaf can have changed height, so
\thmref{c06-fixheight} says at most one node on it is illegal after an insertion -- and \texttt{balance}
on the unwinding recursion finds and fixes it with $\le 2$ rotations.
\end{observation}
\begin{observation}
The heavy child's own lean ($\texttt{BF(child)}$'s sign) separates the single-rotation zig-zig cases
(LL, RR) from the double-rotation zig-zag cases (LR, RL): same-sign lean is zig-zig, opposite-sign needs
the child straightened first (\defref{c06-rotcases}).
\end{observation}
\begin{observation}
Preorder is produced by the standard root-left-right walk (\chref{c03}); the only AVL-specific care is
that the \emph{tree we walk} is the balanced one, so a wrong rotation shows up as a wrong traversal.
\end{observation}

\solstep{Think It Through With Me}
\begin{thinkbox}
\deadend{Insert all keys into a plain BST (\secref{c05-insert-delete}) and rebalance the whole tree once
at the end by rebuilding from the sorted inorder.}
{That produces \emph{a} balanced BST but \emph{not the AVL tree} the judge expects -- AVL's rotations are
order-dependent, so the final shape differs from the median-halving rebuild (\figref{c06-chain}). Wrong
preorder, \textbf{WA}. We must run AVL's local repair after each insertion.}
\deadend{Code four separate rebalancing routines (LL, RR, LR, RL) chosen by inspecting two levels.}
{Reproducible and bug-prone: the child rotation's write-back (Trap~3) is easy to forget, and you duplicate
\texttt{upd} logic. The single \texttt{balance} that composes the two rotations is strictly less code.}
\aha{Write the two rotations once (\coderef{c06-avlrot}), a \texttt{balance} that picks one or two based on
the child's lean, and a Ch.7-style \texttt{insert} that calls \texttt{balance} on the way up. The four cases
emerge from \texttt{BF(n)} and \texttt{BF(n->left/right)} without a case table.}
\end{thinkbox}

\solstep{Solution}
Maintain at each node its key, height and two children. \texttt{H}/\texttt{upd}/\texttt{BF} as in
\coderef{c06-avlrot}. \texttt{insert(n,key)}: recurse left or right by the BST rule, \emph{reassigning}
the child to the returned (possibly rotated) root, then return \texttt{balance(n)}. \texttt{balance(n)}:
\texttt{upd(n)}; if $\bal(n)=+2$ (heavy left) straighten the child if $\bal(\text{left})<0$ then
\texttt{rotateRight(n)}; if $\bal(n)=-2$ symmetric; else return $n$ unchanged. The driver builds the tree
per test case and prints \texttt{preorder(root)}.

\solstep{Proof of Correctness}
\begin{proofbox}[The built tree is the AVL tree; the traversal is its preorder]
By induction on the number of inserted keys. Base: after the first key the single node has $\bal=0$, AVL.
Inductive step: assume the tree is AVL before inserting $x$. \texttt{insert} performs Chapter~7's BST
descent, so the inorder (sorted order) of the tree is exactly ``old keys plus $x$''; the recursive
reassignment keeps every subtree correctly linked. Only the path to the new leaf changed height, so by
\thmref{c06-fixheight} at most one node $w$ on it is illegal, and it is the \emph{lowest}. The unwinding
calls \texttt{balance} at every path node; nodes below $w$ return unchanged ($|\bal|\le1$), and at $w$ the
heavy-side/child-lean tests select precisely the \defref{c06-rotcases} repair, whose validity is
\lemref{c06-onerot}. That repair leaves $w$'s subtree with the same inorder but a legal root, and (for
insertion) its height equal to the pre-insert value, so every ancestor's $\bal$ is unchanged and their
\texttt{balance} calls are no-ops. Hence the returned tree is AVL and stores exactly the keys inserted so
far, completing the induction. The final \texttt{preorder} walk is root-left-right (\chref{c03}), the
definition of preorder, so it matches the judge's expected traversal of that AVL tree. \end{proofbox}

\solstep{Implementation}
\begin{lstlisting}[caption={AVL Tree Insertion -- height-carrying balance plus a preorder walk (GeeksforGeeks)}, label={lst:c06-avl-soln}]
#include <bits/stdc++.h>
using namespace std;
struct Node { int key; int ht; Node *left, *right;
    Node(int k) : key(k), ht(0), left(nullptr), right(nullptr) {} };

int  H (Node* n) { return n ? n->ht : -1; }
void upd(Node* n) { n->ht = 1 + max(H(n->left), H(n->right)); }
int  BF(Node* n) { return H(n->left) - H(n->right); }
Node* rotateLeft(Node* x){ Node* y=x->right; x->right=y->left; y->left=x; upd(x); upd(y); return y; }
Node* rotateRight(Node* y){ Node* x=y->left;  y->left=x->right; x->right=y; upd(y); upd(x); return x; }

Node* balance(Node* n){
    upd(n);
    if (BF(n) == 2) { if (BF(n->left) < 0) n->left = rotateLeft(n->left); return rotateRight(n); }
    if (BF(n) == -2){ if (BF(n->right) > 0) n->right = rotateRight(n->right); return rotateLeft(n); }
    return n;
}
Node* insert(Node* n, int key){
    if (!n) return new Node(key);
    if (key < n->key)      n->left  = insert(n->left,  key);
    else if (key > n->key) n->right = insert(n->right, key);
    else                   return n;
    return balance(n);
}
void preorder(Node* n){ if(!n) return; cout << n->key << " "; preorder(n->left); preorder(n->right); }

int main(){
    ios::sync_with_stdio(false); cin.tie(nullptr);
    int t; cin >> t;
    while (t--) {
        int n; cin >> n; Node* root = nullptr;
        for (int i = 0; i < n; ++i) { int x; cin >> x; root = insert(root, x); }
        preorder(root); cout << "\n";
    }
    return 0;
}
\end{lstlisting}

\solstep{Trace}
\begin{tracebox}[Official-style sample: keys $21,26,30$]
$21$: root. $26$: right of $21$; $\text{BF}(21)=-1$, $\text{BF}(26)=0$ -- balanced. $30$: path
$21\!\to\!26\!\to\!30$. Unwind: $\text{balance}(30)$ no-op; $\text{balance}(26)$: $\text{BF}(26)=-1$
no-op; $\text{balance}(21)$: $\text{BF}(21)=\height(\_)-\height(26\!\to\!30)=-1-1=-2$ (heavy right, child
$26$ has $\text{BF}=-1$, not $>0$ $\Rightarrow$ RR single). \texttt{rotateLeft(21)}: $26$ up, $21$ left,
$30$ right. Final tree $26\to(21,30)$, preorder $\;26\;\;21\;\;30$. \checkmark
\end{tracebox}

\solstep{Complexity}
\begin{complexitybox}
Time $\bigO{n\log n}$: $n$ insertions each a $\bigO{\log n}$ descent-plus-repair (\thmref{c06-avlbound}
bounds the height; \thmref{c06-fixheight} bounds the rotations at $\le2$ per key), plus $\bigO{n}$ for the
single preorder walk. Memory $\bigO{n}$ for nodes and $\bigO{\log n}$ recursion stack. This is exactly
what makes the same driver, run on the sorted worst case $1,\dots,n$, finish in $\bigO{n\log n}$ where the
plain-BST construction of \probref{c06-treecons} would have ground to $\bigO{n^2}$. Verdict: \ac.
\end{complexitybox}

\solstep{Pitfalls and Variants}
\begin{warningbox}
(1) \texttt{upd} order in a rotation (Trap~1). (2) Forgetting \texttt{else return n;} on a duplicate key --
harmless here since keys are distinct, but a re-insert that recurses forever is a real bug elsewhere. (3)
Printing trailing/leading spaces: the judge tokenises on whitespace, but keep the walk clean. \emph{Variant:}
swap \texttt{preorder} for \texttt{inorder} and you should always read the sorted keys -- a fast
sanity check that no rotation lost a subtree. \emph{Extension to deletion} needs the re-balance to run up
the \emph{whole} path (Var.~1), not stop early.
\end{warningbox}

\begin{insightbox}[Key takeaways]
\begin{enumerate}[nosep, leftmargin=1.6em]
  \item \textbf{A local rule buys a global bound.} Keeping $|\bal|\le1$ at every node forces height
        $\le 1.44\,\lg(n+2)$ (\thmref{c06-avlbound}) -- the chain $C$ becomes impossible.
  \item \textbf{Balance is cheap because the fix is local.} Only the access path changed; one
        \texttt{balance} call per level with $\le2$ rotations total (\thmref{c06-fixheight}) and you stop.
  \item \textbf{Case work collapses into two tests.} $\texttt{BF}(n)$ picks the heavy side, and
        $\texttt{BF}(\text{child})$ picks single-vs-double -- no four-branch table (\coderef{c06-avlrot}).
\end{enumerate}
\end{insightbox}

\begin{hardeningbox}
\begin{itemize}[nosep, leftmargin=1.4em]
  \item \textbf{AVL deletion} -- the harder sibling: rebalance \emph{every} ancestor (heights can shrink),
        practised on the GFG deletion judge below.
  \item \textbf{Order-statistics AVL} -- store subtree \emph{size} alongside height; rotations are
        identical, and you gain rank/kth queries -- the augmented-tree preview of \secref{c06-treaps}.
  \item \textbf{Compare with the rebuild} -- LC 1382 balances a static BST by median-halving the inorder
        ($\bigO{n}$, no rotations); AVL balances a \emph{dynamic} set one key at a time. Know when each wins.
\end{itemize}
\end{hardeningbox}

\subsection{Practice Problems}

\begin{stuckbox}
\begin{itemize}[nosep, leftmargin=1.4em]
  \item Did you \texttt{upd} the demoted node before the promoted one in every rotation? (Trap~1.)
  \item Did you reassign the child on the zig-zag's first rotation, \texttt{n->left = rotateLeft(n->left)}?
        (Trap~3.)
  \item On deletion, did you resist stopping early -- a shrunk subtree must propagate upward? (Var.~1.)
\end{itemize}
\end{stuckbox}

\begin{practiceset}
  \practice{AVL Tree Deletion}{GeeksforGeeks}
    {https://www.geeksforgeeks.org/problems/avl-tree-deletion/1}{GFG Hard · L4}
    {Delete a key from an AVL tree and print the inorder of the result. Reuse \coderef{c06-avlrot}
    verbatim; the only new wrinkle is \thmref{c06-fixheight}'s caveat: after Ch.7's successor-splice delete
    (\thmref{c05-delete}) the subtree may be one shorter, so \texttt{balance} the \emph{entire} path back
    to the root, not just the first violator. The four rotation cases at each level are unchanged. This is
    the single most bug-prone routine in the chapter -- test the two-child delete and the null-balance
    child carefully.}
\end{practiceset}

\begin{recallbox}[10-second refresher -- 8.2 AVL Trees]
AVL enforces $|\bal(v)|\le1$ at every node, storing a height per node (\defref{c06-bal}); that local rule
forces height $\le 1.44\lg(n+2)$ via the Fibonacci minimum-node count (\thmref{c06-avlbound}). Insert is
Ch.7's descent plus one \texttt{balance} call on the way up (\coderef{c06-avlins}); a single $\texttt{BF}(n)$
test plus the child-lean test picks the one- or two-rotation fix (\defref{c06-rotcases},
\lemref{c06-onerot}). After an insertion the repair lands at the \emph{lowest} violator and stops, because
the restructured subtree regains its old height -- $\le2$ rotations total (\thmref{c06-fixheight}). Every
operation is thus guaranteed $\bigO{\log n}$, worst case -- the enemy of \probref{c06-treecons}'s chain is
finally dead.
\selftest{
  \item Why does $|\bal|\le1$ forbid a linear-height tree, and what shape \emph{is} the tallest AVL?
  \item Give the balance factors before/after \texttt{rotateLeft(1)} on the chain $1\to2\to3$.
  \item Why may insertion stop at the first repaired ancestor but deletion may not?
  \item In \texttt{balance}, what does $\texttt{BF(n->left)}<0$ tell you to do before \texttt{rotateRight(n)}?
  \item State the AVL height bound with its constant and derive it in one line from $N(h)$'s recurrence.
}
\end{recallbox}

% !TEX root = ../../tree_notes.tex
% =====================================================================
%  8.3  Treaps, Split/Merge, Implicit Keys       sec:c06-treaps
%  Exports: def:c06-treap  thm:c06-exptheight  def:c06-splitmerge
%           lem:c06-mergevalid  prop:c06-implicit  fig:c06-splitfig
%           lst:c06-treap  lst:c06-split  lst:c06-imp  ex:c06-treap-run
%           prob:c06-rpmq  lst:c06-rpmq-soln   (practice: CSES 1190)
% =====================================================================
\section{Treaps, Split/Merge, Implicit Keys}\label{sec:c06-treaps}

\begin{previouslybox}
\begin{enumerate}[nosep, leftmargin=1.6em]
  \item What single local move reshapes a BST without touching its inorder line? (The rotation,
        \thmref{c06-rotation}.)
  \item AVL keeps height $\le 1.44\,\lg(n+2)$ by enforcing what local rule, and how many rotations does
        one insertion need? ($|\bal|\le1$ everywhere; $\le2$, \thmref{c06-fixheight},
        \thmref{c06-avlbound}.)
  \item Why does a \emph{random} insertion order into a plain BST give height $\bigO{\log n}$ in
        expectation? (the random-BST fact we now reuse via a coin flip -- \secref{c05-insert-delete}.)
\end{enumerate}
\end{previouslybox}

\begin{topicintro}
AVL paid for its worst-case bound in \emph{case work}: four rebalancing shapes, a balance factor, a
height field maintained by hand. The \term{treap} throws away the case work and keeps the guarantee by
\emph{randomising}. Give every key an independent random priority; demand the tree be a BST on the keys
\emph{and} a heap on the priorities. With distinct priorities there is exactly \emph{one} such tree --
the Cartesian tree of the key/priority pairs -- and because the priorities are random, that tree is a
\emph{random} BST, whose height is $\bigO{\log n}$ in expectation without you ever balancing anything
(\thmref{c06-exptheight}). The real payoff is structural: a treap is maintained by just two universal
routines, \texttt{split} and \texttt{merge} (\defref{c06-splitmerge}), from which insert, erase, rank,
and range queries all fall out as one-liners -- no rotations, no four cases. Then the trick that makes
this section worth a chapter: drop the explicit keys and use each node's \emph{subtree size} as its
position (\propref{c06-implicit}). A treap becomes a dynamic sequence supporting insert/delete/reverse
\emph{by index} in $\bigO{\log n}$ expected -- the machinery behind the implicit-treap solutions to
sequence problems. We land on LeetCode 307, maintaining a mutable array's range sums with exactly this
size-augmented split/merge.
\textbf{Prerequisites.} \secref{c06-rotations}, \secref{c06-avl} (you should feel AVL's case work before
we abolish it); Chapter~7 for the BST descent. Cartesian trees and the monotonic stack that builds them
are \emph{not} developed here -- that is Chapter~9 (\chref{c26}); we only need that a random priority
uniquely fixes a tree.

\textbf{After this section you can}
\begin{itemize}[nosep]
  \item build a treap from random priorities and explain why its height is $\bigO{\log n}$ expected; \hfill\lv{4}
  \item implement \texttt{split} and \texttt{merge} and derive insert/erase/rank from them; \hfill\lv{4}
  \item use an \emph{implicit-key} treap as a dynamic sequence (insert/delete/reverse by index); \hfill\lv{4}--\lv{5}
  \item augment a treap with an aggregate (a sum) to answer range queries under point updates. \hfill\lv{4}
\end{itemize}
\end{topicintro}

\begin{recognitionbox}[Pattern card -- you want a balanced tree with no rebalancing cases]
\reckey{you need a \emph{customisable} balanced search tree (own the node, carry aggregates, cut/paste
subtrees) but AVL's four-case bookkeeping is too much code, and the library \code{set} cannot be
extended.}
\recuse{
  \item a keyed set/map with aggregate queries you must implement yourself -- treap with a size/sum field
        (\defref{c06-treap});
  \item operations on a \emph{sequence} by position -- insert/delete/reverse/concat a range -- implicit
        treap (\propref{c06-implicit}, the LC 307 route \probref{c06-rpmq});
  \item problems needing subtree \texttt{split}/\texttt{merge} as the primitive (rope editors, persistent
        sequences) -- see \chref{c23} for the link-cut upgrade.
}
\recnotuse{
  \item a plain ordered set/map suffices -- just use \code{std::set}/\code{map} (red-black,
        \secref{c06-redblack}); no need to implement anything;
  \item the array is \emph{fixed-size} with only point updates and prefix/range sums -- a Fenwick tree
        (\chref{c14}) is far less code and fully deterministic;
  \item you need a \emph{deterministic} worst-case height and randomness is disallowed (rare) -- AVL
        (\secref{c06-avl}).
}
One-line test: \emph{do you need to reshape arbitrary subtrees (split/merge) or index by position, and
would rather flip a coin than write rotation cases? Treap.}
\end{recognitionbox}

\subsection{Intuition}

A BST is determined by the \emph{order} keys arrive in; a treap is determined by an order you
\emph{invent}. Attach to each key $k$ a random priority $p(k)$, and insist that the tree satisfy two
rules at once: BST on the keys (so the inorder is sorted) and heap on the priorities (so the root has the
best priority). These pin down a unique tree: the root must be the max-priority key, and it splits the
remaining keys into a left and a right block that recurse identically. Now the miracle -- because the
priorities are random, ``root is the max-priority key'' means the root is a \emph{uniformly random} key,
so the treap is distributed exactly like a BST built by inserting keys in random order: balanced to
depth $\bigO{\log n}$ in expectation (\thmref{c06-exptheight}). You get AVL's height without AVL's
\emph{work}: nothing is ever rotated deliberately; you simply re-roll nothing and maintain the heap
invariant, and randomness does the balancing.

The second idea is the load-bearing one. To edit a treap you never rotate -- you \texttt{split} it at a
key (or a position) into two treaps and \texttt{merge} two treaps whose keys are already all-comparable
back into one. Insert becomes ``split at $x$, merge the left, a fresh node, and the right''; erase is
``split around $x$, drop the middle, merge the rest.'' And if you index nodes not by a stored key but by
their inorder \emph{position} (which $\mathit{size}(\mathit{left})$ gives you for free), the same two
routines turn the treap into a dynamic array: cut out a range, paste it elsewhere, or -- with a lazy flag
-- reverse it in place. That is \propref{c06-implicit}, the crown of this section.

\begin{definition}{Treap = BST on keys, heap on random priorities}{c06-treap}
Each node holds a key $k$ and a priority $p$ drawn independently at random (assume distinct, an event of
probability $1$). The tree is a \term{treap} if it is a binary search tree on the $k$'s \emph{and} a
max-heap on the $p$'s (every node's priority exceeds its children's). Given a set of distinct
key/priority pairs there is \textbf{exactly one} treap: its root is the maximum-priority key, and its two
subtrees are the treaps of the keys below and above it. Building it is the Cartesian-tree construction,
which we take as a black box here (\chref{c26}).
\end{definition}

\noindent\textbf{In plain words.} Two orderings, one tree: sorted by key left-to-right, best-priority on
top. The randomness that decides \emph{who becomes root} is the same randomness that makes the shape
balanced -- you never touch it again.

\begin{theorem}{A treap has height $\bigO{\log n}$ in expectation}{c06-exptheight}
For $n$ keys with i.i.d.\ continuous priorities, the treap's height $H_n$ satisfies
$\mathbb{E}[H_n] = \bigO{\log n}$ (and, in fact, $H_n \sim c\,\lg n$ with high probability).
\end{theorem}
\begin{proofbox}[Treap $\equiv$ random BST]
The max-priority key is uniformly one of the $n$ keys, so it is the root with each rank equally likely;
conditioned on the root being rank $K$, the left and right subtrees are treaps on $K-1$ and $n-K$ keys
with independent random priorities. That is \emph{precisely} the recursion for a binary search tree built
by inserting keys in uniformly random order (the first inserted -- here the max priority -- is the root).
It is classical that a random BST has expected height $\bigO{\log n}$ (equivalently $\sim 4.311\,\ln n$):
the depth of a fixed key is the number of its \emph{ancestors}, and a key is an ancestor of $x$ iff it is
the first (highest priority) among the interval it spans, an event of probability $1/j$, so expected depth
$= \sum_j 1/j = \bigO{\log n}$; a union/marked argument lifts this to the height. Since the treap and the
random BST are the same distribution, so is their height. \end{proofbox}
\begin{insightbox}
This is the whole selling point: \emph{you do not balance a treap, you \textbf{sample} it.} There are no
balance factors and no rotation cases -- insert and delete re-establish the heap property automatically
through \texttt{split}/\texttt{merge} (\defref{c06-splitmerge}). The price is that the bound is
\emph{expected}, not worst-case: an astronomically unlucky coin flip could give a tall tree. In contests a
64-bit priority makes this risk nil; when you need a hard guarantee, pay AVL's case work instead.
\end{insightbox}

\subsection{The universal toolkit: split and merge}

\begin{definition}{\texttt{split} and \texttt{merge}}{c06-splitmerge}
\emph{$\texttt{merge}(a,b)$} --- precondition: every key of $a$ is less than every key of $b$ (both are
treaps). It returns the treap on $a\cup b$: compare the roots' priorities; the winner becomes the root and
recursively merges into its free (right/left) child, then updates aggregate fields.
\emph{$\texttt{split}(t,k,a,b)$} --- cut treap $t$ into $a$ = keys $\le k$ (or first $k$ positions, in the
implicit form) and $b$ = the rest, both treaps, by walking the search path for $k$ and re-threading.
Both run down a single root-to-leaf path, hence expected $\bigO{\log n}$.
\end{definition}

\begin{lemma}{\texttt{merge} preserves the treap invariants}{c06-mergevalid}
If $a$ and $b$ are treaps with $\max(a.\text{keys}) < \min(b.\text{keys})$, then $\texttt{merge}(a,b)$ is
the unique treap on $a\cup b$: it stays a BST (all $a$ precede all $b$ in inorder) and a heap (the winner
of the two root priorities is placed on top, and recursion applies the same rule). Dually,
$\texttt{split}$ returns two valid treaps whose inorder is $t$'s inorder cut at the boundary.
\end{lemma}
\begin{proofbox}
In $\texttt{merge}(a,b)$ suppose $a.\mathit{pri} > b.\mathit{pri}$ (max-heap). Then $a.\mathit{root}$ is the
highest-priority key of $a\cup b$ (as it already beat all of $a$, and $b$'s top lost to it), so the heap
rule forces it to be the treap's root; by the BST rule all $b$'s keys, being larger than every $a$-key,
must lie in the right subtree, and $b$ attaches into $a.\mathit{root}$'s right -- the recursive
$\texttt{merge}(a.\mathit{right}, b)$ is the identical problem one level down. Induction gives the full
tree; the $\texttt{upd}$ of aggregate fields keeps sizes/sums correct. $\texttt{split}$ is the inverse
re-threading along $k$'s search path: subtrees wholly below $k$ are handed to $a$, wholly above to $b$, and
each handed-over pointer keeps a BST block intact and a heap subtree untouched. \end{proofbox}

\begin{lstlisting}[caption={Treap node, aggregate update, and merge -- the two rules maintained by recursion, no rotations}, label={lst:c06-treap}]
mt19937 rng(chrono::steady_clock::now().time_since_epoch().count());
struct Node { int key, pri, sz; Node *l, *r;
    Node(int k) : key(k), pri(rng()), sz(1), l(nullptr), r(nullptr) {} };

int  SZ(Node* t) { return t ? t->sz : 0; }
void upd(Node* t) { if (t) t->sz = 1 + SZ(t->l) + SZ(t->r); }

// Precondition: every key in a < every key in b.
Node* merge(Node* a, Node* b){
    if (!a) return b;
    if (!b) return a;
    if (a->pri > b->pri) { a->r = merge(a->r, b); upd(a); return a; } // a's root is global max-pri
    else                 { b->l = merge(a, b->l); upd(b); return b; }
}
\end{lstlisting}

\begin{lstlisting}[caption={Split by key (and its implicit by-position variant) -- cut a treap into two valid treaps}, label={lst:c06-split}]
// a gets keys <= k, b gets keys > k.
void splitKey(Node* t, int k, Node*& a, Node*& b){
    if (!t) { a = b = nullptr; return; }
    if (t->key <= k) { splitKey(t->r, k, t->r, b); a = t; upd(a); }  // t in a, its right splits
    else             { splitKey(t->l, k, a, t->l); b = t; upd(b); }  // t in b, its left splits
}
// IMPLICIT split: a gets the first cnt positions (inorder), b the rest.
void splitPos(Node* t, int cnt, Node*& a, Node*& b){
    if (!t) { a = b = nullptr; return; }
    if (SZ(t->l) >= cnt) { splitPos(t->l, cnt, a, t->l); b = t; upd(b); }
    else                 { splitPos(t->r, cnt - SZ(t->l) - 1, t->r, b); a = t; upd(a); }
}
\end{lstlisting}

\noindent Every keyed operation is two calls deep. \texttt{insert(x)}: split into $L,R$ with
\texttt{splitKey(root,x,L,R)}, then merge $L$, a fresh node, and $R$. \texttt{erase(x)}: split out
$\{x\}$ into its own treap and \texttt{merge} the two sides around it. \texttt{rank(x)}: split at $x$,
read $\texttt{SZ}(L)$. All expected $\bigO{\log n}$, all rotation-free -- contrast the four cases of
\coderef{c06-avlrot}.

\begin{figure}[H]
\centering
\begin{tikzpicture}[diagram, scale=0.82, transform shape]
  % left: treap on 1..7 with lucky priorities = balanced twin
  \node[troot] (r4) at (0,0)    {$4$};
  \node[tnode] (r2) at (-1.6,-1.3) {$2$};
  \node[tnode] (r6) at (1.6,-1.3)  {$6$};
  \node[tleaf] (r1) at (-2.4,-2.6) {$1$};
  \node[tleaf] (r3) at (-0.8,-2.6) {$3$};
  \node[tleaf] (r5) at (0.8,-2.6)  {$5$};
  \node[tleaf] (r7) at (2.4,-2.6)  {$7$};
  \draw (r4) -- (r2); \draw (r4) -- (r6); \draw (r2) -- (r1); \draw (r2) -- (r3);
  \draw (r6) -- (r5); \draw (r6) -- (r7);
  \node[annot, align=left, font=\scriptsize] at (0,1.1) {treap on $1..7$ (max-pri root $4$)};
  % arrow
  \draw[flow, double] (3.6,-1.3) -- (5.0,-1.3) node[midway, above, font=\scriptsize] {splitPos(t,3)};
  % right: a = {1,2,3}, b = {4,5,6,7}
  \node[tnode] (a2) at (6.4,-0.4) {$2$};
  \node[tleaf] (a1) at (5.6,-1.7) {$1$};
  \node[tleaf] (a3) at (7.2,-1.7) {$3$};
  \draw (a2) -- (a1); \draw (a2) -- (a3);
  \node[annot, font=\scriptsize] at (6.4,0.4) {$a$ = first $3$};
  \node[tnode] (b4) at (10.2,-0.4) {$4$};
  \node[tnode] (b6) at (11.4,-1.7) {$6$};
  \node[tleaf] (b5) at (10.6,-3.0) {$5$};
  \node[tleaf] (b7) at (12.2,-3.0) {$7$};
  \draw (b4) -- (b6); \draw (b6) -- (b5); \draw (b6) -- (b7);
  \node[annot, align=left, font=\scriptsize] at (9.9,0.4) {$b$ = rest (positions $4..7$)};
\end{tikzpicture}
\caption{Implicit \texttt{splitPos} cuts a sequence-treap by \emph{position}, not value: splitting the
treap at $\mathit{cnt}=3$ hands the first three inorder keys $\{1,2,3\}$ to $a$ and $\{4,5,6,7\}$ to $b$,
each still a valid treap (\lemref{c06-mergevalid}). Merging $a$ and $b$ back restores the original -- so
``delete a range'', ``paste after position $k$'', ``reverse $[l,r]$'' are all a split, a local edit, and a
merge. This is the dynamic-array superpower of \propref{c06-implicit}.}\label{fig:c06-splitfig}
\end{figure}

\subsection{Implicit keys: the treap as a dynamic sequence}

\begin{proposition}{Implicit keys make a treap a dynamic array}{c06-implicit}
Suppose no key is stored, and a node's logical index is its inorder position, recovered as
$\mathit{idx}(v)=\texttt{SZ}(v.l)$ at the root. Then $\texttt{splitPos}$ cuts the sequence at an
\emph{index} and $\texttt{merge}$ concatenates two sequences, so a treap supports, each in expected
$\bigO{\log n}$: insert at position, erase at position, extract/concat a range, and (with one lazy flag
swapping children) reverse a range -- operations no keyed BST offers directly.
\end{proposition}
\begin{proofbox}
With implicit indexing the BST rule is replaced by the invariant ``inorder = sequence order,'' which
$\texttt{splitPos}$/ $\texttt{merge}$ maintain because they only ever re-thread along inorder boundaries
($\texttt{merge}$ still requires one sequence to precede the other, which is exactly concatenation). The
heap-on-priorities rule is untouched, so balance is still $\bigO{\log n}$ expected
(\thmref{c06-exptheight}) regardless of the positions you edit. $\texttt{SZ}$ tells you, at any node, how
many elements lie to its left, so a descent by count finds the $k$-th position in $\bigO{h}$. Reversing
$[l,r]$ is: split to isolate the $r-l+1$ middle nodes, toggle a lazy ``reverse'' flag on its root (flip
left/right children, mark children to be flipped on push), and merge back -- the flag makes it $\bigO{\log
n}$ rather than touching every element. \end{proofbox}

\begin{lstlisting}[caption={Implicit sequence treap -- insert/erase by position and a lazy range reverse (skeleton)}, label={lst:c06-imp}]
struct Node { int pri, sz, val; bool rev; Node *l, *r; /* ... */ };
void push(Node* t){                       // propagate a pending reverse before descending
    if (t && t->rev){ swap(t->l, t->r);
        if (t->l) t->l->rev ^= 1; if (t->r) t->r->rev ^= 1; t->rev = false; }
}
// every recursion (splitPos / merge / upd) calls push(t) first -- lazy reverse stays O(log n).
Node* insert(Node* t, int pos, int val){  // insert val before current position pos
    Node *a, *b; splitPos(t, pos, a, b);
    return merge(merge(a, new Node(val)), b);
}
Node* erase(Node* t, int pos){            // delete the element at position pos
    Node *a, *mid, *b; splitPos(t, pos, a, mid); splitPos(mid, 1, mid, b);
    return merge(a, b);                   // mid is dropped
}
Node* reverse(Node* t, int lo, int hi){   // reverse positions [lo, hi)
    Node *a, *m, *b; splitPos(t, hi, m, b); splitPos(m, lo, a, m);
    if (m) m->rev ^= 1;                   // O(1) lazy flip on the isolated block
    return merge(merge(a, m), b);
}
\end{lstlisting}

\noindent\textbf{Lazy note.} The reverse flag is a one-line taste of \emph{lazy propagation}; its full
theory (push/pull discipline, when aggregates survive a pending flip) is \chref{c15}. For LC 307 below we
need no laziness -- only $\mathit{sum}$ augmentation, updated in the very same \texttt{upd} call that
maintains $\mathit{sz}$.

\subsection{Worked Example}

\begin{example}{A lucky draw, a split, and a mutable range sum}{c06-treap-run}
\emph{Treap on $C$.} Draw priorities so that key $4$ is the max, then $2$ the max of $\{1,2,3\}$ and $6$
the max of $\{5,6,7\}$. The unique treap is exactly the balanced twin $4\to(2\to(1,3),\,6\to(5,7))$ of
\figref{c06-chain} -- but AVL had to \emph{rotate four times} to reach it (\exref{c06-avl-run}), whereas
here the coin flip simply \emph{lands} there, and any other distinct-priority draw is also fine (still
$\bigO{\log n}$ expected). No cases were inspected.\\
\emph{Split/merge on the sequence.} Treat that treap as the sequence $[1,2,3,4,5,6,7]$.
$\texttt{splitPos}(t,3)$ yields $a=[1,2,3]$ and $b=[4,5,6,7]$ (\figref{c06-splitfig}); inserting $99$ at
position $1$ is $\texttt{merge}(a.\text{first }1,\ \text{node }99,\ \text{rest})$ giving
$[1,99,2,3]\Vert[4,5,6,7]$; $\texttt{merge}(a,b)$ restores order.\\
\emph{LC 307 micro-run} on $[1,3,5]$: build a treap with $\mathit{sum}$ fields (root sum $9$).
$\texttt{sumRange}(0,2)$: split into $a=[]$, $\text{mid}=[1,3,5]$, $b=[]$; answer $\texttt{mid->sum}=9$;
merge back. $\texttt{update}(1,2)$: split to isolate position $1$, set its value $3\to2$, $\texttt{upd}$
recomputes sums ($\text{mid->sum}$ was $3$, node value $2$; the root sum becomes $1+2+5=8$); merge.
$\texttt{sumRange}(0,2)$ now returns $8$ -- matching the official example.
\end{example}

\begin{tracebox}[One \texttt{update}(1,2) then \texttt{sumRange}(0,2) on the array $1,3,5$, by position]
Start $[1,3,5]$, $\mathit{sz}=(3)$, root $\mathit{sum}=9$.\\
\texttt{update(1,2)}: \texttt{splitPos(root,1, a, b)} $\Rightarrow$ $a=[1]$, $b=[3,5]$;
\texttt{splitPos(b,1, mid, c)} $\Rightarrow$ $mid=[3]$, $c=[5]$; set $mid->val=2$, \texttt{upd}
$\Rightarrow$ $mid$ sum $=2$; \texttt{root}=merge($a$,merge($mid$,$c$))$=[1,2,5]$, root sum $=8$.\\
\texttt{sumRange(0,2)}: \texttt{splitPos(root,0, a, b)} $\Rightarrow$ $a=[]$, $b=[1,2,5]$;
\texttt{splitPos(b,3, mid, c)} $\Rightarrow$ $mid=[1,2,5]$, $c=[]$; answer $=\texttt{mid->sum}=8$;
merge back. \checkmark
\end{tracebox}

\subsection{Variations and Twists}

\begin{itemize}[nosep, leftmargin=1.4em]
  \item \textbf{Min-heap vs max-heap priorities} is pure convention; just be consistent in
        \texttt{merge}'s comparison. What matters is that the priority is independent of the key. \hfill\lv{4}
  \item \textbf{Deterministic priorities from the key's hash} -- no RNG, still ``random enough'' and
        reproducible; a common contest trick to avoid flaky \code{rand}. \hfill\lv{4}
  \item \textbf{Persistent treap} -- since \texttt{split}/\texttt{merge} copy $\bigO{\log n}$ nodes along
        the path, path-copying yields a \emph{functional} set/map with history. \hfill\lv{5}
  \item \textbf{Splay instead of random} -- \secref{c06-splay} reaches $\bigO{\log n}$ amortised
        \emph{deterministically} by rotating accessed nodes to the root, but without split/merge's clean
        algebra. \hfill\lv{4}--\lv{5}
\end{itemize}

\subsection{Common Pitfalls}\label{sec:c06-treaps-pitfalls}

\begin{warningbox}[Trap 1 -- \texttt{merge} with overlapping key ranges]
\lemref{c06-mergevalid}'s precondition ``all keys of $a$ $<$ all keys of $b$'' is a \emph{contract}. Merge
two treaps that interleave and the BST rule silently breaks -- you lose or duplicate keys. In the implicit
setting the contract is ``$a$'s whole sequence precedes $b$'s''; only use \texttt{merge} to concatenate in
order.
\end{warningbox}

\begin{warningbox}[Trap 2 -- forgetting \texttt{upd} after re-threading]
Both \texttt{split} and \texttt{merge} change $\mathit{sz}$ (and $\mathit{sum}$, and any aggregate) on the
path. Skip the \texttt{upd(t)} on the way back and sizes go stale -- the \emph{next} \texttt{splitPos}
counts positions wrong and cuts mid-subtree. Update \emph{after} recursion, at the node you return, exactly
as in \coderef{c06-treap} and \coderef{c06-split}.
\end{warningbox}

\begin{warningbox}[Trap 3 -- the implicit \texttt{splitPos} off-by-one]
When the left subtree has $\texttt{SZ}(t.l) \ge \mathit{cnt}$ elements the cut lies inside $t.l$; otherwise
you recurse into $t.r$ asking for $\mathit{cnt} - \texttt{SZ}(t.l) - 1$ (subtract the left block \emph{and}
$t$ itself). A slipped $-1$ produces a range query that is one element long or short -- a \textbf{WA} that
survives small hand-tests and fails hidden ones.
\end{warningbox}

\subsection{Problem Lab}

The anchor is LeetCode 307, the canonical ``mutable range sum'' -- chosen to be solved with the
\emph{implicit treap} rather than the more common Fenwick/segment tree, so that split/merge and the
$\mathit{sum}$ augmentation are exercised end to end. (A Fenwick tree \chref{c14} is the shorter,
deterministic solution here; the point of doing it as a treap is to prove the toolkit generalises.)

% ---------------------------------------------------------------------------
\solvedproblem{c06-rpmq}{Range Sum Query -- Mutable}{LeetCode 307}
  {https://leetcode.com/problems/range-sum-query-mutable/}{LC Medium L3}
  {implicit treap, split/merge by position, subtree-sum augmentation, point update and range query}

\begin{problemstatement}[Range Sum Query -- Mutable]
Implement a class \code{NumArray} over an integer array \code{nums} supporting two operations, called up
to $3\cdot 10^{4}$ times on an array of length up to $3\cdot 10^{4}$ with values in $[-100,100]$:
\texttt{update(index, val)} sets \code{nums[index]} to \code{val}, and \texttt{sumRange(left, right)}
returns the sum of \code{nums[left..right]} (inclusive). Both should be $\bigO{\log n}$.\\
\emph{Example:} \code{NumArray([1,3,5])}; \code{sumRange(0,2)} $=9$; \code{update(1,2)}; \code{sumRange(0,2)}
$=8$.
\end{problemstatement}

\begin{attemptgate}[Attempt it yourself first -- \lv{4}, give yourself 30 minutes]
Model the array as an implicit treap keyed by \emph{position}; each node carries its subtree $\mathit{sum}$.
\peek{To sum positions $[\mathit{left},\mathit{right}]$, you want that block to be a single subtree so its
stored $\mathit{sum}$ is the answer. Which two operations isolate a contiguous block of positions?}
\nudge{Two \texttt{splitPos}: first cut at \code{left}, then cut the remainder at
\code{right-left+1}. The middle treap \emph{is} the range; read \code{mid->sum}, then \texttt{merge} the
three back. For \texttt{update}, isolate the one position as a size-$1$ treap and change its value.}
\fullhint{Store \code{sum} and \code{sz}; \texttt{upd(t)} sets
\code{sum = SM(l) + val + SM(r)} and \code{sz}. \texttt{sumRange}:
\code{splitPos(root,left,a,b); splitPos(b,right-left+1,mid,c); ans=SM(mid);
root=merge(a,merge(mid,c));}. \texttt{update}:
\code{splitPos(root,index,a,b); splitPos(b,1,mid,c); mid->val=val; upd(mid);
root=merge(a,merge(mid,c));}. Expected $\bigO{\log n}$ per call by \thmref{c06-exptheight}.}
\end{attemptgate}

\solstep{Understanding the Problem}
We must maintain a sequence under point writes and answer contiguous-range sums, both fast. A plain array
gives $\bigO{1}$ update but $\bigO{n}$ sum (or $\bigO{1}$ sum via prefix arrays but $\bigO{n}$ update when
a value changes all downstream prefixes). We want $\bigO{\log n}$ for \emph{both}. The implicit treap fits
because a range is a contiguous run of \emph{positions}, and \texttt{splitPos} isolates a position range
into one subtree; if each node caches its subtree sum, the answer is that one cached value -- no traversal.
Update is ``isolate the single position, rewrite it, fix sums on the way back up.''

\solstep{Key Observations}
\begin{observation}
A subtree of an implicit treap corresponds to a \emph{contiguous slice} of the sequence, because inorder =
sequence order. So caching $\mathit{sum}$ per node answers every range-sum-in-a-subtree instantly.
\end{observation}
\begin{observation}
$\texttt{splitPos}(\cdot,\mathit{left})$ then $\texttt{splitPos}(\cdot,\mathit{right}-\mathit{left}+1)$
cuts exactly $[\mathit{left},\mathit{right}]$ into the middle treap; merging the three pieces back restores
the whole array. Point update is the same cut to width $1$.
\end{observation}
\begin{observation}
Since priorities are random, the tree has expected height $\bigO{\log n}$ (\thmref{c06-exptheight}) and each
split/merge walks one path -- so both operations are expected $\bigO{\log n}$ with \emph{no} balancing code
and no lazy propagation (values, not ranges, change here).
\end{observation}

\solstep{Think It Through With Me}
\begin{thinkbox}
\deadend{Rebuild a prefix-sum array after every update.}
{$\bigO{n}$ per update, and up to $3\cdot 10^{4}$ updates $\Rightarrow$ $\bigO{nq}\sim 9\cdot 10^{8}$,
borderline \tle and against the $\bigO{\log n}$ goal. Prefix sums only pay off when updates are rare.}
\deadend{Do it as a keyed treap with $\mathit{key}=\mathit{index}$.}
{Works, but you must keep the indices as literal keys and they never renumber -- fine for \emph{this}
problem since positions are fixed. The implicit (position) view is used here because it generalises to
insert/erase/reverse, which the keyed view cannot do without re-keying every node.}
\aha{Use an implicit treap keyed by subtree size; augment each node with $\mathit{sum}$. Range sum $=$ two
splits + read middle sum + two merges; update $=$ split to width $1$, rewrite, merge back. All $\bigO{\log
n}$ expected, all rotation-free.}
\end{thinkbox}

\solstep{Solution}
Build the treap by merging nodes in order (each new element appended on the right), maintaining $\mathit{sz}$
and $\mathit{sum}$ in a single \texttt{upd}. \texttt{update(index,val)}: split the root at $\mathit{index}$
into $a,b$; split $b$ at $1$ into $mid,c$; set $mid$'s value to $\mathit{val}$ and \texttt{upd} it; reassemble
\texttt{root=merge(a,merge(mid,c))}. \texttt{sumRange(left,right)}: split at $\mathit{left}$; split the tail
at $\mathit{right}-\mathit{left}+1$; read \texttt{SM(mid)}; reassemble. Sums are $\bigO{1}$ to read because
every subtree is a contiguous slice (\thmref{c06-exptheight} keeps the pieces logarithmic in depth).

\solstep{Proof of Correctness}
\begin{proofbox}[Sums are exact and positions are preserved]
Invariant: for every node $t$, $\mathit{sz}(t)$ equals the number of sequence elements in $t$'s subtree and
$\mathit{sum}(t)$ equals the sum of their values, in inorder. \texttt{upd} establishes it from the children,
and \texttt{merge}/\texttt{splitPos} call \texttt{upd} on exactly the re-threaded nodes (Trap~2), so the
invariant holds after every operation; untouched subtrees keep correct fields by the induction hypothesis.
Because inorder equals sequence order and \texttt{splitPos}(t,k,a,b) is defined to put the first $k$
positions in the left result, after \texttt{splitPos(root,left,a,b)} and
\texttt{splitPos(b,right-left+1,mid,c)} the treap $mid$'s inorder is \emph{precisely} the slice
$\mathit{nums[left..right]}$ (a contiguous slice, so by the invariant its root's $\mathit{sum}$ is the sum of
that slice). Thus the returned \texttt{SM(mid)} is the correct range sum; merging $a,mid,c$ back in that
order restores the identical sequence (\lemref{c06-mergevalid}), leaving the structure ready for the next
call. \texttt{update} isolates a width-$1$ $mid$ whose inorder is exactly $\mathit{nums[index]}$; rewriting
its value and \texttt{upd}-ing makes every ancestor's $\mathit{sum}$ reflect the change, and re-merging
restores the sequence. Hence all answers equal the array's true range sums. \end{proofbox}

\solstep{Implementation}
\begin{lstlisting}[caption={Range Sum Query -- Mutable via an implicit, sum-augmented treap (LeetCode 307)}, label={lst:c06-rpmq-soln}]
#include <bits/stdc++.h>
using namespace std;
mt19937 rng(chrono::steady_clock::now().time_since_epoch().count());

struct Node { int val, pri, sz, sum; Node *l, *r;
    Node(int v) : val(v), pri(rng()), sz(1), sum(v), l(nullptr), r(nullptr) {} };
int  SZ(Node* t){ return t ? t->sz : 0; }
int  SM(Node* t){ return t ? t->sum : 0; }
void upd(Node* t){ if(t){ t->sz = 1+SZ(t->l)+SZ(t->r); t->sum = SM(t->l)+t->val+SM(t->r); } }

Node* merge(Node* a, Node* b){
    if(!a) return b; if(!b) return a;
    if(a->pri > b->pri){ a->r = merge(a->r, b); upd(a); return a; }
    else               { b->l = merge(a, b->l); upd(b); return b; }
}
void splitPos(Node* t, int cnt, Node*& a, Node*& b){   // a = first cnt positions
    if(!t){ a=b=nullptr; return; }
    if(SZ(t->l) >= cnt){ splitPos(t->l, cnt, a, t->l); b=t; upd(b); }
    else               { splitPos(t->r, cnt-SZ(t->l)-1, t->r, b); a=t; upd(a); }
}

class NumArray {
    Node* root = nullptr;
public:
    NumArray(vector<int>& nums){
        for (int v : nums) { Node* t = new Node(v); root = merge(root, t); }
    }
    void update(int index, int val){
        Node *a, *mid, *b;
        splitPos(root, index, a, mid);
        splitPos(mid, 1, mid, b);      // mid is exactly position index
        mid->val = val; upd(mid);
        root = merge(a, merge(mid, b));
    }
    int sumRange(int left, int right){
        Node *a, *mid, *b;
        splitPos(root, left, a, mid);
        splitPos(mid, right-left+1, mid, b);   // mid is nums[left..right]
        int res = SM(mid);
        root = merge(a, merge(mid, b));
        return res;
    }
};
\end{lstlisting}

\solstep{Trace}
\begin{tracebox}[Official example, array $1,3,5$]
Constructor merges $[1]$, then $[1,3]$, then $[1,3,5]$; whichever becomes root, its $\mathit{sum}=9$.\\
\texttt{sumRange(0,2)}: \texttt{splitPos(root,0, a, mid)} $\Rightarrow$ $a=\emptyset$, $mid=[1,3,5]$;
\texttt{splitPos(mid,3, mid, b)} $\Rightarrow$ $mid=[1,3,5]$, $b=\emptyset$; $res=\mathit{SM}(mid)=9$;
merge back. Returns $\boxed{9}$.\\
\texttt{update(1,2)}: \texttt{splitPos(root,1, a, mid)} $\Rightarrow$ $a=[1]$, $mid=[3,5]$;
\texttt{splitPos(mid,1, mid, b)} $\Rightarrow$ $mid=[3]$, $b=[5]$; $mid->val=2$, $mid->sum=2$;
rebuild $[1,2,5]$ with root sum $8$.\\
\texttt{sumRange(0,2)}: isolate $[1,2,5]$, $\mathit{SM}=8$, restore. Returns $\boxed{8}$. \checkmark
\end{tracebox}

\solstep{Complexity}
\begin{complexitybox}
Build $\bigO{n\log n}$ expected ($n$ appends, each a $\bigO{\log n}$ merge); a tighter linear build via a
Cartesian stack exists but is \chref{c26}'s tool. Each \texttt{update} and \texttt{sumRange} does two splits
and two merges, each expected $\bigO{\log n}$ by \thmref{c06-exptheight}, so $\bigO{\log n}$ per call,
$\bigO{q\log n}$ overall ($q\le 3\cdot 10^{4}$). Memory $\bigO{n}$ nodes. Deterministic rivals Fenwick
\chref{c14} / segment tree \chref{c13} do the same $\bigO{\log n}$ with less code \emph{for this problem};
the treap earns its keep the moment you also need insert/erase/reverse by position. Verdict: \ac.
\end{complexitybox}

\solstep{Pitfalls and Variants}
\begin{warningbox}
(1) The $-1$ in \texttt{splitPos}(t->r, \texttt{cnt-SZ(t->l)-1}, ...) (Trap~3). (2) \texttt{upd} after every
re-thread (Trap~2). (3) If a whole-array \texttt{sumRange} is queried often, reading the root's $\mathit{sum}$
without splitting is a valid $\bigO{1}$ shortcut when $\mathit{left}=0,\mathit{right}=n-1$. \emph{Variants:}
swap $\mathit{sum}$ for \emph{any} associative aggregate (min, max) to get range-min under updates; swap the
values for the four-tuple $(\mathit{total},\mathit{pref},\mathit{suff},\mathit{best})$ and you have the
max-subarray monoid -- exactly the CSES 1190 practice below.
\end{warningbox}

\begin{insightbox}[Key takeaways]
\begin{enumerate}[nosep, leftmargin=1.6em]
  \item \textbf{Randomise, don't rebalance.} Heap-ordered random priorities make the treap a random BST
        with expected height $\bigO{\log n}$ -- no balance factors, no rotation cases (\thmref{c06-exptheight}).
  \item \textbf{Split and merge are the whole algebra.} Insert, erase, rank, and range extraction are
        two-line compositions (\defref{c06-splitmerge}, \lemref{c06-mergevalid}).
  \item \textbf{Index by size, get a dynamic array.} Implicit keys turn the treap into a sequence with
        insert/delete/reverse by position -- and sum-augmentation answers range queries instantly
        (\propref{c06-implicit}, \probref{c06-rpmq}).
\end{enumerate}
\end{insightbox}

\begin{hardeningbox}
\begin{itemize}[nosep, leftmargin=1.4em]
  \item \textbf{Range reverse / rotate} a subarray by isolating it and toggling the lazy flag -- see the
        lazy-propagation chapter \chref{c15} for discipline.
  \item \textbf{Ordered statistics} -- keyed treap with $\mathit{sz}$ gives rank and select in $\bigO{\log n}$,
        replacing a Fenwick over a \emph{dynamic} key set.
  \item \textbf{Rope / editor buffers} -- concatenating and splitting sequences is \texttt{merge}/\texttt{split}
        verbatim; this is why editors use treap-family trees.
\end{itemize}
\end{hardeningbox}

\subsection{Practice Problems}

\begin{stuckbox}
\begin{itemize}[nosep, leftmargin=1.4em]
  \item Did you read the middle treap's cached aggregate \emph{without} descending into it? (A range should be
        one subtree -- \thmref{c06-exptheight} keeps it shallow anyway.)
  \item In \texttt{splitPos}, did you subtract the left block \emph{and} the node when going right? (Trap~3.)
  \item After editing a value, did you \texttt{upd} the isolated node before re-merging? (Trap~2.)
\end{itemize}
\end{stuckbox}

\begin{practiceset}
  \practice{Subarray Sum Queries}{CSES 1190}
    {https://cses.fi/problemset/task/1190}{CSES · L3}
    {Point updates on an array; after each, print the maximum subarray sum (empty allowed, i.e.\ at least
    $0$). An implicit treap where each node stores the monoid $(\mathit{total},\mathit{pref},\mathit{suff},
    \mathit{best})$ answers the whole-array query as \texttt{root->best}; \texttt{merge} combines two
    children's tuples in constant time, and an update is the LC 307 isolate-rewrite-merge. The same problem is
    the classic segment-tree monoid application (\chref{c13}) -- solve it \emph{both} ways to feel why the
    treap's split/merge is more flexible (you can also \emph{insert} positions, which a segment tree cannot).}
\end{practiceset}

\begin{recallbox}[10-second refresher -- 8.3 Treaps, Split/Merge, Implicit Keys]
A treap is a BST on keys and a heap on random priorities, hence uniquely determined and, being a random BST,
of expected height $\bigO{\log n}$ (\defref{c06-treap}, \thmref{c06-exptheight}) -- balance by sampling, not
rotating. Two rotation-free routines do everything: \texttt{merge} (needs all of $a$ before all of $b$) and
\texttt{split} (\defref{c06-splitmerge}, \lemref{c06-mergevalid}); insert/erase/rank are one-liners over them.
Index by subtree size instead of stored keys and the treap becomes a dynamic sequence -- insert, delete,
reverse a range by position (\propref{c06-implicit}, \coderef{c06-imp}). Cache a $\mathit{sum}$ per node and
range sums come free, which is LC 307 done the treap way (\probref{c06-rpmq}).
\selftest{
  \item Why is the treap on $n$ random priorities distributed like a random BST, and what height follows?
  \item State \texttt{merge}'s precondition and what breaks if you violate it (Trap~1).
  \item Give the split/merge recipe to isolate positions $[\mathit{lo},\mathit{hi})$ and reverse them in
        $\bigO{\log n}$.
  \item In LC 307, why can you read \texttt{mid->sum} without traversing $mid$?
  \item Name a deterministic structure that solves LC 307 in less code, and one it \emph{cannot} extend to.
}
\end{recallbox}

% !TEX root = ../../tree_notes.tex
% =====================================================================
%  8.4  Red-Black Trees (conceptual, N-lite)     sec:c06-redblack
%  Exports: def:c06-rbcolor  thm:c06-rbheight  fig:c06-rbfix  ex:c06-rb-run
%           (no solved problem; practice: CF 1354D)
% =====================================================================
\section{Red-Black Trees (Why \code{std::set} Is Fast)}\label{sec:c06-redblack}

\begin{previouslybox}
\begin{enumerate}[nosep, leftmargin=1.6em]
  \item AVL keeps height $\le 1.44\,\lg(n+2)$ by enforcing $|\bal|\le 1$ and rotating up to twice per
        insertion. What did the balancing \emph{cost} you? (balance factors, height upkeep, the four
        rotation cases -- \defref{c06-bal}, \thmref{c06-fixheight}.)
  \item Chapter~7 asserted $\code{std::set}$/$\code{map}$ operations are $\bigO{\log n}$ and \emph{deferred}
        the proof. To which structure? (a self-balancing tree -- now named: red-black,
        \propref{c05-logn}.)
  \item What does $\bigO{\log n}$ \emph{depend on} -- the keys, or the height? (the height;
        \defref{c06-heightcost}.)
\end{enumerate}
\end{previouslybox}

\begin{topicintro}
This section is deliberately \emph{conceptual}: you will almost never implement a red-black tree by hand,
because your standard library already ships one -- but you should be able to \emph{explain} why a
\code{std::set} never degrades into the chain $C$. Red-black trees are AVL's pragmatic cousin. Instead of
demanding every node be within one level of balance, they paint nodes \term{red} or \term{black} and impose
five cheap colouring rules. Those rules let a path be \emph{twice} as long as another -- so the tree is
lazier than AVL -- yet still cap the height at $2\,\lg(n+1)$ (\thmref{c06-rbheight}), which is exactly the
$\bigO{\log n}$ promise Chapter~7 borrowed and now cashes. The colouring is also engineered to make edits
\emph{cheap}: an insertion that breaks a rule is repaired by recolouring and at most two rotations, and
recolouring a red node is a single-bit flip -- often cheaper than an AVL rotation, which is why the
\emph{write-heavy} containers chose red-black over AVL. We show the five invariants, prove the height
bound, walk the insertion repair cases, and connect back to \secref{c05-contests}. There is no
build-it-yourself lab here by design; the practice instead makes you \emph{use} an order statistic, the
one question a bare red-black \code{set} refuses to answer quickly.
\textbf{Prerequisites.} \secref{c06-rotations} (rotations), \secref{c06-avl} (what a strict balance rule
costs); \chref{c05}'s container section.

\textbf{After this section you can}
\begin{itemize}[nosep]
  \item state the five red-black invariants and why forbidding adjacent reds bounds the height; \hfill\lv{3}
  \item prove $\height \le 2\,\lg(n+1)$ from black-height, repaying \propref{c05-logn}; \hfill\lv{4}
  \item narrate the insertion repair: recolour when the uncle is red, rotate when it is black; \hfill\lv{4}
  \item explain the AVL-vs-red-black tradeoff and know \code{std::set} \emph{is} a red-black tree. \hfill\lv{3}
\end{itemize}
\end{topicintro}

\begin{recognitionbox}[Pattern card -- you need a balanced tree but you own none]
\reckey{an ordered set/map with $\bigO{\log n}$ insert/erase/find where \emph{using} the library beats
re-implementing, \emph{and} you must reason about the guarantee it gives.}
\recuse{
  \item ``store keys, query neighbours/coverage fast'' with a \code{std::set}/\code{map} -- that container
        is red-black (\secref{c05-contests}, Range Module);
  \item write-heavy workloads where a rotation is dearer than a flip -- red-black's lazy rebalancing wins
        over AVL;
  \item you must \emph{explain} why an indexed \code{map}'s guarantee holds under any input order.
}
\recnotuse{
  \item you need the $k$-th smallest \emph{fast} -- a red-black \code{set} does \textbf{not} expose rank
        cheaply; use an order-statistic tree or a Fenwick over values (the practice at the end of
        this section, \chref{c14});
  \item you need the tightest possible search height and mostly read -- AVL's $1.44\lg n$ beats red-black's
        $2\lg n$ (\thmref{c06-avlbound} vs \thmref{c06-rbheight});
  \item you want to code the simplest balanced tree yourself -- a treap (\secref{c06-treaps}) has no colour
        cases at all.
}
One-line test: \emph{are you leaning on a black-box balanced container and need to trust (or explain) its
$\bigO{\log n}$? That box is red-black; understand the five rules and you understand it.}
\end{recognitionbox}

\subsection{Intuition}

AVL said: ``the two subtrees of every node may differ by at most one level,'' and had to rotate often to
obey it. Red-black says something weaker but easier to maintain: ``along any root-to-leaf path, red nodes
may \emph{slip in} extra levels, but never two reds in a row, and every path must contain the same number
of \emph{black} nodes.'' A red edge is a borrowed half-level -- a child hanging at the same black-depth as
its parent. Because reds can never stack, the longest a path can stretch is to alternate black, red,
black, red\ldots, which is at most \emph{twice} the all-black shortest path. So the height is within a
factor $2$ of logarithmic -- bad enough to be lazy, good enough to be fast.

The magic is that this looser rule is \emph{cheap to restore}. When an insertion violates ``no two adjacent
reds,'' usually you can fix it just by recolouring a few nodes (a colour bit is O(1)) and pushing the
problem one level up -- only sometimes do you actually rotate. AVL, by contrast, must rotate to fix even
the smallest height imbalance. That is precisely why the language runtimes that prize fast \emph{writes}
(GCC's \code{std::map}, Java's \code{TreeMap}, and many others) implement red-black trees: fewer
structural changes per edit, at the cost of a slightly taller tree on reads.

\begin{figure}[H]
\centering
\begin{tikzpicture}[diagram, scale=0.82, transform shape]
  % helper styles via inline options
  \tikzset{bl/.style={circle, draw=black!70, fill=black!75, text=white, minimum size=6mm, font=\small\bfseries},
           rd/.style={circle, draw=warningred, fill=warningred!35, minimum size=6mm, font=\small\bfseries}}
  % (a) uncle RED -> recolour
  \node[annot, font=\scriptsize\bfseries] at (1.3,1.3) {(a) uncle red: recolour};
  \node[bl] (a2) at (1.3,0.4) {$20$};
  \node[rd] (a1) at (0.4,-0.7) {$10$};
  \node[rd] (a3) at (2.2,-0.7) {$30$};
  \node[rd] (an) at (0.0,-1.8) {$5$};
  \draw (a2) -- (a1); \draw (a2) -- (a3); \draw (a1) -- (an);
  \node[annot, align=left, font=\scriptsize] at (1.6,-2.5) {grandparent+uncle $\to$ black,\\ parent's line up $\to$ red};
  % arrow a->b
  \draw[flow] (3.4,-0.5) -- (4.1,-0.5);
  % (b) uncle BLACK, line -> single rotate
  \node[annot, font=\scriptsize\bfseries] at (6.2,1.3) {(b) black uncle, line: rotate};
  \node[bl] (b4) at (6.2,0.4) {$40$};
  \node[rd] (b2) at (5.1,-0.7) {$20$};
  \node[bl] (b6) at (7.3,-0.7) {$60$};
  \node[rd] (b1) at (4.4,-1.8) {$10$};
  \node[bl] (b3) at (5.8,-1.8) {$30$};
  \draw (b4) -- (b2); \draw (b4) -- (b6); \draw (b2) -- (b1); \draw (b2) -- (b3);
  \node[annot, align=left, font=\scriptsize] at (6.2,-2.6) {zig-zig: rotate the grandparent,\\ recolour parent/grandparent};
  % (c) bend below (mention)
  \node[annot, align=left, font=\scriptsize\itshape] at (9.3,-0.6) {(c) black uncle, bend\\ (zig-zag): the LR/RL\\ double rotation of\\ \defref{c06-rotcases} +\\ recolour. Same shapes\\ as AVL, different\\ \emph{when}.};
\end{tikzpicture}
\caption{Red-black insertion repair, given a red parent (which already breaks the no-adjacent-red rule with
the freshly-inserted red node). \textbf{(a)} If the \emph{uncle is red}, no rotation: recolour parent and
uncle black, grandparent red, and treat the grandparent as the new ``red node to check'' -- the fix hops
one level up. \textbf{(b)} If the uncle is \emph{black} and the red nodes form a straight \emph{line}, rotate
about the grandparent and swap two colours. \textbf{(c)} If they form a \emph{bend}, do the LR/RL double
rotation first to make a line, then case (b). Case (a) is red-black's signature -- it fixes by
\emph{flipping bits}, which AVL never gets to do.}\label{fig:c06-rbfix}
\end{figure}

\subsection{Definitions and Notation}

\begin{definition}{The five red-black invariants}{c06-rbcolor}
Every node is coloured \term{red} or \term{black}; the external leaves ($\nil$) are black. A binary search
tree is \term{red-black} if:
\begin{enumerate}[nosep, leftmargin=1.6em]
  \item every node is red or black;
  \item the root is black;
  \item every $\nil$ leaf is black;
  \item a red node has only black children (no two reds on any edge);
  \item every path from a node down to any $\nil$ leaf in its subtree contains the same number of black
        nodes (its \term{black-height} $\mathit{bh}(v)$).
\end{enumerate}
\end{definition}

\noindent\textbf{In plain words.} Rules 1--3 are bookkeeping. Rule 4 stops two reds stacking, so no path
doubling up; rule 5 makes all paths agree on how many \emph{black} steps they take. Together they say
``paths can differ, but only by red sprinkles that never double,'' and that is enough to keep the tree
logarithmic while making edits cheap.

\begin{theorem}{Red-black height is at most $2\,\lg(n+1)$}{c06-rbheight}
A red-black tree with $n$ internal (non-$\nil$) nodes has height $\height \le 2\,\lg(n+1)$. Consequently
membership, insertion and deletion in it are $\bigO{\log n}$ worst case -- the guarantee
\propref{c05-logn} promised for \code{std::set}/\code{map}.
\end{theorem}
\begin{proofbox}
Let $\mathit{bh}$ be the black-height of the root. \emph{(Lower bound on path length.)} The shortest root-to-leaf
path is all-black, so it has length at least $\mathit{bh}$ edges. \emph{(Upper bound.)} By rule 4 no two
consecutive nodes are red, so the longest path alternates black/red and has at most $2\mathit{bh}$ edges -- each of
the $\mathit{bh}$ black nodes is followed by at most one red. Hence the longest path is at most twice the shortest:
$\height \le 2\mathit{bh}$. \emph{(Black-height vs $n$.)} A subtree of black-height $b$ contains at least
$2^{b}-1$ internal nodes (induction: both children of the black root have black-height $\ge b-1$, so
$\ge 2(2^{b-1}-1)+1 = 2^{b}-1$). Thus $n \ge 2^{\mathit{bh}}-1$, giving $\mathit{bh} \le \lg(n+1)$. Combining,
$\height \le 2\mathit{bh} \le 2\,\lg(n+1)$. \end{proofbox}
\begin{insightbox}
Compare the constants: AVL $\le 1.44\,\lg(n+2)$ (\thmref{c06-avlbound}) versus red-black
$\le 2\,\lg(n+1)$ -- red-black's trees are up to $\sim 40\%$ taller in the worst case, so \emph{reads}
touch a few more nodes. But rule 4's slack means a red insertion is frequently repaired by \emph{recolouring}
(case (a) of \figref{c06-rbfix}) with no structural change at all. For \emph{balanced read/write} library
containers that tradeoff favours red-black, which is why \code{std::map} is red-black, not AVL.
\end{insightbox}

\subsection{The policy: fix upward, mostly by colour}

The repair in \figref{c06-rbfix} is applied after planting a red node at a failed-search leaf (the same
plant as \secref{c05-insert-delete}); only rule 4 can break, and only against a red parent. The loop
walks \emph{up} two levels at a time (to the grandparent) while the parent is red, choosing case (a), (b)
or (c) by the uncle's colour and the line-vs-bend geometry. It terminates in $\bigO{\log n}$ steps because
each recolouring pushes the violation toward the root, and at most \emph{two rotations} end the whole
repair (the classic bounds: insertion $\le 2$ rotations, deletion $\le 3$, plus $O(\log n)$ recolourings).
The BST descent and the rotations are literally the ones from \coderef{c06-rotcode}; only the
\emph{decision of when and which} differs -- red-black's answer is ``let colour absorb the slack, rotate
only when it cannot.''

\subsection{Worked Example}

\begin{example}{Colouring the chain $C$: where red-black is lazier than AVL}{c06-rb-run}
Insert $1,2,\dots,7$ (the order that made AVL rotate four times, \exref{c06-avl-run}). Red-black plants
each as a \emph{red} leaf and repairs upward. Ascending keys make every new node a right child, so we
tend to see the \emph{line} case (b): $1$ (black root), insert $2$ red -- fine; insert $3$ red onto red
parent $2$ with a black (absent, hence black-$\nil$) uncle -- case (b), \texttt{rotateLeft(1)} and swap
colours, giving a black $2$ with red children $1,3$. Continue: the tree at each stage is a valid
red-black tree of height at most $2\lg(n+1)$ and \emph{frequently only recolouring} is invoked where AVL
would have rotated (case (a) -- red uncle -- appears as soon as both sides grow). The final 7-node tree is
some red-black tree of height $\le 2\lg 8 = 6$... and in fact red-black lands at height $3$ with a couple
of red nodes slipped in -- \emph{looser} than AVL's height-$2$ balanced twin of \figref{c06-chain}, and
perfectly acceptable: still $\bigO{\log n}$. The lesson is visible on the same input: AVL insists on the
twin, red-black settles for ``twice the shortest'' and pays for it with fewer, cheaper edits. Inorder is
$1,2,\dots,7$ throughout -- colour never touches the sorted line (\thmref{c06-rotation}).
\end{example}

\begin{tracebox}[Why ``no two reds'' bounds the shape -- one recolour cascade]
Suppose a red node $z$ has red parent $p$ and red uncle $u$ (case a). Recolour $p,u$ black and the
grandparent $g$ red. Black-height is preserved on every path \emph{through} $g$ (they each gain a black
$p$ but $g$ turns red, net zero), so rule 5 is not broken. Now $g$ is red and may itself have a red
parent -- so we simply continue the loop at $g$. This is the ``fix travels up by colour'' mechanism: an
arbitrary number of these, then a single line/bend rotation that terminates, is all it takes -- $O(\log n)$
colour flips but $\le 2$ rotations total.
\end{tracebox}

\subsection{Variations and Twists}

\begin{itemize}[nosep, leftmargin=1.4em]
  \item \textbf{Left-leaning red-black (LLRB / Sedgewick)} -- a one-invariant red-black variant you can
        actually write in a page; same bound, fewer cases, at the cost of a small asymmetry. \hfill\lv{4}
  \item \textbf{Top-down red-black} -- one-pass insertion that pushes redness down as it descends, avoiding
        the upward fixup loop; this is roughly how \code{std::map} insert is structured. \hfill\lv{5}
  \item \textbf{Why not just always AVL?} If your workload is read-dominated and you can afford the tighter
        height (search-critical indexes), AVL's $1.44\lg n$ beats $2\lg n$; the choice is a genuine
        read/write engineering tradeoff, not a correctness one. \hfill\lv{3}--\lv{4}
  \item \textbf{Augmented red-black} -- add a subtree \emph{size} to nodes and the \code{set} gains
        order-statistics -- the missing ability probed by the practice problem. \hfill\lv{4}--\lv{5}
\end{itemize}

\subsection{Common Pitfalls}\label{sec:c06-redblack-pitfalls}

\begin{warningbox}[Trap 1 -- thinking red-black is ``balanced like AVL'']
It is not \emph{height}-balanced; it is \emph{black-height}-balanced. Two paths may legitimately differ by
up to a factor of two. Do not reason ``AVL's $1.44\lg n$'' for a \code{set}; the correct worst-case constant
is $2$ (\thmref{c06-rbheight}). Proofs and performance estimates must use the looser bound.
\end{warningbox}

\begin{warningbox}[Trap 2 -- expecting $k$-th smallest from \code{std::set} in $\bigO{\log n}$]
\code{std::set} gives lower/upper bound and neighbour queries fast, but the $k$-th element via
\texttt{std::next(s.begin(),k)} is $\bigO{k}$ -- the tree stores no subtree sizes. That is not a red-black
failure, just a non-augmented one: reach for an order-statistic tree (GNU \code{pb\_ds}) or a Fenwick over
the value range (\chref{c14}) -- see the practice below.
\end{warningbox}

\begin{warningbox}[Trap 3 -- recolouring that breaks black-height]
Every recolour must keep rule 5 true. Flipping a node without flipping its ``sibling'' changes the black
count on one path and silently destroys the balance guarantee (the tree can rot). Case (a) recolours
\emph{both} parent and uncle precisely to preserve black-height -- never recolour just one side.
\end{warningbox}

\subsection{Why You Almost Never Implement This}

The point of \secref{c06-redblack} in a contest toolkit is \emph{leverage, not code}. When you insert into
\code{std::set}, \code{std::map}, or \code{std::multiset}, you are already standing on a red-black tree:
the $\bigO{\log n}$ bounds Chapter~7 used throughout (\propref{c05-logn}, the whole of
\secref{c05-contests}) are \thmref{c06-rbheight} in disguise. The times you hand-roll a balanced tree are
exactly when the container's node layout is not enough -- you need an \emph{aggregate} (a sum, a size) at
each node, or \texttt{split}/\texttt{merge}. For those, prefer the treap (\secref{c06-treaps}): same
guarantee modulo expectation, a fraction of the cases. Red-black is the structure to \emph{understand}, not
to \emph{write}, which is why there is no lab here.

\begin{center}
\small
\renewcommand{\arraystretch}{1.25}
\begin{tabular}{@{}p{0.15\linewidth} p{0.36\linewidth} p{0.36\linewidth}@{}}
\toprule
 & \textbf{AVL (\secref{c06-avl})} & \textbf{Red-black (\secref{c06-redblack})} \\
\midrule
Rule     & $|\bal(v)|\le 1$ at every node (def.\ \defref{c06-bal}) & five colour invariants (def.\ \defref{c06-rbcolor}) \\
Height   & $\le 1.44\,\lg(n+2)$ (\thmref{c06-avlbound}) & $\le 2\,\lg(n+1)$ (\thmref{c06-rbheight}) \\
On edit  & rotate to restore height, $\le 2$ per insert   & recolour if possible, $\le 2$ rotations per insert \\
Good for & read-heavy, tighter search                       & write-heavy, fewer structural changes \\
Where    & in-memory indexes you tune                       & \code{std::set}/\code{map}, OS schedulers \\
\bottomrule
\end{tabular}
\end{center}

\subsection{Practice Problems}

\begin{stuckbox}
\begin{itemize}[nosep, leftmargin=1.4em]
  \item Can a red node have a red child? (No -- rule 4; that is the single invariant doing all the height
        work.)
  \item Does red-black give AVL's $1.44\lg n$, or the looser $2\lg n$? (The looser one,
        \thmref{c06-rbheight}.)
  \item For the practice: does \code{std::set} know the $k$-th smallest quickly? (No -- you must supply rank;
        see Trap~2.)
\end{itemize}
\end{stuckbox}

\begin{practiceset}
  \practice{Multiset}{Codeforces 1354D}
    {https://codeforces.com/problemset/problem/1354/D}{CF 1600 · L3}
    {A multiset starts with given values; a query either \emph{inserts} a value or \emph{removes the
    $k$-th smallest} (order statistic). A \code{std::multiset} handles insert and erase, but the
    $k$-th-smallest is the trap: \texttt{std::next} walks $\bigO{k}$, which \tle s (Trap~2). The clean
    fix is to store \emph{frequencies over the value domain} in a Fenwick tree
    (\chref{c14}) and find the smallest value whose prefix count reaches $k$ by binary lifting on the tree
    -- $\bigO{\log(\max\text{value})}$ per query -- or an order-statistic tree (GNU \code{pb\_ds}), which is
    exactly a red-black \code{set} with the subtree-size augmentation of Var.~4. This is the one question a
    bare \code{set} cannot answer, and it is why the chapter pairs the red-black guarantee
    (\thmref{c06-rbheight}) with augmentation (\secref{c06-treaps}).}
\end{practiceset}

\begin{recallbox}[10-second refresher -- 8.4 Red-Black Trees]
Red-black replaces AVL's tight height rule with five colour invariants (\defref{c06-rbcolor}); forbidding
adjacent reds makes the longest path at most twice the shortest, so $\height\le 2\lg(n+1)$
(\thmref{c06-rbheight}) -- the $\bigO{\log n}$ behind \code{std::set}/\code{map} that Chapter~7 borrowed
(\propref{c05-logn}). Insertion repairs a red-red violation by recolouring when the uncle is red (case (a)
of \figref{c06-rbfix}) and rotating when it is black, so edits cost fewer structural changes than AVL at
the price of a taller tree -- the reason write-heavy libraries pick red-black. You should \emph{use} it and
\emph{explain} it; to \emph{implement} a custom balanced node, use the treap (\secref{c06-treaps}).
\selftest{
  \item List the five invariants and say which two carry almost all the height argument.
  \item Prove $\height \le 2\mathit{bh}$ and $\mathit{bh} \le \lg(n+1)$ in two sentences.
  \item When does insertion recolour instead of rotate, and why is that cheaper?
  \item Give the worst-case height constants for AVL and red-black, and name the workload each suits.
  \item Why can \code{std::set} not return the $k$-th smallest in $\bigO{\log n}$, and what two structures fix it?
}
\end{recallbox}

% !TEX root = ../../tree_notes.tex
% =====================================================================
%  8.5  Splay Trees (intro; full block stays Ch25 c23)  sec:c06-splay
%  Exports: def:c06-splay  def:c06-zigzig  thm:c06-amortized
%           fig:c06-splayfig  ex:c06-splay-run   (practice: CF 343D)
% =====================================================================
\section{Splay Trees: Balance by Locality (Introduction)}\label{sec:c06-splay}

\begin{previouslybox}
\begin{enumerate}[nosep, leftmargin=1.6em]
  \item AVL guarantees $\bigO{\log n}$ by \emph{checking} heights; treaps by \emph{flipping coins};
        red-black by \emph{colouring}. Each stores extra per-node state. What if we stored \emph{none}?
        (the motivation for splay -- it keeps no balance factor, no priority, no colour.)
  \item Which single $\bigO{1}$ move lets us reshape a BST without disturbing the sorted inorder?
        (\thmref{c06-rotation}, \coderef{c06-rotcode}.)
  \item What does \emph{amortised} mean, and where have we used it before? (charging a rare expensive op
        against many cheap ones; \chref{c05}'s container discussion.)
\end{enumerate}
\end{previouslybox}

\begin{topicintro}
Every structure so far spends a \emph{worst-case} effort on every operation to keep \emph{all} keys
equally reachable. Splay trees, Sleator and Tarjan's invention, make a different bet: real access
patterns are \emph{skewed} -- a few hot keys dominate -- so why not spend rotations to drag whatever you
just touched to the root, letting hot data settle near the top and cold data sink? The rule is
astoundingly simple: after any access to node $x$, \term{splay} $x$ to the root with a bottom-up sequence
of rotations (\defref{c06-splay}), using only the zig / zig-zig / zig-zag moves of
\defref{c06-zigzig}. There is no auxiliary field and no coin flip. The payoff is a genuinely surprising
\emph{amortised} guarantee: $m$ operations cost $\bigO{(n+m)\log n}$, i.e.\ $\bigO{\log n}$ \emph{per
operation on average}, even though \emph{no single operation} is bounded better than linear
(\thmref{c06-amortized}). This section is the \emph{introduction}: the splay idea, the three moves, the
potential-function \emph{statement} of the amortised bound, and a hands-on splay of the chain $C$. The
full analysis -- the Access Lemma and its consequences (the static optimality, the working-set
property) -- and the link-cut trees that run splay on \emph{dynamic} trees live in Chapter~25
(\chref{c23}), which already expects exactly what you learn here.
\textbf{Prerequisites.} \secref{c06-rotations} (the rotation, the zig/zig-zag geometry); \secref{c06-avl}
for contrast. The amortised method is assumed lightly from \chref{c05}.

\textbf{After this section you can}
\begin{itemize}[nosep]
  \item perform a splay with the zig, zig-zig and zig-zag rotations and explain why zig-zig is two
        same-direction single rotations; \hfill\lv{3}--\lv{4}
  \item state the $k$-th-access working-set intuition and why splay exploits locality; \hfill\lv{3}
  \item state the amortised $\bigO{\log n}$ bound and name the potential ($\sum \log \mathit{size}$); \hfill\lv{4}--\lv{5}
  \item recognise a ``hot keys dominate / move-to-front'' workload where splay beats a rigid AVL. \hfill\lv{4}
\end{itemize}
\end{topicintro}

\begin{recognitionbox}[Pattern card -- self-adjusting by access pattern]
\reckey{a search tree whose query distribution is \emph{non-uniform and shifting} (hot keys recur, recent
keys recur), where a fixed balance policy wastes effort keeping cold keys shallow.}
\recuse{
  \item caches, dictionaries, ``move to front'' buffers -- recent / frequent items float to the root for
        free;
  \item you want a balanced BST with \emph{zero} per-node bookkeeping (no heights, colours or priorities);
  \item the entry point to \emph{dynamic trees}: link-cut trees (\chref{c23}) are splay trees whose nodes
        stand for preferred paths.
}
\recnotuse{
  \item you need a \emph{worst-case} $\bigO{\log n}$ per operation -- splay only gives it \emph{amortised};
        one access can be $\bigO{n}$ (use AVL/red-black, \secref{c06-avl}/\secref{c06-redblack});
  \item you need customisable subtree aggregates via clean concatenation -- the treap's
        \texttt{split}/\texttt{merge} (\secref{c06-treaps}) is far more convenient;
  \item plain ordered-set needs with no locality angle -- the library \code{set} already balances
        (\secref{c06-redblack}).
}
One-line test: \emph{do recent/frequent keys dominate the queries, and is an average-case (not per-call)
$\bigO{\log n}$ acceptable? Then let the tree reshape itself around the traffic -- splay.}
\end{recognitionbox}

\subsection{Intuition}

A treap is balanced \emph{on average over its random construction}; a splay tree is balanced
\emph{on average over the sequence of queries you actually ask}. The whole idea is one selfish rule:
whenever you reach a node, promote it to the root. If you search for $x$ and then immediately search for
$x$ again, the second search is now $\bigO{1}$ -- $x$ sits at the top. Frequently-used keys are constantly
re-promoted and thus stay near the root; unfrequently-used keys are almost never promoted and quietly
drift down. The tree ends up \emph{matched to your habits}, not to a worst-case adversary.

Promotion is not free -- a single splay walks all the way to the root, potentially $\bigO{n}$ rotations --
but the clever part is the \emph{order} of rotations. Splaying does not naively rotate $x$ up one level at
a time; it works two levels at a time (the zig-zig and zig-zag moves of \defref{c06-zigzig}) so that each
rotation roughly \emph{halves} the depth of $x$ \emph{and} flattens the whole path it traversed. That
``path halving'' is the geometric engine behind \thmref{c06-amortized}: an expensive splay leaves the tree so
much better shaped for the nodes on its path that the cost, spread over the accesses, averages to
$\bigO{\log n}$.

\begin{definition}{Splaying a node to the root}{c06-splay}
A \term{splay tree} is an ordinary binary search tree that stores \emph{no} extra fields. To
\term{splay} a node $x$: while $x$ is not the root, apply the moves of \defref{c06-zigzig} bottom-up
until $x$ becomes the root. \texttt{access(x)} --- a search that stops at $x$ (or where it would insert),
then splay; \texttt{insert(x)} --- normal BST insert then splay the new node; \texttt{delete(x)} --- splay
$x$, then join its left and right subtrees (splay the max of the left, attach the right). Every operation
\emph{ends} by splaying the touched key.
\end{definition}

\begin{definition}{Zig, zig-zig and zig-zag}{c06-zigzig}
Let $p$ be $x$'s parent and $g$ its grandparent. One splay step looks at the last two edges:
\begin{itemize}[nosep]
  \item \textbf{ZIG} ($p$ is the root): a single rotation about $p$ in the direction that lifts $x$. This
        is the final step and the whole move when $x$'s parent is the root.
  \item \textbf{ZIG-ZIG} ($x$ and $p$ on the \emph{same} side of $g$, a zig-zig line): rotate about $g$
        \emph{first}, then about $p$, both in the same direction. $x$ rises two levels and $p$ stays its
        parent.
  \item \textbf{ZIG-ZAG} ($x$ and $p$ on \emph{opposite} sides of $g$, a zig-zag bend): rotate about $p$
        in $x$'s direction, then about $g$ in $x$'s new direction -- $x$ rises two levels through the bend.
\end{itemize}
Repeat until $x$ is the root. Zig-zig's ``parent first'' ordering is the crucial design choice; the naive
``rotate $x$ up one at a time'' does \emph{not} achieve the amortised bound.
\end{definition}

\begin{figure}[H]
\centering
\begin{tikzpicture}[diagram, scale=0.8, transform shape]
  % ---- zig-zig (same side) ----
  \node[annot, font=\scriptsize\bfseries] at (1.1,0.7) {zig-zig (same side)};
  \node[troot] (zg) at (1.4,0)    {$g$};
  \node[tnode] (zp) at (0.5,-1.1) {$p$};
  \node[tleaf] (zx) at (-0.4,-2.2) {$x$};
  \draw (zg) -- (zp); \draw (zp) -- (zx);
  \draw[flow] (2.5,-1.1) -- (3.5,-1.1);
  \node[troot] (zx2) at (4.7,0)    {$x$};
  \node[tnode] (zp2) at (4.0,-1.1) {$p$};
  \node[tleaf] (zg2) at (3.3,-2.2) {$g$};
  \draw (zx2) -- (zp2); \draw (zp2) -- (zg2);
  \node[annot, font=\scriptsize, align=left] at (5.6,-1.4) {rotate $g$,\\ then $p$};
  % ---- zig-zag (bend) ----
  \node[annot, font=\scriptsize\bfseries] at (9.6,0.7) {zig-zag (bend)};
  \node[troot] (yg) at (9.4,0)    {$g$};
  \node[tnode] (yp) at (8.5,-1.1) {$p$};
  \node[tleaf] (yx) at (9.4,-2.2) {$x$};
  \draw (yg) -- (yp); \draw (yp) -- (yx);
  \draw[flow] (10.7,-1.1) -- (11.7,-1.1);
  \node[troot] (bx) at (12.9,-1.1) {$x$};
  \node[tleaf] (bp) at (12.0,-2.2) {$p$};
  \node[tleaf] (bg) at (13.8,-2.2) {$g$};
  \draw (bx) -- (bp); \draw (bx) -- (bg);
  \node[annot, font=\scriptsize, align=left] at (14.2,-1.0) {rotate $p$,\\ then $g$};
\end{tikzpicture}
\caption{The two two-level splay moves (the ZIG is just the single closing rotation). In \textbf{zig-zig}
the three nodes were collinear; rotating first about the grandparent, then the parent, lifts $x$ two
levels \emph{and} promotes $p$, leaving $x\to p\to g$ as a short left spine -- the whole local path
flattened. In \textbf{zig-zag} $x$ sat in the middle; one rotation each way about $p$ then $g$ brings the
bent $x$ to the top of the triple. The ``rotate the grandparent first'' order in zig-zig is precisely what
makes the amortised bound work.}\label{fig:c06-splayfig}
\end{figure}

\subsection{Why zig-zig, not rotate-$x$-up?}

Suppose we always just rotated $x$ one level at a time (skipping zig-zig). Consider repeatedly accessing
the deepest leaf of a long chain: each such access rotates $x$ to the root, but in doing so it drags
\emph{every other node on the path down one level}, so the second-deepest node becomes the new deepest.
The next access to \emph{it} reproduces the same linear path -- you get $\bigO{n}$ every time and the
amortisation collapses. Zig-zig fixes exactly this: by rotating the \emph{grandparent} first, the two ends
of every three-node segment both rise, so nodes \emph{not} accessed but sitting on the path also get pulled
up toward the root. This is the difference between the two schemes being $\bigO{n}$ versus $\bigO{\log n}$
amortised, and it is the entire art of \thmref{c06-amortized}.

\subsection{The Amortised Guarantee}

\begin{theorem}{Splay operations are $\bigO{\log n}$ amortised}{c06-amortized}
Starting from an $n$-node splay tree, any sequence of $m$ operations (access / insert / delete, each a
splay) costs $\bigO{(n+m)\log n}$ total, hence $\bigO{\log n}$ amortised per operation. No individual
operation is better than $\bigO{n}$ worst case; the bound is an average over the whole sequence.
\end{theorem}
\begin{proofbox}[The potential idea only -- the full proof is \chref{c23}]
Assign a \term{potential} $\Phi = \sum_{v} \log(\mathit{size}(v))$, the log of every subtree's node count,
where $\mathit{size}(v)$ is the number of nodes below (and including) $v$. The amortised cost of a splay is
its rotation count plus the change in $\Phi$. Splaying $x$ to the root grows the subtrees of the nodes on
the access path (they gain children as they are promoted), so $\Phi$ \emph{increases} for them but shrinks
everywhere the path thins; a careful accounting -- the \term{Access Lemma} -- bounds a single splay's
amortised cost by $3\big(\log \mathit{size}(\mathit{root}) - \log \mathit{size}(x)\big) + 1 =
\bigO{\log n}$. Summing over $m$ splays, the telescoping potential leaves $\bigO{(n+m)\log n}$ actual
rotations. We state the shape of the argument to make the guarantee credible; the Access Lemma and its
stronger consequences (dynamic finger, static optimality, working set) are proved in full in
\chref{c23}. \end{proofbox}
\begin{insightbox}
Note the tradeoff that \thmref{c06-amortized} makes explicit. AVL gives a \emph{worst-case}
$\bigO{\log n}$ that never lies to you but also never exploits locality (\thmref{c06-fixheight}); splay
gives an \emph{amortised} $\bigO{\log n}$ that can spike to $\bigO{n}$ on one call but in exchange is
\emph{asymptotically optimal for skewed access} -- a single hot key costs $\bigO{1}$ amortised once it is
saturated near the root. Pick the guarantee your problem can actually use.
\end{insightbox}

\subsection{Worked Example}

\begin{example}{Splaying the deep end of the chain $C$}{c06-splay-run}
Take the right-skew chain $C$: $1\to2\to3\to\cdots\to7$ (each the parent's right child) -- the shape AVL
was built to forbid and the treap avoids by luck. \texttt{access(7)} finds $7$ at depth $6$, then splays
it. All edges bend the same way (a pure zig-zig line), so each step rotates the grandparent then the
parent, and $7$ climbs two levels at a time: over the sequence $7$ reaches the root in $\bigO{6}$
rotations, and -- the point -- the intermediate nodes $2,3,4,5$ are pulled up too, so the tree is markedly
shallower for a follow-up access to \emph{them}. Now access the chain's \emph{oldest} hot key: if the
workload keeps asking $7$, it is at the root and every further query is $\bigO{1}$. Contrast with the same
chain in AVL (\exref{c06-avl-run}), where reaching the balanced twin required rotations \emph{at
insertion time}; splay did no work up front and instead spent it lazily, at access time, targeting the key
you actually asked for. Inorder stays $1,2,\dots,7$ throughout -- every move is a
\thmref{c06-rotation}-preserving rotation.
\end{example}

\begin{tracebox}[Zig-zig steps for $x$ on a right-right run $g\to p\to x$]
$x$ is the right child of $p$, which is the right child of $g$ (a zig-zig line). Step 1:
\texttt{rotateLeft(g)} -- $p$ rises above $g$, and $x$ is now $p$'s right child one level up. Step 2:
\texttt{rotateLeft(p)} -- $x$ rises above $p$; the three-node block now reads $x\to p\to g$ going down
$x$'s left side. $x$ gained two levels and $g,p$ each gained a shallower home. Repeat on the next grandparent
until $x$ is root. The mirror applies to a left-left run; a left-right or right-left run is the zig-zag
(rotate $p$ then $g$). \checkmark
\end{tracebox}

\subsection{Variations and Twists}

\begin{itemize}[nosep, leftmargin=1.4em]
  \item \textbf{Top-down splaying} -- perform the rotations during the \emph{descent} using a zig-zag-only
        join, avoiding the upward pointer walk and recursion; the standard efficient implementation. \hfill\lv{4}--\lv{5}
  \item \textbf{Splay as a move-to-front list} -- the same ``touch it, float it up'' logic underlies
        self-organising lists (MTF is $\bigO{\log n}$-competitive-ish by the same amortised lens). \hfill\lv{4}
  \item \textbf{Link-cut trees} -- replace the fixed parent/child with \emph{preferred paths} each stored in
        a splay tree, so you can \texttt{link}/\texttt{cut}/path-query a \emph{changing} tree in
        $\bigO{\log n}$ amortised. That is \chref{c23} in one sentence. \hfill\lv{5}
  \item \textbf{Partial / lazy splay} -- splay only to a fixed depth to bound worst-case spikes in
        real-time systems, sacrificing a little of the amortised tightness. \hfill\lv{5}
\end{itemize}

\subsection{Common Pitfalls}\label{sec:c06-splay-pitfalls}

\begin{warningbox}[Trap 1 -- rotate $x$ one level at a time (skip zig-zig)]
The amortised bound \thmref{c06-amortized} is false for naive one-level splaying (see the chain argument in
\secref{c06-splay} above). In a zig-zig you must rotate the \emph{grandparent} first; only then does a
splay flatten the whole path instead of just fetching $x$.
\end{warningbox}

\begin{warningbox}[Trap 2 -- calling the guarantee ``$\bigO{\log n}$ per operation'']
It is \emph{amortised}, not worst-case. A single access can cost $\bigO{n}$, so a hard real-time or
adversarially-timed setting may prefer AVL/red-black (\thmref{c06-fixheight}, \thmref{c06-rbheight}). Do
not cite splay when the problem needs a per-call bound.
\end{warningbox}

\begin{warningbox}[Trap 3 -- forgetting that every edit also splays]
Splay's shape is defined by \emph{what was last touched}. If you insert or delete without splaying the
affected node, you break the invariant the analysis assumes and lose the bound -- and, in practice,
degrade toward the chain $C$ you were trying to avoid.
\end{warningbox}

\subsection{Practice Problems}

\begin{stuckbox}
\begin{itemize}[nosep, leftmargin=1.4em]
  \item In a zig-zig, did you rotate the \emph{grandparent} before the parent? (Trap~1.)
  \item Can you state \thmref{c06-amortized}'s potential in one line? ($\Phi=\sum_v \log \mathit{size}(v)$.)
  \item Is your target a worst-case or an amortised bound? If worst-case, splay is the wrong tool.
\end{itemize}
\end{stuckbox}

\begin{practiceset}
  \practice{Water Tree}{Codeforces 343D}
    {https://codeforces.com/problemset/problem/343/D}{CF 2100 · L5}
    {A rooted tree, all vertices initially empty; three operations: pour water into a vertex \emph{and its
    whole subtree}, flush the path from a vertex up to the root, and query whether a vertex holds water.
    This is the chapter's ``keep an ordered index'' idea at full strength. Flatten the tree with an Euler
    tour so a subtree becomes the interval $[\tin(v), \tout(v)]$ (\chref{c24}); a clean solution keeps a
    \code{std::set} of currently-wet vertices keyed by $\tin$ -- pouring a subtree inserts $v$ and
    \texttt{lower\_bound}-and-erases every wet descendant (each vertex is erased at most as often as it is
    inserted, so the erasures amortise to $\bigO{\log n}$ each -- the balanced BST doing the heavy
    lifting), flushing records a ``dry at time $t$'' stamp on $v$'s ancestors, and a query compares the
    newest wet-fill against the newest dry on the path. The fully general route is Euler-tour $+$
    lazy segment tree with range-assign over timestamps (\chref{c13}, \chref{c15}); reach for splay /
    link-cut (\chref{c23}) when the \emph{tree itself} changes between operations. Solve it once with the
    \code{set}, and you will feel \thmref{c06-rbheight}'s $\bigO{\log n}$ index earning its keep.}
\end{practiceset}

\begin{recallbox}[10-second refresher -- 8.5 Splay Trees (Introduction)]
A splay tree stores no balance/priority/colour; after touching $x$ you \term{splay} it to the root with the
zig / zig-zig / zig-zag rotations (\defref{c06-splay}, \defref{c06-zigzig}) -- rotate the grandparent first
in a zig-zig, which is what flattens the path. Access patterns then self-balance: hot keys float up, cold
keys sink. The reward is amortised $\bigO{\log n}$ per operation (total $\bigO{(n+m)\log n}$) via the
potential $\Phi=\sum_v \log\mathit{size}(v)$ and the Access Lemma (\thmref{c06-amortized}) -- average, not
worst-case. Full analysis, top-down splaying, and link-cut dynamic trees are Chapter~25 (\chref{c23}).
\selftest{
  \item Draw one zig-zig and one zig-zag three-node configuration and the result of each.
  \item Why does rotating $x$ up one level at a time lose the amortised bound?
  \item State the splay potential and the total cost bound over $m$ operations.
  \item Give a workload where splay beats AVL, and one where its amortised-only guarantee disqualifies it.
  \item Name the chapter that turns splay into a structure for \emph{changing} trees.
}
\end{recallbox}

% !TEX root = ../../tree_notes.tex
% =====================================================================
%  Chapter 8 review -- the 15-section revision contract (chapter_end.txt).
%  Revision-only: tables, one-liners, pointers back via \thmref/\probref/
%  \secref/\coderef/\propref/\defref/\exref/\figref/\lemref/\corref.
%  No derivations.  Summit: LC Hard 493 Reverse Pairs
%  (a keyed order-statistic treap reads "how many seen keys c satisfy
%   nums[i] > 2c" on one guaranteed-O(log n) descent -- the balanced
%   upgrade of Chapter 7's fragile rank walk, with merge/Fenwick rivals).
% =====================================================================

\chaptersection{Chapter 8 Summary}

Five sections, one sentence: \emph{restore the shape without ever breaking the inorder line}. Chapter~7
spent the invariant $\text{left}<\text{node}<\text{right}$ to make a tree its own index, but every answer
cost $\bigO{h}$ and $h$ was not yours to control. This chapter takes the wheel. The rotation view (8.1)
proves a single local pointer surgery is the \emph{only} legal fix and that it leaves the sorted inorder
identical; the strict view (8.2) is AVL, which keeps $|\bal|\le1$ and repairs any insertion with $\le2$
rotations, buying $h\le1.44\,\lg(n+2)$; the randomised view (8.3) is the treap, which trades case-work for a
coin flip and, via \code{split}/\code{merge}, becomes a dynamic sequence; the pragmatic view (8.4) is
red-black, whose colour invariants repay Chapter~7's deferred $\bigO{\log n}$ promise for \code{std::set}; and
the self-adjusting view (8.5) is splay, which stores \emph{nothing} and earns $\bigO{\log n}$ only amortised.
Every entry below is a pointer, not a re-derivation.

\begin{center}
\renewcommand{\arraystretch}{1.25}
\begin{tabular}{@{}p{0.25\linewidth} p{0.47\linewidth} p{0.20\linewidth}@{}}
\toprule
\textbf{Result} & \textbf{Statement (with conditions)} & \textbf{Where} \\
\midrule
Height is the whole cost & Every search/insert/delete/rank is $\bigO{h}$; balance means forcing
  $h=\bigO{\log n}$ \emph{by construction}, whatever the arrival order. & \defref{c06-heightcost} \\
Rotation preserves inorder & A single left/right rotation re-weights a subtree while the sorted inorder
  line is \emph{identical} -- the only legal local move. & \thmref{c06-rotation} \\
Four violation cases & LL and RR need one rotation; LR and RL (a bend) need two -- child first, then
  anchor. & \defref{c06-rotcases} \\
AVL invariant & Every node obeys $|\bal(v)|\le1$ with $\bal=\height(L)-\height(R)$; a strict rule,
  repaired on every edit. & \defref{c06-bal} \\
AVL height bound & $N(h)=N(h-1)+N(h-2)+1$ (Fibonacci) forces $h\le1.44\,\lg(n+2)$. & \thmref{c06-avlbound} \\
One-insert repair & After one insertion only the lowest violator on the access path is fixed, in $\le2$
  rotations, then you stop early. & \thmref{c06-fixheight} \\
Treap expected height & BST on keys \emph{and} heap on random priorities is $\bigO{\log n}$ tall \emph{in
  expectation} over the coin flips. & \thmref{c06-exptheight} \\
Split / merge & \code{split} cuts a treap by key or position, \code{merge} joins two with all-left
  $<$ all-right; every keyed edit is two calls deep. & \defref{c06-splitmerge}, \lemref{c06-mergevalid} \\
Implicit keys & Replacing the key by left-subtree size turns a treap into a \emph{dynamic array}: insert,
  erase, reverse by position. & \propref{c06-implicit} \\
Red-black invariants & Five colour rules; no two reds stack and all paths share a black-height -- lazier
  than AVL, cheap to restore. & \defref{c06-rbcolor} \\
Red-black height & Longest path $\le2\times$ shortest $\Rightarrow h\le2\,\lg(n+1)$: the exact guarantee
  \propref{c05-logn} promised for \code{set}/\code{map}. & \thmref{c06-rbheight} \\
Splay amortised & No stored balance; every access rotates to the root. $m$ operations cost
  $\bigO{(n+m)\log n}$ -- $\bigO{\log n}$ \emph{amortised}, not worst-case. & \thmref{c06-amortized} \\
\bottomrule
\end{tabular}
\end{center}

\medskip
{\small\textbf{Algorithms from this chapter -- rotate, then rebalance by rule, by chance, by colour, or
by habit.}}
\begin{center}
\small
\renewcommand{\arraystretch}{1.2}
\setlength{\tabcolsep}{3pt}
\begin{tabular}{@{}p{0.19\linewidth} l l p{0.29\linewidth} p{0.155\linewidth}@{}}
\toprule
\textbf{Algorithm} & \textbf{Time} & \textbf{Memory} & \textbf{When to use (deciding phrase)} & \textbf{Code} \\
\midrule
Left / right rotation
  & $\bigO{1}$ & $\bigO{1}$
  & ``fix a height violation without touching the sorted line''
  & \coderef{c06-rotcode} \\
AVL insert (balance on return)
  & $\bigO{\log n}$ & $\bigO{\log n}$
  & ``keep a hand-rolled BST shortest; read-heavy''
  & \coderef{c06-avlins} \\
Treap insert/erase via split+merge
  & $\bigO{\log n}$ exp. & $\bigO{n}$
  & ``a balanced tree you can actually code cold''
  & \coderef{c06-split} \\
Implicit-sequence ops
  & $\bigO{\log n}$ exp. & $\bigO{n}$
  & ``insert/reverse/range-sum a sequence by position''
  & \coderef{c06-imp} \\
Order-statistic rank (Summit)
  & $\bigO{\log n}$ exp. & $\bigO{n}$
  & ``count keys meeting a threshold, online, worst-case fast''
  & \coderef{c06-summit} \\
\bottomrule
\end{tabular}
\end{center}

\chaptersection{Master Decision Flowchart}

Post-study ladder: walk top to bottom, stop at the first branch that matches. It decides \emph{which
balancing philosophy} the situation asks for -- a follower never hand-rolls a treap when the library
\code{set} already balances, nor trusts a plain BST when keys arrive sorted.

\begin{lstlisting}[language={}, basicstyle=\ttfamily\small, numbers=none, frame=single]
  You need an ORDERED set/map (or a dynamic sequence). Which balancer?
  |
  |-- keys can arrive in sorted / adversarial order and you MUST keep depth ~log n
  |     +-- never trust a plain BST: its height is the enemy        [8.1, chain C]
  |     +-- a rotation is the only inorder-preserving local fix      [8.1]
  |
  |-- you are implementing the tree yourself, read-heavy, want the shortest tree
  |     +-- AVL: |BF|<=1 per node, <=2 rotations per insert          [8.2, GFG AVL]
  |     +-- track balance factors up the path; fix the LOWEST one    [8.2]
  |
  |-- you want the simplest correct balanced tree, or you must cut/paste ranges
  |     +-- treap: BST on keys + heap on random priority             [8.3, LC 307]
  |     +-- split/merge are the whole toolkit; no rotation cases     [8.3]
  |
  |-- it is a SEQUENCE (positions, not values): insert at i, reverse a range
  |     +-- implicit-key treap; key = left-subtree size              [8.3]
  |
  |-- you only need an ordered set/map and can use the library
  |     +-- std::set / std::map: they ARE red-black, O(log n) proven [8.4]
  |     +-- but no fast rank -- kth-smallest needs an order-stat tree[8.4, CF 1354D]
  |
  |-- hot keys dominate, or you want zero stored bookkeeping
  |     +-- splay: rotate what you touch to the root, amortised      [8.5]
  |     +-- dynamic trees / link-cut -> Ch25 [c23]; Cartesian/monotonic -> Ch9 [c26]
  |
  +-- the query is a rank / count-under-threshold, online, worst-case O(log n)
      +-- augment a balanced tree with subtree size (Summit 493), or
        let an array index do it: Fenwick [c14] / segment tree [c13] / merge split.
\end{lstlisting}

\chaptersection{Common Mistakes Checklist}

\begin{itemize}[nosep, leftmargin=1.4em]
  \item \textbf{Rotating the wrong way or forgetting a subtree.} A left rotation promotes the \emph{right}
        child and re-parents that child's left subtree as the old root's right; mirror for right. Dropping
        the moved subtree $B$ silently deletes keys (\thmref{c06-rotation}, \coderef{c06-rotcode}).
  \item \textbf{Not updating heights bottom-up after a rotation.} Recompute the \emph{old} root first, then
        its new parent -- in that order -- or a stale height mislabels the very balance factor you are about
        to read (\coderef{c06-avlrot}).
  \item \textbf{Treating LR/RL like LL/RR.} When the heavy child leans the \emph{opposite} way, one rotation
        only converts the bend into a line; you must rotate the child first. Check the grandchild's balance
        factor (\defref{c06-rotcases}, \figref{c06-zigzag}).
  \item \textbf{Rotating at the highest violator on insertion.} Only the \emph{lowest} node with
        $|\bal|=2$ needs a turn; after it, the subtree height is restored and every ancestor is already fine
        -- fix once and stop (\thmref{c06-fixheight}).
  \item \textbf{Breaking \code{merge}'s precondition.} \code{merge(a,b)} is only valid when every key in $a$
        precedes every key in $b$; call it on two arbitrary treaps and the result is not a BST
        (\lemref{c06-mergevalid}).
  \item \textbf{Splitting an implicit treap by value.} With implicit keys the number you pass to \code{split}
        is a \emph{position} count compared against $\mathit{size}(\text{left})$, not a stored key
        (\propref{c06-implicit}).
  \item \textbf{Assuming \code{std::set} answers ``$k$th smallest'' fast.} A red-black \code{set} exposes
        \code{lower\_bound} in $\bigO{\log n}$ but rank in $\bigO{n}$; you must augment or switch to a
        Fenwick/order-stat tree (\thmref{c06-rbheight}, practice CF 1354D).
  \item \textbf{Splaying by rotating $x$ up one level at a time.} The naive single-step splay is
        $\bigO{n}$ \emph{every} access; the amortised bound needs zig-zig to rotate the \emph{grandparent}
        first (\defref{c06-zigzig}, \thmref{c06-amortized}).
\end{itemize}

\chaptersection{WA Clinic}

Verdicts are symptoms; the causes are the short list this chapter drilled. Diagnose top-down -- every
leaf names the training problem.

\begin{figure}[H]
\centering
\begin{tikzpicture}[diagram,
    v/.style={boxnode, fill=primaryblue!12, draw=primaryblue, text width=22mm, align=center},
    b/.style={boxnode, fill=difforange!12, draw=difforange, text width=12mm},
    l/.style={boxnode, fill=successgreen!10, draw=successgreen, text width=58mm,
              font=\scriptsize, align=left}]
  \node[v] (verdict) {\textbf{Verdict on the judge}};
  \node[b, below left=8mm and 6mm of verdict] (wa) {\textbf{WA}};
  \node[b, below right=1mm of verdict] (tle) {\textbf{TLE}};
  \node[b, below left=8mm and 1mm of verdict, xshift=30mm] (re) {\textbf{RE}};
  \node[b, below right=8mm and 6mm of verdict] (mle) {\textbf{MLE}};
  \node[l, below right=1mm and 2mm of wa.south] (wa1)
    {a rotation drops the moved subtree $B$ and loses keys (\thmref{c06-rotation}); LR done as a single
     turn (\defref{c06-rotcases}); \code{merge} called on treaps that are not all-left-before-all-right
     (\lemref{c06-mergevalid}); implicit split compares a value instead of $\mathit{size}(\text{left})$
     (\propref{c06-implicit})};
  \node[l, below right=1mm and 2mm of tle.south] (tle1)
    {a plain BST fed sorted keys -- the chain $C$ again, $\bigO{n}$ per descent (\defref{c06-heightcost},
     \figref{c06-chain}); the naive rotate-one-level splay, $\bigO{n}$ per access (\thmref{c06-amortized});
     the Summit on a fragile unbalanced rank tree is $\bigO{n^2}$ (\probref{c06-summit})};
  \node[l, below right=1mm and 2mm of re.south] (re1)
    {dereferencing a null child in a rotation before the base case; stale height reads after forgetting
     \code{upd} bottom-up (\coderef{c06-avlrot}); 32-bit overflow in a $2\cdot$-threshold compare
     (\probref{c06-summit})};
  \node[l, below right=1mm and 2mm of mle.south] (mle1)
    {recursion depth on a treap whose priorities were a bad draw -- fixed by a fresh random seed
     (\thmref{c06-exptheight}); leaking split/merge temporaries in a range-reverse loop (\coderef{c06-imp})};
  \draw[flow] (verdict) -- (wa); \draw[flow] (verdict) -- (tle);
  \draw[flow] (verdict) -- (re); \draw[flow] (verdict) -- (mle);
  \draw[flow] (wa) -- (wa1); \draw[flow] (tle) -- (tle1);
  \draw[flow] (re) -- (re1); \draw[flow] (mle) -- (mle1);
\end{tikzpicture}
\caption*{Chapter~8 verdict tree -- only the failures this chapter's problems can produce.}
\end{figure}

\medskip
\noindent Backup table (same content, terse):
\begin{center}
\small
\renewcommand{\arraystretch}{1.2}
\begin{tabular}{@{}llp{0.42\linewidth}l@{}}
\toprule
\textbf{Verdict} & \textbf{Usual cause} & \textbf{One-line fix} & \textbf{Trained in} \\
\midrule
WA  & rotation drops subtree $B$   & re-parent $B$ before re-attaching  & \thmref{c06-rotation} \\
WA  & LR treated as LL/RR          & rotate the child first, then anchor & \defref{c06-rotcases} \\
WA  & bad \code{merge} inputs      & only merge when all $a$ keys precede $b$ & \lemref{c06-mergevalid} \\
TLE & sorted feed, plain BST       & balance by construction (AVL/treap) & \figref{c06-chain} \\
TLE & Summit on fragile rank tree  & balance the tree, or Fenwick/merge & \probref{c06-summit} \\
RE  & stale height after rotation  & \code{upd} old root then new parent & \coderef{c06-avlrot} \\
\bottomrule
\end{tabular}
\end{center}

\chaptersection{Quick Reference Card}

\begin{quickref}[Chapter 8 at a glance]
\textbf{One move, four policies.} The rotation is the only inorder-preserving local repair
(\thmref{c06-rotation}); AVL, treap, red-black and splay are simply \emph{when to rotate and how much to
care}. Keeping $h=\bigO{\log n}$ is the whole point -- every Chapter-7 descent is then guaranteed fast
(\defref{c06-heightcost}).

\medskip
\textbf{Two repair orders.} \emph{AVL}: walk heights up, fix the lowest $|\bal|=2$ in $\le2$ rotations,
stop (\thmref{c06-fixheight}). \emph{Red-black}: recolour when the uncle is red, rotate when it is black --
mostly flips, at most two turns (\defref{c06-rbcolor}, \figref{c06-rbfix}).

\medskip
\textbf{Three running objects.} The chain $C$ (keys $1,\dots,7$ ascending, height $6$) is the villain; its
balanced twin has height $2$ (\figref{c06-chain}). The zig-zag $Z$ ($30,10,40,20,25$) forces the LR double
rotation (\figref{c06-zigzag}). The same keys as a treap under random priorities are short \emph{in
expectation} (\thmref{c06-exptheight}).

\medskip
\textbf{The payoff.} A controlled height makes augmentation \emph{safe}: a size-augmented keyed treap reads
an order statistic in guaranteed $\bigO{\log n}$, the honest upgrade of Chapter~7's fragile rank walk -- the
Summit (\probref{c06-summit}) is that upgrade at LC Hard scale.
\end{quickref}

\begin{lstlisting}[caption={Chapter 8 code kit: the rotation, AVL balance and insert, treap merge and split, an implicit split, and an order-statistic rank}, label={lst:c06-kit}]
#include <bits/stdc++.h>
using namespace std;
struct Node { int key, pri, sz, ht; Node *l, *r;
  Node(int k) : key(k), pri(rand()), sz(1), ht(0), l(nullptr), r(nullptr) {} };

// (1) The atomic move -- left about x, right about y.                      [8.1]
Node* rotateLeft(Node* x){ Node* y = x->r; x->r = y->l; y->l = x; upd(x); upd(y); return y; }
Node* rotateRight(Node* y){ Node* x = y->l; y->l = x->r; x->r = y; upd(y); upd(x); return x; }

// (2) AVL balance -- four cases, at most two rotations.                    [8.2]
int H(Node* n){ return n ? n->ht : -1; }
void upd(Node* n){ n->ht = 1 + max(H(n->l), H(n->r)); }
int BF(Node* n){ return H(n->l) - H(n->r); }
Node* balance(Node* n){ upd(n);
  if (BF(n)==2){ if (BF(n->l)<0) n->l=rotateLeft(n->l); return rotateRight(n); }
  if (BF(n)==-2){ if (BF(n->r)>0) n->r=rotateRight(n->r); return rotateLeft(n); }
  return n; }
// (3) AVL insert -- Ch7 plant, then balance on the way back up.            [8.2]
Node* insert(Node* n, int k){ if(!n) return new Node(k);
  if(k<n->key) n->l=insert(n->l,k); else if(k>n->key) n->r=insert(n->r,k);
  return balance(n); }

// (4) Treap merge / split -- rotation-free rebalancing.                    [8.3]
int SZ(Node* t){ return t ? t->sz : 0; }
Node* merge(Node* a, Node* b){ if(!a) return b; if(!b) return a;
  if(a->pri > b->pri){ a->r=merge(a->r,b); t_upd(a); return a; }
  else               { b->l=merge(a,b->l); t_upd(b); return b; } }
void splitPos(Node* t, int cnt, Node*& a, Node*& b){    // a = first cnt positions
  if(!t){ a=b=nullptr; return; }
  if(SZ(t->l)>=cnt){ splitPos(t->l,cnt,a,t->l); b=t; t_upd(b); }
  else             { splitPos(t->r,cnt-SZ(t->l)-1,t->r,b); a=t; t_upd(a); } }
// (5) Order-statistic rank on a keyed treap -- count keys c with 2c < x.   [Summit]
int countSmallerHalf(Node* t, long long x){ int res=0;
  while(t){ if(2LL*t->key < x){ res += 1 + SZ(t->l); t=t->r; } else t=t->l; } return res; }
\end{lstlisting}

\chaptersection{Write-From-Memory Drills}

Five signature skeletons, each lifted from a listing already in the chapter. Cover the key line,
reproduce it from memory, then check the answer one line below.

\begin{lstlisting}[caption={Drill 1 -- the left rotation (from \coderef{c06-rotcode})}]
Node* y = x->right;
x->right = /* FILL: y->left */;   // subtree B moves under x
y->left  = x;                      // x becomes y's left child
upd(x); /* FILL: upd(y) */;        // old root first, then its new parent
return y;                          // y is the new subtree root
\end{lstlisting}
\noindent{\scriptsize Key: the rotation re-parents exactly one subtree ($B$) and must refresh the
\emph{changed} node's height before its parent's -- the inorder line never moves (\thmref{c06-rotation}).}

\begin{lstlisting}[caption={Drill 2 -- AVL four-case balance (from \coderef{c06-avlrot})}]
Node* balance(Node* n){ upd(n);
  if (BF(n) == 2) { if (/* FILL: BF(n->left) < 0 */) n->left = rotateLeft(n->left);
                    return rotateRight(n); }
  if (BF(n) == -2){ if (BF(n->right) > 0) n->right = /* FILL: rotateRight(n->right) */;
                    return rotateLeft(n); }
  return n; }
\end{lstlisting}
\noindent{\scriptsize Key: heavy child leaning the \emph{other} way is the LR/RL bend -- rotate the child
first to make a line, then the anchor (\defref{c06-rotcases}, \thmref{c06-fixheight}).}

\begin{lstlisting}[caption={Drill 3 -- AVL insert (from \coderef{c06-avlins})}]
if (n == nullptr) return new Node(key);         // plant, exactly as Ch 7
if (key < n->key)      n->left  = insert(n->left,  key);
else if (key > n->key) n->right = insert(n->right, key);
else                   return n;                // distinct keys: no re-insert
return /* FILL: balance(n) */;                  // fix THIS node on the return path
\end{lstlisting}
\noindent{\scriptsize Key: the only new line is the \code{balance} call on the way back up; the descent and
the plant are Chapter~7 verbatim (\secref{c06-avl}).}

\begin{lstlisting}[caption={Drill 4 -- treap merge (from \coderef{c06-treap})}]
// precondition: every key in a precedes every key in b
if (!a) return b;  if (!b) return a;
if (a->pri > b->pri) { a->r = /* FILL: merge(a->r, b) */; upd(a); return a; }
else                 { b->l = merge(a, b->l);           upd(b); return b; }
\end{lstlisting}
\noindent{\scriptsize Key: the higher-priority node is the root; recurse down its right (or the other's
left) spine. No rotation ever -- the priorities \emph{are} the shape (\thmref{c06-exptheight}).}

\begin{lstlisting}[caption={Drill 5 -- implicit split by position (from \coderef{c06-split})}]
void splitPos(Node* t, int cnt, Node*& a, Node*& b){   // a = first cnt positions
    if (!t){ a = b = nullptr; return; }
    if (SZ(t->l) >= cnt) { splitPos(t->l, cnt, a, t->l); /* FILL: b = t */; upd(b); }
    else { splitPos(t->r, cnt - SZ(t->l) - 1, t->r, b); a = t; upd(a); }
}
\end{lstlisting}
\noindent{\scriptsize Key: compare the \emph{position} count to $\mathit{size}(\text{left})$, not to any
key -- that single swap makes a treap a dynamic array (\propref{c06-implicit}).}

\chaptersection{Complete Problem Checklist}

Every problem this chapter offers -- 3 fully solved in-section, 5 practice, one Summit, five mixed.
Tick \emph{only} after an unaided solve; a hinted attempt counts as half. Cite by name and number --
URLs appear once, in the solved and practice entries.

\begin{center}
\footnotesize
\renewcommand{\arraystretch}{1.25}
\begin{tabular}{@{}c p{0.19\linewidth} l p{0.20\linewidth} p{0.24\linewidth} c@{}}
\toprule
\textbf{\#} & \textbf{Problem} & \textbf{Source} & \textbf{Key technique} & \textbf{What to look for} & \textbf{Done?} \\
\midrule
\multicolumn{6}{l}{\textbf{Section 8.1 -- why balance; the rotation}} \\
W8.1.1 & Tree Construction & CF 675D & \lv{4} set \code{lower\_bound} parent & pick the \emph{later} of succ/pred; $-1$ root & $\square$ \\
P8.1.2 & Balance a BST & LC 1382 & \lv{3} inorder $\to$ rebuild from middle & sorted line becomes a tree by halving & $\square$ \\
\midrule
\multicolumn{6}{l}{\textbf{Section 8.2 -- AVL trees}} \\
W8.2.1 & AVL Tree Insertion & GFG & \lv{4} balance on return + preorder & update heights bottom-up; fix lowest violator & $\square$ \\
P8.2.2 & AVL Tree Deletion & GFG & \lv{4} three-case delete + rebalance up & deletion may need rotations at \emph{several} levels & $\square$ \\
\midrule
\multicolumn{6}{l}{\textbf{Section 8.3 -- treaps, split/merge, implicit keys}} \\
W8.3.1 & Range Sum Query -- Mutable & LC 307 & \lv{3} implicit sum-augmented treap & split to width 1 to write; read $\mathit{sum}$ of the middle & $\square$ \\
P8.3.2 & Subarray Sum Queries & CSES 1190 & \lv{3} point-update range-max & treap or segment tree; max needs the right merge rule & $\square$ \\
\midrule
\multicolumn{6}{l}{\textbf{Sections 8.4--8.5 -- red-black (use) and splay (intro)}} \\
P8.4.1 & Multiset & CF 1354D & \lv{3} order statistic by value & $k$th smallest needs a count, not a \code{set} descent & $\square$ \\
P8.5.2 & Water Tree & CF 343D & \lv{5} subtree fill / path flush / query & splay or Euler-tour $+$ segtree; path-to-root is the hard part & $\square$ \\
\midrule
\multicolumn{6}{l}{\textbf{Summit (chapter end)}} \\
S8 & Reverse Pairs & LC 493 & \lv{5} keyed treap rank $2c<x$ & right-to-left, count seen under threshold; mind $2\cdot$ overflow & $\square$ \\
\bottomrule
\end{tabular}
\end{center}

\chaptersection{Recommended Study Order}

Five focused days. Day 1 fixes the rotation as the one legal move; day 2 the AVL insert cold; day 3 the
treap and its split/merge engine; day 4 the two \emph{reading} sections (red-black guarantee, splay
amortisation); day 5 the Summit. Solve in this order even where it differs from section order -- each
day's tool is already in hand the day before.

\begin{center}
\small
\renewcommand{\arraystretch}{1.25}
\begin{tabular}{@{}llp{0.55\linewidth}@{}}
\toprule
\textbf{Day} & \textbf{Focus} & \textbf{Tasks} \\
\midrule
1 & Read + do 8.1 & rotate the head triple of $C$ by hand (\exref{c06-rot-run}); resolve $Z$'s LR case; CF 675D \\
2 & Read + do 8.2 & insert $1,\dots,7$ eagerly on $C$, then $Z$ (\exref{c06-avl-run}); GFG AVL insertion; write \code{balance} cold \\
3 & Read + do 8.3 & draw the treap on $C$ with lucky priorities, split/merge $Z$ (\exref{c06-treap-run}); LC 307; CSES 1190 \\
4 & Read + do 8.4--8.5 & colour the chain $C$ red-black (\exref{c06-rb-run}); splay the deep end of $C$ (\exref{c06-splay-run}); CF 1354D \\
5 & Summit + revise & LC 493 with the keyed-treap rank; walk the decision flowchart on five questions \\
\bottomrule
\end{tabular}
\end{center}

\chaptersection{Difficulty Progression}

Tiers are competency gates -- each row names the \emph{skill} the tier unlocks, not just a rating
bucket.

\begin{center}
\small
\renewcommand{\arraystretch}{1.3}
\begin{tabular}{@{}lp{0.36\linewidth}p{0.34\linewidth}@{}}
\toprule
\textbf{Tier} & \textbf{Problems} & \textbf{What passing this tier proves} \\
\midrule
$\lv{2}$ \ $\bigstar\bigstar$
  & rotate $C$ and $Z$ by hand
  & you see the rotation as the only inorder-preserving move and never drop subtree $B$ \\
$\lv{3}$ \ $\bigstar\bigstar\bigstar$
  & LC 1382, LC 307, CSES 1190
  & rebuild balanced from the sorted line; drive a sequence with \code{split}/\code{merge} \\
$\lv{4}$ \ $\bigstar\bigstar\bigstar\bigstar$
  & CF 675D, GFG AVL insert, GFG AVL delete, CF 1354D
  & code AVL cold; reason about the \code{set} guarantee and its missing rank \\
$\lv{5}$ \ $\bigstar\bigstar\bigstar\bigstar\bigstar$
  & CF 343D, LC 493 (Summit)
  & exploit a controlled height for online rank; see the fragile BST become quadratic \\
\bottomrule
\end{tabular}
\end{center}

\chaptersection{Reduction Gallery}

Balance is a disguise; each line collapses a family onto the one idea -- \emph{keep the descent short by
construction, then spend it}. No new theory -- just where the rotation is already hiding.

\begin{itemize}[nosep, leftmargin=1.4em]
  \item ``insert keys that may arrive sorted; the tree must stay fast'' (chain $C$) $\longrightarrow$
        a balancing policy, not a smarter search (\defref{c06-heightcost}, \figref{c06-chain}).
  \item ``fix a height violation'' $\longrightarrow$ one rotation on the lowest violator, child-first if it
        bends (\thmref{c06-rotation}, \defref{c06-rotcases}).
  \item ``implement an AVL tree'' (GFG) $\longrightarrow$ Chapter 7's insert with a \code{balance} call on
        the return path (\coderef{c06-avlins}, \probref{c06-avl}).
  \item ``rebuild a \emph{balanced} BST from a BST'' (LC 1382) $\longrightarrow$ its inorder, then halve --
        the sorted line becomes the tree (\secref{c06-rotations}).
  \item ``mutate an element, re-sum a range, online'' (LC 307) $\longrightarrow$ an implicit, sum-augmented
        treap: split out the position, rewrite, merge back (\propref{c06-implicit}, \probref{c06-rpmq}).
  \item ``point update, range max'' (CSES 1190) $\longrightarrow$ the same split/merge engine with a
        different aggregate -- or a segment tree (\chref{c13}).
  \item ``the $k$-th smallest, online, worst-case fast'' (CF 1354D) $\longrightarrow$ a size-\emph{count}
        over values; a red-black \code{set} alone will not do it (\thmref{c06-rbheight}).
  \item ``fill a subtree, flush a path to the root'' (CF 343D) $\longrightarrow$ splay the node up, or
        Euler-tour the tree and range-update a segtree (\secref{c06-splay}, \chref{c24}).
  \item ``count pairs $i<j$ with $a_i>2a_j$'' (LC 493) $\longrightarrow$ a right-to-left online rank on a
        balanced keyed treap -- the guaranteed form of Chapter 7's fragile walk (\probref{c05-summit},
        \probref{c06-summit}).
\end{itemize}

\chaptersection{Mixed Problem Set}

These refuse to name their balancer; decide \emph{which structure keeps the descent short enough} before
opening the code. A hint tells you the shape of the answer, never the tool.

\begin{mixedset}
  \practice{Second Order Statistics}{Codeforces 22A}
    {https://codeforces.com/problemset/problem/22/A}{CF 800 · L2}
    {Print the second \emph{distinct} smallest key. A \code{std::set} sorts and de-duplicates on the way
    in, so the answer is the second element -- one \code{begin()} step. The lesson this chapter taught you
    to distrust: a \emph{hand}-rolled BST on sorted-feeling input can rot into a chain, but the red-black
    \code{set} is short by construction (\thmref{c06-rbheight}). Two lines, no rotation.}
  \practice{Sliding Window Median}{CSES 1076}
    {https://cses.fi/problemset/task/1076}{CSES · L4}
    {Median of every length-$k$ window, sliding by one. Two \code{multiset}s held as a lower/upper pair
    with a size invariant is the classic balance -- the same ``keep two halves within one of each other''
    discipline as AVL's $|\bal|\le1$, but on containers (\secref{c06-avl}). Insert, evict, and \emph{rebalance
    the split} so the sizes differ by at most one; the median is the top of the lower half.}
  \practice{Sliding Cost}{CSES 1077}
    {https://cses.fi/problemset/task/1077}{CSES · L4}
    {Total cost to move every window element to the median, summed over windows. Get the median as in
    CSES 1076, then need the \emph{sum} of each half too -- a size-and-sum-augmented structure, exactly the
    augmentation that makes an implicit treap answer range sums (\propref{c06-implicit}). Or keep an order
    statistic over values with a Fenwick pair (\chref{c14}): rank \emph{and} prefix sum, both logarithmic.}
  \practice{Find Median from Data Stream}{LeetCode 295}
    {https://leetcode.com/problems/find-median-from-data-stream/}{LC Hard L4}
    {Online: add a number, report the running median. Two heaps (a max-heap low, min-heap high) with a
    size invariant is the intended $\bigO{\log n}$ answer and needs no balance -- but notice it answers only
    \emph{median}, not arbitrary rank. A keyed treap with subtree sizes answers \emph{any} $k$th smallest
    online; that guaranteed-height augmentation is the Chapter-8 generalisation of this classic
    (\probref{c06-rpmq} is the same engine with a different aggregate).}
  \practice{Sequentially Ordinal Rank Tracker}{LeetCode 2102}
    {https://leetcode.com/problems/sequentially-ordinal-rank-tracker/}{LC Hard L5}
    {Add scenic spots, then ask for the $k$-th by a \emph{composite} score. This is a pure order-statistic
    query over a changing set -- a red-black \code{set} cannot answer rank fast (\thmref{c06-rbheight}), so
    you need a size-augmented balanced tree keyed by the score pair, or a Fenwick over compressed scores
    (\chref{c14}). Ordering spots is the treap's heap-on-priority idea wearing a comparator; read the
    $k$th with the same left-size descent as the Summit (\probref{c06-summit}).}
  \practice{Count of Range Sum}{LeetCode 327}
    {https://leetcode.com/problems/count-of-range-sum/}{LC Hard L4}
    {A subarray sum is a difference of two prefix sums, so count pairs inside the
    window as you walk: keep earlier prefixes in an order-statistic structure and
    query the rank of $[p-\mathit{upper},\,p-\mathit{lower}]$ twice per prefix --
    the same augmentation the Summit reads in one descent (\probref{c06-summit}).}
  \practice{My Calendar III}{LeetCode 732}
    {https://leetcode.com/problems/my-calendar-iii/}{LC Hard L4}
    {Book every event and report the current maximum $k$-overlap: stamp $+1$ and
    $-1$ at the two endpoints of an ordered \code{map} keyed by time and take the
    best running prefix sum -- the \code{map} does the ordering for you
    (\secref{c06-avl}), no rotations in sight.}
  \practice{Longest Continuous Subarray With Absolute Diff Less Than or Equal to Limit}{LeetCode 1438}
    {https://leetcode.com/problems/longest-continuous-subarray-with-absolute-diff-less-than-or-equal-to-limit/}{LC Medium L3}
    {A window is legal while its max minus its min fits the limit: two monotone
    deques ride along holding exactly those extremes -- the rank-free cousin of
    \probref{c06-rpmq}'s ordered container, shrinking from the left by what the
    deque heads say.}
  \practice{Count Integers in Intervals}{LeetCode 2276}
    {https://leetcode.com/problems/count-integers-in-intervals/}{LC Hard L5}
    {Add ranges, then ask how many covered points a query range holds: keep the
    disjoint intervals in an ordered structure \emph{augmented} with the covered
    length of each node's span -- \propref{c06-implicit}'s store-more-than-keys
    move aimed at ranges instead of ranks.}
  \practice{Number of Flowers in Full Bloom}{LeetCode 2251}
    {https://leetcode.com/problems/number-of-flowers-in-full-bloom/}{LC Hard L4}
    {How many flowers are alive at each asked time: sort the start and end events
    and walk the queries in key order with two cursors and a running count --
    before building anything, ask which ordered structure you are simulating;
    online, that question is \probref{c06-rpmq}'s.}
\end{mixedset}

\chaptersection{The Summit -- LC 493 Reverse Pairs}

\begin{notebox}[Why this is the chapter's top rope]
Chapter~7's Summit (LC 315) asked, for every element, ``how many keys already in my tree are strictly
smaller?'' -- a \emph{rank} read off a size-augmented BST in one $\bigO{h}$ descent, correct but
$\bigO{n^2}$ on sorted input, the cliffhanger this whole chapter exists to retire. Reverse Pairs is the
same online rank with a twist: count $i<j$ with $\mathit{nums}[i]>2\cdot\mathit{nums}[j]$. Process
right-to-left and it is again ``how many already-seen keys $c$ satisfy $2c<\mathit{nums}[i]$,'' a rank
query with a $2\cdot$-threshold -- so you still descend a size-augmented tree, but now \emph{the height must
be yours}. On a plain BST a sorted feed is a chain and the $\bigO{nh}$ collapses to $\bigO{n^2}$; on a keyed
treap the height is $\bigO{\log n}$ \emph{in expectation by construction}, so the identical algorithm becomes
$\bigO{n\log n}$. That is the chapter's thesis in one problem: \emph{the augmentation is not the win -- the
controlled height is}. Merge-sort and a Fenwick tree reach the same $\bigO{n\log n}$ by refusing to let a
descent's length depend on input order at all (\chref{c01}, \chref{c14}).
\end{notebox}

\solvedproblem{c06-summit}{Reverse Pairs}{LeetCode 493}
  {https://leetcode.com/problems/reverse-pairs/}{LC Hard L5}
  {right-to-left online rank, keyed order-statistic treap, count seen values under a 2x threshold, height control vs the fragile BST}

\begin{problemstatement}[Reverse Pairs]
Given an integer array \code{nums}, return the number of \term{reverse pairs}: pairs $(i,j)$ with
$i<j$ and $\mathit{nums}[i] > 2\cdot\mathit{nums}[j]$. $1\le n\le 5\cdot 10^{4}$ and
$-2^{31}\le \mathit{nums}[i]\le 2^{31}-1$.
\end{problemstatement}

\begin{attemptgate}[Attempt it yourself first -- \lv{5}, give yourself 45 minutes]
A fair first attempt is the $\bigO{n^2}$ double loop (for each $i$, count $j>i$ with
$\mathit{nums}[i]>2\mathit{nums}[j]$) -- correct, times out. The upgrade is to keep the ``seen to the
right'' in a \emph{balanced, size-augmented} structure and ask its rank under the doubled threshold.
\peek{Read the array right-to-left so everything you have already inserted lies to the \emph{right} of the
current index. If those values lived in a size-augmented tree, what exactly must each node let you count
when you want ``how many stored keys $c$ satisfy $2c<\mathit{nums}[i]$?''}
\nudge{At a node with key $c$: if $2c<\mathit{nums}[i]$, then that node \emph{and its whole left subtree}
(all keys $\le c$) satisfy the test -- add $1+\mathit{size}(\text{left})$ and step right; otherwise step
left. This is \propref{c05-kth}'s rank descent with the comparison changed from $c<x$ to $2c<x$; using a
\emph{treap} keyed by $c$ keeps the descent $\bigO{\log n}$ instead of $\bigO{n}$.}
\fullhint{Insert values right-to-left into a keyed treap whose nodes store $\mathit{size}$. For the current
$x=\mathit{nums}[i]$ call \code{countSmaller(root, x)}: walk down, and when $2\cdot$\code{node.key}$<x$ add
$1+\mathit{size}(\text{node.left})$ and go right, else go left; accumulate into \code{ans}. Then
\code{splitKey} at $x$ and merge the fresh node back in. Every comparison must use \code{long long} because
$\mathit{nums}[j]$ can be $-2^{31}$. Total $\bigO{n\log n}$ expected, immune to input order. Full build below.}
\end{attemptgate}

\solstep{Understanding the Problem}
The naive loop is $\bigO{n^2}$. The reduction is to an \emph{online rank}: read the array right-to-left so
the running structure holds precisely the elements after index $i$, then $\mathit{count}=|\{j>i:
\mathit{nums}[i]>2\mathit{nums}[j]\}|$ is ``how many stored keys $c$ have $2c<\mathit{nums}[i]$.'' Chapter~7
answered the $c<x$ version with a size-augmented BST and paid $\bigO{n^2}$ when the input was sorted, because
that descent's length is the tree's height and the height was not controlled. Here we change only one thing:
make the height logarithmic \emph{by construction} -- a keyed treap (\thmref{c06-exptheight}) -- and the
identical rank descent becomes $\bigO{\log n}$, hence the whole pass $\bigO{n\log n}$. That single swap from
``hope the input is random'' to ``guarantee the height'' is the entire chapter distilled.

\solstep{Key Observations}
\begin{observation}
Right-to-left processing makes the stored multiset exactly $\{\mathit{nums}[j]:j>i\}$, so a rank query on it
counts reverse pairs with first endpoint $i$. Summing over $i$ is the answer.
\end{observation}
\begin{observation}
The rank test is $2c<\mathit{nums}[i]$, a monotone predicate in the key $c$: if it holds at a node it holds
for that node's whole left subtree (all keys $\le c$), so a size-augmented descent accumulates
$1+\mathit{size}(\text{left})$ and steps right -- the mirror of \propref{c05-kth} with a doubled comparison
(\defref{c06-heightcost}).
\end{observation}
\begin{observation}
Using a \emph{treap} keyed by value keeps the descent expected $\bigO{\log n}$ no matter the arrival order
(\thmref{c06-exptheight}); the fragile plain-BST version of this exact algorithm is the $\bigO{n^2}$
Chapter~7 Summit (\probref{c05-summit}). $2\mathit{nums}[j]$ overflows 32-bit at the endpoints -- every
comparison is in \code{long long}.
\end{observation}

\solstep{Think It Through With Me}
\begin{thinkbox}
\deadend{Sort a copy and binary-search each element for its doubled rank.}
{A static sort gives rank among \emph{all} elements, but you need rank among those \emph{to the right}.
The query is inherently online -- you must insert incrementally, and reading rank must be fast \emph{and}
correct after every insert.}
\deadend{Keep seen values in a sorted \code{vector} and \code{lower\_bound} $x/2$.}
{Reading is $\bigO{\log n}$ but \emph{inserting} into a vector is $\bigO{n}$ -- one quadratic traded for
another, exactly the trap in Chapter~7's Summit. You need insert \emph{and} rank both logarithmic, which a
balanced tree gives.}
\aha{A keyed, size-augmented \emph{treap}: one descent per element both routes the insert (split/merge) and
reads ``how many stored keys satisfy $2c<x$'' by summing left sizes while stepping right. The height is
$\bigO{\log n}$ by construction, so the whole pass is $\bigO{n\log n}$ -- the guaranteed upgrade of the
fragile BST of \probref{c05-summit}.}
\end{thinkbox}

\solstep{The Approach Ladder}
\begin{center}\renewcommand{\arraystretch}{1.3}
\begin{tabular}{@{}c p{0.235\linewidth} p{0.30\linewidth} l c@{}}
\toprule
\textbf{\#} & \textbf{Approach} & \textbf{Idea} & \textbf{Time} & \textbf{Verdict} \\
\midrule
1 & \lv{2} double loop & test every pair $(i,j)$ & $O(n^2)$ & \tle \\
2 & \lv{3} size-augmented \emph{plain} BST & Chapter 7's rank walk, doubled test & $O(n\,h)$ & \tle\ on sorted \\
3 & \lv{5} keyed order-statistic \emph{treap} & same rank, height guaranteed & $O(n\log n)$ exp & \ac \\
4 & \lv{5} merge-sort split / Fenwick & count cross pairs, or index the rank & $O(n\log n)$ & \ac \\
\bottomrule
\end{tabular}\end{center}
Rung 2 is the honest Chapter-7 answer and its verdict is ``correct, elegant, worst-case quadratic'' -- the
problem this chapter was written to solve. Rung 3 is \emph{this} chapter's solution: the identical algorithm,
but the treap's \thmref{c06-exptheight} makes the descent logarithmic. Rung 4 leaves the pointer world behind
-- merge counting (\chref{c01}) and a Fenwick over coordinate-compressed values (\chref{c14}) both hit
$\bigO{n\log n}$ worst case because an array index cannot rot.

\solstep{Solution}
Process \code{nums} from the last element to the first, keeping a keyed treap whose nodes store the number
of nodes in their subtree ($\mathit{size}$). For the current value $x$, call \code{countSmaller(root, x)}:
descend the treap as a BST; at each node with key $c$, if $2c<x$ then $c$ and its whole left subtree qualify,
so add $1+\mathit{size}(\text{left})$ and step right, otherwise step left. Add the returned count to
\code{ans}. Then insert $x$: \code{splitKey} the treap into keys $<x$ and keys $\ge x$, and
\code{merge} in a fresh node. Duplicates live as separate nodes, so sizes count occurrences. Every threshold
comparison uses \code{long long}. (Rung 4 gives the guaranteed-$\bigO{n\log n}$ merge and Fenwick variants.)

\solstep{Proof of Correctness}
\begin{proofbox}[The treap rank equals the count of right-side elements under the doubled threshold]
\emph{Invariant:} before handling index $i$, the treap holds exactly the values $\mathit{nums}[j]$ for
$j>i$, and each node's $\mathit{size}$ equals the number of nodes in its subtree (maintained by the
\code{upd} on every split and merge, \lemref{c06-mergevalid}). Consider \code{countSmaller(root, x)} with
$x=\mathit{nums}[i]$. By the BST invariant, a node's left subtree holds keys $\le$ the node key. When the
descent is at a node with key $c$ and $2c<x$, every key $c'$ in that node's left subtree satisfies
$c'\le c$, hence $2c'\le 2c<x$ -- so the node and its entire left subtree (exactly
$1+\mathit{size}(\text{left})$ stored keys) qualify, and they are all disjoint from the right subtree, which
the descent then explores for further qualifiers. When $2c\ge x$, the node and its right subtree (keys $\ge
c$) all fail, so the only qualifiers lie in the left subtree, which the descent enters. The two branches
therefore count, without duplication or omission, precisely the stored keys $c$ with $2c<x$. By the
invariant those are exactly the indices $j>i$ with $2\mathit{nums}[j]<\mathit{nums}[i]$, i.e.\ the reverse
pairs whose first endpoint is $i$. Summing over $i$ from $n-1$ down to $0$ counts every reverse pair exactly
once. \end{proofbox}

\solstep{Implementation}
\begin{lstlisting}[caption={Reverse Pairs -- keyed, size-augmented treap with a 2x-threshold rank (LeetCode 493; O(n log n) expected)}, label={lst:c06-summit}]
#include <bits/stdc++.h>
using namespace std;
mt19937 rng(chrono::steady_clock::now().time_since_epoch().count());
struct Node { int key, pri, sz; Node *l, *r;
    Node(int k) : key(k), pri(rng()), sz(1), l(nullptr), r(nullptr) {} };
int  SZ(Node* t){ return t ? t->sz : 0; }
void upd(Node* t){ if (t) t->sz = 1 + SZ(t->l) + SZ(t->r); }

Node* merge(Node* a, Node* b){                 // all keys in a <= all keys in b
    if (!a) return b; if (!b) return a;
    if (a->pri > b->pri){ a->r = merge(a->r, b); upd(a); return a; }
    else                { b->l = merge(a, b->l); upd(b); return b; }
}
void splitKey(Node* t, int k, Node*& a, Node*& b){   // a: keys < k, b: keys >= k
    if (!t){ a = b = nullptr; return; }
    if (t->key < k){ a = t; splitKey(t->r, k, t->r, b); upd(a); }
    else           { b = t; splitKey(t->l, k, a, t->l); upd(b); }
}
// count stored keys c with 2*c < x  (== x > 2*c), one guaranteed-log descent
int countSmaller(Node* t, long long x){
    int res = 0;
    while (t){
        if (2LL * t->key < x){ res += 1 + SZ(t->l); t = t->r; }   // node + left subtree qualify
        else                 { t = t->l; }                          // only the left subtree can qualify
    }
    return res;
}
class Solution {
public:
    int reversePairs(vector<int>& nums){
        Node* root = nullptr; long long ans = 0;
        for (int i = (int)nums.size() - 1; i >= 0; --i){
            ans += countSmaller(root, nums[i]);           // seen j>i with nums[i] > 2*nums[j]
            Node *a, *b;
            splitKey(root, nums[i], a, b);
            root = merge(merge(a, new Node(nums[i])), b); // insert nums[i]
        }
        return (int)ans;
    }
};
\end{lstlisting}
\begin{notebox}[The rung-4 guaranteed answers, for the record]
\textbf{Merge sort:} recurse on both halves, then before merging count cross pairs -- both halves are
sorted, so for each right value advance a left pointer while $\mathit{nums}[i]\le 2\mathit{nums}[j]$ and add
the leftover left count. Counting uses the doubled test, merging the raw values: $\bigO{n\log n}$ worst case
(\chref{c01}). \textbf{Fenwick:} coordinate-compress all $\mathit{nums}$ and their doubles; scan
right-to-left, \code{ans += query(rank(nums[i])-1)} over the keys already present, then add one at
$\mathit{rank}(2\mathit{nums}[i])$. Each update/query is $\bigO{\log n}$, immune to input order
(\chref{c14}). Both sidestep the tree entirely -- the same robustness the treap buys with rotations-by-chance.
\end{notebox}

\solstep{Trace}
\begin{tracebox}[Official example one -- nums equal to 1, 3, 2, 3, 1]
Process right-to-left; the running tree holds the seen keys, \code{countSmaller} returns those with $2c<x$.
\begin{itemize}[nosep,leftmargin=1.3em]
  \item $i=4$, $x=1$: tree empty $\Rightarrow$ \code{countSmaller}$=0$; insert $1$. \checkmark
  \item $i=3$, $x=3$: at key $1$, $2\cdot1=2<3$ so add $1+\mathit{SZ}(nil)=1$; go right (null). count $=1$;
        insert $3$. \checkmark (only the trailing $1$ pairs with this $3$).
  \item $i=2$, $x=2$: keys $\{1,3\}$. $2\cdot1=2<2$? no $\Rightarrow$ go left off node $1$, no more; count
        $=0$; insert $2$. \checkmark
  \item $i=1$, $x=3$: keys $\{1,2,3\}$. $2\cdot1=2<3$ add $1$; up to key $2$: $2\cdot2=4<3$? no $\Rightarrow$
        its left offers nothing further; count $=1$; insert $3$. \checkmark
  \item $i=0$, $x=1$: need $2c<1$ $\Rightarrow$ none; count $=0$. \checkmark
\end{itemize}
Total $0+1+0+1+0=2$ -- matches the reference \checkmark, the two pairs being $(1,4)$ and $(3,4)$.
\end{tracebox}
\begin{tracebox}[Official example two -- nums equal to 2, 4, 3, 5, 1]
\begin{itemize}[nosep,leftmargin=1.3em]
  \item $i=4$, $x=1$: empty, count $0$; insert $1$.
  \item $i=3$, $x=5$: key $1$ has $2<5$, add $1$; count $1$; insert $5$.
  \item $i=2$, $x=3$: key $1$ has $2<3$, add $1$; key $5$ has $10<3$? no; count $1$; insert $3$.
  \item $i=1$, $x=4$: key $1$ has $2<4$, add $1$; deeper keys $3,5$ give $6,10<4$? no; count $1$; insert $4$.
  \item $i=0$, $x=2$: need $2c<2$ $\Rightarrow$ none; count $0$.
\end{itemize}
Total $0+1+1+1+0=3$ -- matches the reference \checkmark: pairs $(1,4),(2,4),(3,4)$.
\end{tracebox}

\solstep{Complexity}
\begin{complexitybox}
Time $\bigO{n\log n}$ \emph{expected}: each element is one rank descent plus one split and two merges, all
$\bigO{h}$, and a treap's $h=\bigO{\log n}$ in expectation over the random priorities
(\thmref{c06-exptheight}) -- not over the input, which is the point. The plain-BST twin is
$\bigO{n\,h}$ worst-case $\bigO{n^2}$ (\probref{c05-summit}); merge and Fenwick are
$\bigO{n\log n}$ \emph{worst} case. Space $\bigO{n}$. At $n=5\cdot10^{4}$ the quadratic form is
$\approx2.5\cdot10^{9}$ -- \tle -- so the controlled height is exactly what separates \ac{} from \tle.
\end{complexitybox}

\solstep{Verification and Edge Cases}
\begin{center}\renewcommand{\arraystretch}{1.2}
\begin{tabular}{@{}p{0.30\linewidth} p{0.62\linewidth}@{}}
\toprule
\textbf{Case} & \textbf{What the code does} \\
\midrule
Sorted ascending input & treap still short by priorities -- $\bigO{n\log n}$; the case that kills the plain-BST twin \\
All-equal elements & each key $c$ has $2c<x$ false against an equal $x$, so count stays $0$ (no pair) \\
Negative / minimum values & $2\cdot(-2^{31})$ needs \code{long long}; the descent compares $2LL\cdot c$, never overflows \\
Single element & tree empty at the only step, answer $0$ \\
Duplicates & each occurrence is a separate node; $\mathit{size}$ counts them, so ranks stay exact \\
\bottomrule
\end{tabular}\end{center}

\solstep{Pitfalls and Variants}
\begin{warningbox}
Ways to break it: (i) comparing $c$ against $x/2$ instead of $2c$ against $x$ -- integer division truncates
toward zero and miscounts around negatives; keep the multiply and use \code{long long}. (ii) counting only
$1+\mathit{size}(\text{left})$ after going right but forgetting to \emph{continue} into the right subtree --
the qualifiers are not all on one side; the descent must keep walking. (iii) processing left-to-right -- the
seen set must be the elements \emph{after} the current index. (iv) trusting a plain BST: it is correct but
$\bigO{n^2}$ on sorted input -- the whole reason this is a \emph{balanced} tree. \emph{Variants / the real
alternatives:} the merge-sort cross-count (\chref{c01}) and the Fenwick over compressed doubled values
(\chref{c14}) both guarantee $\bigO{n\log n}$; a size-augmented red-black tree is the third way -- someone
must keep $h$ logarithmic, whether by rotation, chance, colour, or an array index.
\end{warningbox}

\begin{insightbox}[Key takeaways -- three things to remember]
\begin{enumerate}[nosep, leftmargin=1.6em]
  \item \emph{The augmentation was never the win -- the height is.} LC 315 and LC 493 use the \emph{same}
        rank descent; only the tree's height differs. Chapter~7 could not promise it; a balanced tree does
        (\defref{c06-heightcost}, \thmref{c06-exptheight}).
  \item \emph{Change one comparison, keep the whole machinery.} $c<x$ became $2c<x$ and every line -- split,
        merge, size, descent -- survived intact. A balanced, augmented tree is a \emph{platform}, not a
        trick (\propref{c05-kth}, \propref{c06-implicit}).
  \item \emph{A pointer tree is one honest rival away from an array index.} Merge and Fenwick hit the same
        $\bigO{n\log n}$ worst case with no tree at all -- know when the guaranteed height is worth the
        pointers and when it is not (\chref{c13}, \chref{c14}).
\end{enumerate}
\end{insightbox}

\begin{hardeningbox}
\begin{itemize}[nosep, leftmargin=1.4em]
  \item \textbf{Worst-case $\bigO{n\log n}$}: the merge-sort cross-count (notebox) is the follow-up -- same
        right-to-left essence, no tree to rot. Full divide-and-conquer is \chref{c01}.
  \item \textbf{Deterministic balance}: swap the treap for an AVL with a \code{size} field -- the rank
        descent is unchanged and the height is guaranteed, not expected (\thmref{c06-avlbound}).
  \item \textbf{Generalises}: this ``online rank under a monotone predicate'' recurs for any order statistic
        -- $k$th smallest, inversion count, and the dynamic-array aggregates of the implicit treap
        (\propref{c06-implicit}); the Euler-tour flatten of \chref{c24} is the same idea on trees.
\end{itemize}
\end{hardeningbox}

\chaptersection{Tool Kit Inventory}

The machine assembled so far -- cumulative through Chapter~8. Chapters~1--7's rows are carried
forward; Chapter~8's rows are appended. Pointers only, never prose.

\begin{center}
\small
\renewcommand{\arraystretch}{1.25}
\begin{tabular}{@{}p{0.24\linewidth} l c p{0.27\linewidth} p{0.16\linewidth}@{}}
\toprule
\textbf{Tool} & \textbf{Since} & \textbf{Level} & \textbf{Trigger phrase} & \textbf{Where} \\
\midrule
\multicolumn{5}{l}{\itshape Chapters 1--7 -- carried forward} \\
Three-slot DFS (pre/in/post) & 3.1 & \lv1 & a node sequence in a named order & \coderef{c03-dfs} \\
BST search / insert / delete descent & 7.1 & \lv2 & ``a BST is given -- does/where does $x$ go?'' & \coderef{c05-search} \\
Order statistics by streaming inorder & 7.3 & \lv2 & kth / min-gap / successor on the sorted line & \coderef{c05-mindiff} \\
Ordered-map interval set (\code{set/map}) & 7.4 & \lv4 & ``add/query/remove intervals online'' & \coderef{c05-range} \\
Size-augmented BST rank (fragile) & 7 (Summit) & \lv5 & ``count smaller to the right -- per-insert rank'' & \coderef{c05-summit} \\
\midrule
\multicolumn{5}{l}{\itshape Chapter 8 -- new this chapter} \\
The rotation (LL/RR/LR/RL) & 8.1 & \lv2 & ``fix a violation, keep the inorder'' & \coderef{c06-rotcode} \\
AVL insert/delete with balance factors & 8.2 & \lv3 & ``hand-rolled BST, shortest, read-heavy'' & \coderef{c06-avlins} \\
Treap split/merge engine & 8.3 & \lv4 & ``balanced tree you can code cold; cut/paste ranges'' & \coderef{c06-split} \\
Implicit-sequence treap & 8.3 & \lv4 & ``insert/reverse/range-sum a sequence by position'' & \coderef{c06-imp} \\
Red-black guarantee for \code{set/map} & 8.4 & \lv3 & ``why is \code{std::set} $\bigO{\log n}$?'' & \thmref{c06-rbheight} \\
Splay: self-adjust by locality & 8.5 & \lv4 & ``hot keys dominate; amortised, no stored state'' & \thmref{c06-amortized} \\
Order-statistic treap rank (Summit) & Summit & \lv5 & ``count under threshold, online, worst-case fast'' & \coderef{c06-summit} \\
\bottomrule
\end{tabular}
\end{center}

\chaptersection{Related Topics}

\begin{itemize}[nosep, leftmargin=1.4em]
  \item \chref{c05} -- backward: this chapter is the answer to that one's cliffhanger. The BST descent, the
        three-case delete, inorder$\equiv$sorted and the \code{set} API are assumed and \emph{never}
        re-derived; every ``$\bigO{h}$'' there becomes a guaranteed ``$\bigO{\log n}$'' here, and its Summit
        (\probref{c05-summit}) is reborn as ours (\probref{c06-summit}).
  \item \chref{c03} -- backward: the rotation is proved by appealing to the inorder walk, taken as a black
        box from that chapter (\thmref{c05-inorder}); we only lean on the fact that a rotation leaves it
        unchanged (\thmref{c06-rotation}).
  \item \chref{c26} -- forward: Cartesian trees and the monotonic stack -- the array$\to$tree bridge this
        chapter \emph{deliberately leaves out}, and the direct successor on the next-chapter pointer.
  \item \chref{c23}, \chref{c24} -- forward: the full splay analysis and link-cut (dynamic-tree) machinery
        live in \chref{c23}; \secref{c06-splay} is only its prerequisite intro. Euler-tour trees
        (\chref{c24}) answer the Water-Tree-style path/subtree queries that a splay also handles.
  \item \chref{c13}, \chref{c14} -- forward: the segment tree and Fenwick tree answer the \emph{same} rank
        and range questions over dense integer keys with array locality and no pointers; the Summit's notebox
        and the mixed set hand their robust forms to those chapters.
\end{itemize}

\begin{lastword}[Last Word -- Chapter 8]
\emph{``Chapter 7 borrowed a rule and spent it everywhere; this chapter paid the interest -- because the rule's
worth depended entirely on a number Chapter 7 could not control.''}

\medskip
A balanced tree is a BST that \emph{knows its own height} (\defref{c06-heightcost}). From that one
guarantee everything in this chapter is a corollary: the rotation is the unique inorder-preserving repair
(\thmref{c06-rotation}); AVL enforces $|\bal|\le1$ and fixes any edit in $\le2$ turns, buying
$h\le1.44\,\lg(n+2)$ (\thmref{c06-avlbound}, \thmref{c06-fixheight}); the treap replaces every case with a coin
flip and wins $\bigO{\log n}$ in expectation, plus a \code{split}/\code{merge} engine that turns it into a
dynamic array (\thmref{c06-exptheight}, \propref{c06-implicit}); red-black relaxes the rule to a colouring and
still proves $h\le2\lg(n+1)$ -- the very promise \propref{c05-logn} deferred, now cashed for
\code{std::set} (\thmref{c06-rbheight}); and splay stores nothing at all yet earns $\bigO{\log n}$ amortised by
rotating habit toward the root (\thmref{c06-amortized}).

\medskip
Carry three reflexes forward:
\begin{enumerate}[nosep, leftmargin=1.6em]
  \item \emph{Never trust a plain BST's height.} If keys can arrive sorted, a ``$\bigO{\log n}$'' claim is a
        lie unless a rotation, a priority, a colour, or an array index \emph{keeps} it logarithmic
        (\defref{c06-heightcost}, \figref{c06-chain}).
  \item \emph{One move, many policies.} Every structure here is the same local surgery
        (\thmref{c06-rotation}) applied under a different rule -- so learn the rotation cold and the four cases
        (\defref{c06-rotcases}), and the rest is bookkeeping you can \emph{choose}.
  \item \emph{Controlled height makes augmentation safe.} A size field on a fragile tree is a nicer-looking
        $\bigO{n^2}$; on a balanced tree it is $\bigO{\log n}$ rank, select, and range -- the Summit
        (\probref{c06-summit}) is precisely that difference.
\end{enumerate}

\medskip
And three habits, drilled until boring: on any ``ordered set under adversarial input'' problem, first ask
``is the height mine to control?'' -- if not, pick a balancer before you code; write the AVL \code{balance}
routine and the treap \code{split}/\code{merge} from memory once each, because they are the two engines every
later dynamic-tree chapter reuses; and reach for \code{std::set}/\code{std::map} freely \emph{knowing} it is
red-black -- then remember the one question it refuses to answer fast, rank, which is exactly what an
augmented tree or a Fenwick is for.

\medskip
\textbf{Next:} \chref{c26} turns from trees you \emph{keep} balanced to trees you \emph{read out of an array}.
A Cartesian tree is built in $\bigO{n}$ by a monotonic stack, trades rotations for ``the smallest element is the
root,'' and unlocks largest-rectangle, subarray minimums, and next-smaller links -- the array$\to$tree bridge
this chapter deliberately left unbuilt so you would feel the difference between balancing a search tree and
deriving one from the shape of the data.
\end{lastword}


\closechapter
